% ============================================================
% PAPER 1 of 11 -- obstruction_calculus
% Assembled 2026-09-29 by content-and-merits partition.
% Title: Obstruction Certificates under Incomplete Observation
% Source: diffs/obstr_v57.tex
% Partition note: Single-source. The theory core: obstruction certificates, the necessity direction.
%% STATUS: structurally assembled as a NEW version. Editorial integration
%%         (de-duplication of shared framework, sibling cross-citations)
%%         is NOT yet done. See PAPERS_MANIFEST.md.
%% ============================================================

%% v60 (2026-09-29): added subsection 3.8, mechanized verification of the discrete
%%   core. EVERY claim in it was checked before being written: Lean 4.14.0 installed,
%%   lake build run (60/60 jobs, exit 0), #print axioms inspected for the named
%%   declarations (sorryAx absent), and a comment-stripped source scan run over all
%%   of Formalizations/ (no sorry/admit/axiom). Scope limits stated: continuous-time
%%   results NOT formalized; completeness is wrt policy trees, not stationary
%%   policies. Do not strengthen these claims without re-running the checks.
%% v59 (2026-09-29): added subsection 3.7, a worked case in which a certificate
%%   returns an exact verdict on a continuum (the two-stock fibre benchmark),
%%   verified in exact rational arithmetic. Closes the plan's gap: at least one
%%   worked case where a certificate bites on a system whose kernel cannot be
%%   computed by enumeration. Scope stated in the subsection: fibre feasibility,
%%   not the dynamic kernel.
% SIBLING PAPERS: Paper 2 (Exact Probabilistic Sufficiency), Paper 3 (Computational Certification of Obstruction), Paper 4 (Minimax Dual Certificates), Paper 5 (Exact Belief Computation).
% Cross-cite, do not re-derive: the selector principle, the epistemic kernel
% and the observation structure are defined in Paper 1. Papers 2-5 cite it.

% SIBLING PAPERS: Paper 2 (Exact Probabilistic Sufficiency), Paper 3 (Computational Certification of Obstruction), Paper 4 (Minimax Dual Certificates), Paper 5 (Exact Belief Computation).
% Cross-cite, do not re-derive: the selector principle, the epistemic kernel
% and the observation structure are defined in Paper 1. Papers 2-5 cite it.

\documentclass[10pt,twocolumn]{article}
\usepackage[a4paper,margin=15mm]{geometry}
\usepackage{amsmath,amssymb}
\usepackage{amsthm}
\theoremstyle{definition}
\newtheorem{theorem}{Theorem}
\newtheorem{proposition}{Proposition}
\newtheorem{corollary}{Corollary}
\newtheorem{remark}{Remark}
\newtheorem{example}{Example}
\newtheorem{definition}{Definition}
\newtheorem{openproblem}{Open Problem}
\usepackage{graphicx}
\usepackage{booktabs,array,calc}
\usepackage[font=small,labelfont=bf]{caption}
\usepackage[colorlinks=true]{hyperref}
\usepackage{etoolbox}
\providecommand{\tightlist}{\setlength{\itemsep}{0pt}\setlength{\parskip}{0pt}}
\setcounter{secnumdepth}{-1}
\emergencystretch=3em
\allowdisplaybreaks
\AtBeginDocument{%
  \BeforeBeginEnvironment{equation}{\small}%
  \BeforeBeginEnvironment{equation*}{\small}%
  \BeforeBeginEnvironment{align}{\small}%
  \BeforeBeginEnvironment{align*}{\small}%
  \BeforeBeginEnvironment{gather}{\small}%
  \BeforeBeginEnvironment{gather*}{\small}%
  \BeforeBeginEnvironment{multline}{\small}%
  \BeforeBeginEnvironment{multline*}{\small}%
  \BeforeBeginEnvironment{flalign}{\small}%
  \BeforeBeginEnvironment{flalign*}{\small}%
}


\begin{document}
\title{Certifying nonviability under incomplete observation: the obstruction calculus and its probabilistic counterpart, in exact arithmetic}
\begin{abstract}
\noindent\textbf{The general claim.} Under incomplete observation, the question ``is any
observation-based policy viable?'' has two answers that the literature treats as different
subjects: a combinatorial one, phrased on information states, and a probabilistic one, phrased
on distributions over them. This paper develops both and shows they are the \emph{same object at
two resolutions}, not competing approximations. The probabilistic theory does not soften the
combinatorial one into a degree of belief; it \emph{prices} it. The value-one level of the
belief-state safety recursion coincides exactly with the viable-set recursion of the calculus,
and the $\alpha$-vectors that represent the value are exactly the indicators of the maximal
jointly survivable subsets of the posterior support --- an antichain, and therefore
Sperner-bounded. So the probabilistic reading inherits the certificate structure rather than
replacing it: where the calculus says \emph{nonviable}, the belief value is bounded away from one
by a deficit that is itself exact and attained on natural instances. \emph{Certifying that
something cannot be done is a finite, checkable, exact act --- and it survives the passage from
sets to measures without any loss of exactness.}

\noindent\textbf{Why two parts.} Each part supports half of that claim. Part I develops the
instrument on \emph{information states}: sound sufficient conditions for nonviability ---
equivalently, necessary conditions for viability --- finitely checkable in the common-action,
fibre-certification and finite-horizon forms, with analytic drift-and-timing conditions
elsewhere. Part II asks what those certificates imply when the information state is a
\emph{distribution}, and carries the whole theory in exact rational arithmetic. Neither is the
other's corollary. The calculus is complete in finite systems by backward recursion and says
nothing about probability; the belief value assigns a number to every belief and says nothing
about the shape of the certificate. Their exact agreement at the value-one level is a theorem
relating them, not a reduction of one to the other, and the paper keeps the two theories
distinct so that the agreement stays a result rather than a definition.

\noindent\textbf{What is found.} Part I establishes six mechanisms --- two finitely checkable,
two closed-form conditional, one minimal, and a sixth under a policy-class restriction --- as a
nonexhaustive taxonomy of information-theoretic failure, each paired with the response
correspondence it empties and the \emph{design consequence} it licenses. They do not exhaust the
complement of the epistemic kernel; in finite systems backward recursion is complete. Part II
gives six matching results: the support identity above; exactness and parametric closed forms
over rational input data; the survival of the certificates as deficit lower bounds, attained on
natural instances under both symmetric and asymmetric priors; cell-by-cell agreement of the
certified delayed-regime class with its closed form, with the class declaration responsible for
exactly the 36 timing combinations on the grid; and an adaptive-observation audit that prices
late resolution exactly. Across both parts, \emph{the observation layer, not the forecast layer,
sets the value} --- the design rules are the ones the deterministic audits impose, now with a
price attached.

\noindent\textbf{Structure.} Part I is the obstruction calculus and Part II its probabilistic
counterpart. Each is self-contained --- its own results, tables, discussion and conclusions ---
and each is reproduced in full from its companion paper without condensation, opening with that
paper's own abstract. A cross-part conclusion follows, stating what transfers between the
set-valued and the probabilistic readings and what does not.
\end{abstract}
\maketitle

\tableofcontents

\newpage

\part*{Part I --- The obstruction calculus on information states}
\label{part:calculus}
\addcontentsline{toc}{part}{Part I --- The obstruction calculus}

\maketitle

\begin{abstract}

Under perfect measurement the viability kernel --- the set of states
from which some feedback keeps the system within its constraints --- is
characterized by tangency conditions. Under incomplete observation the
sufficiency direction has a canonical answer in Veliov's
output-feedback regulation condition and the estimation-tube reduction.
The complementary direction --- certifying that no observation-based
policy is viable --- has received less systematic treatment.

This paper develops that instrument: a calculus of obstruction
certificates, sound sufficient conditions for nonviability (equivalently, necessary conditions for viability), finitely
checkable in the common-action, fibre-certification, and finite-horizon forms and analytic drift-and-timing conditions elsewhere. They do not exhaust the complement of the epistemic kernel, the kernel's observation-based counterpart; in finite systems backward recursion is complete. Five mechanisms are established --- two finitely checkable, two
closed-form conditional, one minimal --- with a sixth under a
policy-class restriction. The central common-action obstruction shows
that when the safe controls of compatible states intersect emptily, no
observation-based policy is viable, though every compatible state is
individually viable under full information. The others are a finite-time
exit certificate under an Isaacs-type drift condition, an
epistemic-emptiness construction in which a constant observation merges
states with incompatible admissible controls, a delayed-information
obstruction with an explicit timing bound, and a certification limit:
an exact observation-only certifier exists iff safe-set
membership is constant on the observation fibres. The sixth, a
certainty-equivalence trap, empties the kernel of the uncorrected controller under a biased
observation even when the perfect-information kernel is nonempty.
Two certification-theoretic completions are established: over polytope-declared
blind-window classes the timing certificate is exactly computable as a linear
program, with the Farkas dual as the obstruction witness, sound beyond the
declaration with an unbounded implementable-class gap; and the residual dynamic
gap is shown to decompose exactly, its post-observation recourse phase
undecidable by any certificate pair that reads only the window sub-model.

The calculus is positioned against barrier certificates and estimation
tubes, and its monitoring and institutional consequences ---
observation timing, coarseness, aggregation, and bias --- are drawn.

\end{abstract}

\textbf{Keywords:} viability kernel; obstruction certificates; robust
viability; epistemic viability; output feedback; sustainability
governance; partial observation

\textbf{Mathematics Subject Classification:} 49J53; 93B03; 93C41; 91B76

\subsection{1. Introduction}\label{calc-introduction}

\subsubsection{1.1 The problem}\label{calc-the-problem}

Viability theory characterizes the states of a constrained control system from which at least one control can keep every future state within a constraint set (Aubin, 1991). For sustainability problems the constraints are floors --- stock levels, service thresholds, safety margins --- and viability is the formal counterpart of maintaining a development path rather than merely optimizing it (B\'en\'e, Doyen, and Gabay, 2001; De Lara and Doyen, 2008;
Doyen et al., 2012; Doyen and Gajardo, 2020). The theory under
\emph{perfect measurement} is complete in the relevant sense: the
viability kernel exists as the largest viable subset, and Nagumo-type
tangency conditions certify it.

Sustainability governance, however, operates under \emph{incomplete
observation}. Stocks are assessed at discrete review intervals.
Indicators are coarser than the state. Some stocks are unobserved
altogether. And the information that does arrive may arrive late. This paper addresses that side: for an incomplete observation structure, it develops instruments that certify that no observation-based policy is viable --- that the infeasibility is due to the information structure, not to the dynamics.

Floors in sustainability assessment are typically constraints on the
productive base --- the state of the system that yields the measured
output --- rather than on the output itself. Such a base can be read as
natural capital, as a stock, or as a slowly regenerating flow of
services; these are overlapping readings of the same asset rather than
mutually exclusive ones, and what makes a use of it sustainable is
whether the use falls on the yield or on the base itself. Across
assets, whether the base behaves as a flow or as a stock is a
continuum, set by its regeneration timescale relative to the rate of
use --- a continuum that runs from a season, through a year, to
geological time for a mineral deposit. If a base regenerates too
slowly for the rate at which it is taken, drawdown is still
liquidation.

A base can also erode while measured output is maintained, because
use draws on the base rather than on its yield: the measured quantity
can be held up while the base declines, and partial observation may
not detect the erosion until a floor is breached. The signal that would reveal the erosion is
accordingly often the one that is coarse, postponed, or absent.
Incomplete observation is therefore structural to the certification
problem, not merely a practical limitation; the sufficiency literature
and the obstruction calculus of this paper address the same underlying
indeterminacy rather than competing ones.

The question splits naturally. The sufficiency direction has a canonical literature: Veliov (1993) gives conditions for the existence of an output-feedback regulation map under incomplete and inexact measurement, reducing under perfect measurement to the classical viability condition (Haddad, 1981); the estimation-tube programme (Cardaliaguet, Quincampoix, and Saint-Pierre, 2007) absorbs imperfect information into an estimation space of sets of compatible states with no loss of value; and barrier certificates certify safety (Prajna and Jadbabaie, 2004; Prajna, Jadbabaie, and Pappas, 2007), with converse and hybrid necessary-and-sufficient characterizations (Prajna and Rantzer, 2005; Maghenem and Sanfelice, 2019). What that literature does not supply is the converse direction: a calculus of obstruction certificates, each a sufficient condition for nonviability whose negation is a necessary condition for observation-based viability. An obstruction certificate is a checkable witness that no observation-based policy exists --- finite in the finite-state and polyhedral cases --- and the argument it makes is that a prescribed class of observation-based policies fails not because a particular policy is bad, but because the information structure leaves no room for any policy. In the polyhedral common-action and certification forms (Theorem~\ref{calc-thm:common-action} and Proposition~\ref{calc-prop:fibre}) it is a finite, checkable test; the drift certificates of Theorem~\ref{calc-thm:exit} and Theorem~\ref{calc-thm:delayed} are not finite objects, one exhibiting an enforcement selection and the other bounding the timing uniformly over every policy. Veliov's condition tells us when output feedback can work; the obstruction calculus tells us when it cannot, and for the borderline cases it identifies which quantitative feature of the observation design --- timing, coarseness, bias, aggregation --- is responsible.

\subsubsection{1.2 Contributions}\label{calc-contributions}

Five obstruction mechanisms are developed, each with a complete proof,
and a sixth is exhibited under a policy-class restriction:

\begin{enumerate}
\def\labelenumi{\arabic{enumi}.}
\item
  \textbf{The instantaneous common-action obstruction} \emph{(finitely checkable)} (Theorem~\ref{calc-thm:common-action},
  Example~\ref{calc-ex:hidden-mode}). If the safe-control sets of the states compatible with the
  current information set intersect emptily, then no observation-based
  policy is viable, even though every compatible state is individually
  viable under full information. The failure is purely informational: no
  stochasticity, no estimation error.
\item
  \textbf{The delayed-information obstruction} \emph{(closed-form conditional)} (Theorem~\ref{calc-thm:delayed}). If the
  information needed to distinguish safe from unsafe compatible states
  arrives only at time \(T_{\mathrm{obs}}\), and the disturbance can
  force constraint violation earlier than \(T_{\mathrm{obs}}\), then no
  observation-based policy is viable. The timing bound
  \(\inf q / \varepsilon < T_{\mathrm{obs}}\) is stated explicitly:
  information may be accurate but arrive too late.
\item
  \textbf{The finite-time exit certificate} \emph{(closed-form conditional)} (Theorem~\ref{calc-thm:exit}). Under a
  Dini-drift condition of Isaacs type --- the disturbance chooses after
  the control --- the disturbance can force exit from the safe set
  within an explicit time bound, for \emph{every} admissible control of
  any information structure. This is the base obstruction: when it
  applies, the viability question is closed without any
  observation-theoretic argument.
\item
  \textbf{Epistemic emptiness by admissibility} \emph{(minimal construction)} (Proposition~\ref{calc-prop:emptiness}, a minimal
  construction). A constant observation merges states whose
  \emph{admissible} controls differ: the state-dependent control set
  forces a single common action, and that action is unsafe at the
  boundary, so the unresolved initial observation fibre lies outside the
  epistemic kernel, even though every physical state is viable under
  full information. The mechanism is admissibility --- the common-action
  obstruction of Theorem~\ref{calc-thm:common-action} at minimal size --- and Example~\ref{calc-ex:hidden-mode} gives the hidden-mode instance.
\item
  \textbf{The fibre certification criterion} \emph{(finitely checkable)} (Proposition~\ref{calc-prop:fibre}, Corollary~\ref{calc-cor:certainly-safe}).
  An exact observation-only safety certifier --- a function of the
  observation that returns ``safe'' exactly on the safe set --- exists
  if and only if safe-set membership is constant on every observation
  fibre. The certainly-safe set is the sound relaxation: the largest set
  of observations that can be labeled safe without further information.
\item
  \textbf{The certainty-equivalence trap} \emph{(policy-class restriction)} (Remark~\ref{calc-rem:ce-trap}). Even an
  \emph{injective} observation empties the kernel of a system whose
  perfect-information kernel is nonempty for the canonical
  certainty-equivalence controller --- the perfect-information zeroing
  law applied to the uncorrected observation. The trap is a
  policy-specific failure: a controller that uses the observation map's
policy-specific failure: a controller that uses the observation map's
structure (the bias correction) restores viability.
\item
  \textbf{Computable instantiation and the exact residual} \emph{(finitely checkable;
  exact residual)} (Theorem~\ref{calc-thm:lp-instant}, Example~\ref{calc-ex:relax-gap},
  Propositions~\ref{calc-prop:decomposition} and \ref{calc-prop:window-nogo}). Over
  polytope-declared blind-window classes the timing certificate is exactly
  computable as a linear program --- polynomial size, exact in rational
  arithmetic, Farkas-certified on failure --- sound beyond the declaration,
  with the implementable-class gap shown unbounded by an explicit instance on
  which the relaxation extinguishes the certificate. The residual gap of the
  completeness question decomposes exactly into the window phase, decided
  completely by the certificates, and the post-observation recourse phase,
  decided by no window-measurable certificate pair.
\end{enumerate}

The calculus is organized by a single selector principle (Proposition~\ref{calc-prop:selector}, Section 2.4):
observation-based viability is the existence of one response ---
admissible, safe, and recursively viable --- for every compatible state,
and each certificate proves that a specific response correspondence is
empty; the obstructions are nested in a ladder
(Proposition~\ref{calc-prop:ladder}). Five refinements complete the picture:
backward belief recursion is sound and complete on any finite horizon
in finite systems (Theorem~\ref{calc-thm:finite-horizon}); a uniform-margin
condition connects the instantaneous and tube obstructions
(Proposition~\ref{calc-prop:uniform-margin}); the timing obstruction has a
sharp threshold form (Remark~\ref{calc-rem:sigma}); a Helly-type sparse
witness bounds the common-action obstruction by \(m+1\) compatible
states (Proposition~\ref{calc-prop:helly}); and a comparison-function form
sharpens the exit-time bound (Remark~\ref{calc-rem:comparison}). The epistemic kernel is monotone in
action sets, disturbances, information, and policy class
(Proposition~\ref{calc-prop:monotone}).

Throughout, ``obstruction'' means a \emph{certified} failure. Theorem~\ref{calc-thm:exit} and Theorem~\ref{calc-thm:delayed} each exhibit the violating constraint, an admissible disturbance,
and the quantitative bound that make failure inevitable, for every policy in the prescribed class; Proposition~\ref{calc-prop:emptiness} and Theorem~\ref{calc-thm:common-action} exhibit the incompatible
admissible controls --- the merged observation whose incompatible safe
controls certify emptiness --- with no quantitative bound claimed there;
and Proposition~\ref{calc-prop:fibre} is the converse-side characterization of the certification
limit itself. The calculus is positioned against
the barrier-certificate literature (Section 6.2) and the estimation-tube
programme (Section 6.3), and its consequences for the design of
monitoring and indicators in sustainability governance are drawn in
Section 6.4.

\subsubsection{1.3 Related work}\label{calc-related-work}

The viability-theory background is Aubin (1991), Aubin, Bayen, and Saint-Pierre (2011), and the kernel approximation theory of Saint-Pierre (1994); robust viability is Aubin and Frankowska (1990) and Frankowska (1989). Under incomplete measurement, the sufficiency theorem is Veliov (1993), viability with a priori unknown but observable parameters is Quincampoix and Veliov (1994), and the estimation-set reduction is Cardaliaguet, Quincampoix, and Saint-Pierre (2007).  Doyen (2000) studies output feedback for uncertain nonlinear systems under state and control constraints and obtains necessary and sufficient conditions for a property close in spirit to ours: a set \(K\) is a \emph{guaranteed viability domain} when there exists a Lipschitz selection
\(u(\cdot)\) of the control constraint map such that \(K\) is invariant under the
differential inclusion \(\dot x \in \{f(x,u(h(x,w)),w): w \in W(x)\}\), the disturbance
acting on both the dynamics and the output; his characterization is by non-emptiness, for
some \(\lambda>0\), of the Lipschitz kernel \(L_{\lambda}(R^{\lambda}_{K})(y)\) at every
output \(y\), and a viable feedback is synthesized by Steiner selection.  Doyen's result is
the nearest neighbour to the question posed here, and it differs from it on three axes, each
visible in his own statement.  First, the \emph{policy class}: the definition quantifies over
Lipschitz selections, and the closed loop is \(u(h(x,w))\), a memoryless function of the
current output --- so that a policy which is non-Lipschitz, or which uses the history of
observations rather than the current output alone, lies outside the range of his necessity.
The obstruction certificates are asserted for observation-based policies generally and are
neither implied by nor subsumed under his condition.  Second, the \emph{property}: his is
exact invariance of a closed set, one particular viability property among those addressed
here.  Third, the \emph{direction}: the characterization is an equivalence whose constructive
content is a feedback-building corollary, while the systematic obstruction side ---
certifying that no observation-based policy is viable --- is not developed there, and Doyen
records the reason: ``to our knowledge, no general viability result is available in this
differential game context with imperfect and/or partial information.''  It is that
undeveloped side, and not the synthesis, that the obstruction calculus is meant to serve. On the failure side, Aubin (2001) characterizes the complement of the kernel --- Poincar\'e's ``shadow'' --- together with the capture basins, and Aubin and Catt\'e (2002) supply the bilateral-fixed-point and discriminating-kernel calculus; the certificates of Section 3 are one-sided, observation-constrained analogues of that programme. In verification, barrier certificates originate with Prajna and Jadbabaie (2004), with worst-case and converse results in Prajna, Jadbabaie, and Pappas (2007) and Prajna and Rantzer (2005) and hybrid characterizations in Maghenem and Sanfelice (2019). The dominant computational tradition in safety verification reaches the same question from the opposite side.  Hamilton--Jacobi reachability formulates safety as a two-player
differential game and recovers the backward reachable tube as the zero sublevel set of the
viscosity solution of a time-dependent Hamilton--Jacobi--Isaacs equation (Mitchell, Bayen,
and Tomlin, 2005; Margellos and Lygeros, 2011; for a survey, Bansal, Chen, Herbert, and
Tomlin, 2017), the analytic foundation being the viscosity-solution theory of Crandall, Ishii,
and Lions (1992).  It does not reach the setting considered here, for a reason that is
structural rather than technical.  The value function is posed on the physical state
\(x \in \mathbb{R}^{n}\), and it is the state space itself that the grid --- or the
collocation points of a neural solver --- discretizes.  Under incomplete observation the
sufficient statistic is not a state but a belief, a probability measure over states, and the
value function would accordingly have to live on an infinite-dimensional space of measures;
level-set and physics-informed solvers are discretization schemes for finite-dimensional
domains and do not carry over.  The obstruction calculus sidesteps the difficulty by not
solving a differential equation at all: the certificates of Section 3 are finite algebraic
objects, checkable without discretizing any continuum.  Two further contrasts are worth
recording.  Even with perfect observation the Hamilton--Jacobi certificate is only as exact
as its discretization: grid-based solvers scale exponentially in the state dimension and are
in practice confined to five or six dimensions, and the learning-based approaches that lift
that bound (Bansal and Tomlin, 2021) purchase scalability with approximation --- in one
benchmark against a six-dimensional grid reference, a neural backward reach--avoid tube
recovers roughly four fifths of the true unsafe set.  And in \emph{direction} the tradition
is, like viability theory proper, a sufficiency machine: it computes the states from which
some control against all disturbance preserves safety, which is precisely the direction
Veliov's condition and the estimation-tube reduction already serve.  The obstruction
certificates supply the complementary direction, and do so in the observation-constrained
regime in which neither tradition returns an answer. The sustainability application domain is B\'en\'e, Doyen, and Gabay (2001), De Lara and Doyen (2008), Doyen et al. (2012), and Doyen and Gajardo (2020). To our knowledge, the obstruction certificates of Theorem~\ref{calc-thm:common-action}, Theorem~\ref{calc-thm:delayed}, and Proposition~\ref{calc-prop:fibre} --- in particular the certification criterion and the timing bound --- have not been stated in this form; the elementary facts on which they rest (quantifier commutation and Dini comparison) are classical. Relative to Doyen (2000) the residual contribution is correspondingly exact, and is the content of the three distinctions above: a certificate that is checkable, that holds against \emph{every} observation-based policy --- not merely against Lipschitz memoryless selections of the current output --- and that applies to viability properties beyond the exact maintenance of a closed domain. It is in that sense that the obstruction calculus does not compete with his characterization but occupies the side of it that he leaves open.

\subsubsection{1.4 Organization}\label{calc-organization}

Section 2 fixes the framework: the control system, the viability kernel
hierarchy, observation structures, information states, the safe-control
correspondence, and the selector principle (Proposition~\ref{calc-prop:selector}). Section 3 develops the obstruction calculus (Theorem~\ref{calc-thm:finite-horizon}--\ref{calc-thm:exit}, Propositions~\ref{calc-prop:uniform-margin}--\ref{calc-prop:emptiness}, Remarks~\ref{calc-rem:sigma}--\ref{calc-rem:comparison}, and Example~\ref{calc-ex:hidden-mode}): finite-horizon completeness, the instantaneous common-action obstruction with its ladder and sparse witness, the delayed-information obstruction with its threshold form, and --- as supporting certificates --- the finite-time exit certificate and the admissibility obstruction, together with a recourse certificate for the post-observation mode (Proposition~\ref{calc-prop:recourse}, Example~\ref{calc-ex:three-branch}). Section 4 develops the certification limits
(Proposition~\ref{calc-prop:fibre}, Corollary~\ref{calc-cor:certainly-safe}, Remark~\ref{calc-rem:ce-trap}). Section 5 reviews the sufficiency landscape, cited rather than re-derived. Section
6 discusses the position of the calculus relative to barrier
certificates and estimation tubes, states the monotonicity of the
epistemic kernel (Proposition~\ref{calc-prop:monotone}), and draws the governance
consequences. Sections 7--10 develop the new material: complete characterizations and the residual gap (Section 7), the selector principle with its dual witnesses and the monitoring-adequacy converse (Section 8), a calibrated case study with a coverage audit (Section 9), and the probabilistic belief-state development (Section 10). Section 11 concludes. Brief sketches and short proofs of the main-text statements are given in place; the extended discussions and worked instances of the Sections 3, 4, and 10 certificates, the complete proofs of the Section 3 certificates, the expanded sufficiency review, and the auxiliary constructions are collected in the Supplementary Material; the results of Sections 7 and 9 are proved in the main text. The running
instances are fixed here for orientation: the two-floor conflict
(Supplementary S1), on which the worked instances of the Sections 3, 4,
and 10 certificates are built, and the calibrated delayed-hidden-regime
case study (Section 9).



\subsection{2. Framework}\label{calc-framework}

\subsubsection{2.1 The control system and the kernel
hierarchy}\label{calc-the-control-system-and-the-kernel-hierarchy}

This section fixes the mathematical objects used throughout the paper:
the control system, the disturbance structure, the constraint set, and
the four kernels (perfect-information, robust, epistemic, and
institutionally restricted) that organize the analysis.

Fix a state space \(X \subseteq \mathbb{R}^n\), a control set \(U\), a
disturbance set \(D\), and dynamics
\[\dot x(t) = f(x(t), u(t), d(t)), \qquad u(t) \in U(x(t)), \quad d(t) \in D(x(t)),\]
with \(f\) continuous, \(U\) and \(D\) set-valued with closed graph and
nonempty compact values, and admissible controls and disturbances taken
to be measurable selections --- the standing solution-existence conditions:
every selection pair admits a Carath\'eodory solution, and the
convexified relaxed inclusion of Section 3.5 satisfies the Marchaud
conditions of the viability literature (Aubin, 1991). Let \(\mathcal{V} \subseteq X\) be a closed
constraint set. A \emph{causal policy} \(\pi\) is a map from the
observation record to controls, respecting the specified information structure. The standard objects are:

\begin{itemize}
\tightlist
\item
  \(\mathrm{Viab}(\mathcal{V}; U, \pi_{\mathrm{perf}})\): the
  \textbf{viability kernel} under full-information (state-feedback)
  policies --- the largest closed subset of \(\mathcal{V}\) from which
  there exists a state-feedback control keeping the trajectory in
  \(\mathcal{V}\) for all time --- no disturbance, or a known
  disturbance trajectory (Aubin, 1991).
\item
  \(\mathrm{RViab}(\mathcal{V})\): the \textbf{robust kernel}, the
  largest subset of \(\mathcal{V}\) from which there exists a
  state-feedback policy keeping every trajectory --- for \emph{every}
  admissible disturbance realization --- in \(\mathcal{V}\) (Aubin and
  Frankowska, 1990).
\item
  \(\mathrm{ERViab}_{\mathcal{I}}(\mathcal{V})\): the \textbf{robust
  epistemic kernel} --- the object of this paper (Definition~\ref{calc-def:kernel}). Its
  elements are \emph{states of information}, not physical states
  (Section 2.3). The non-robust counterpart
  \(\mathrm{EViab}_{\mathcal{I}}(\mathcal{V})\) is the non-robust contrast class (Section 2.3).
\end{itemize}

Projected to physical state space --- with
\(\mathrm{proj}_{\exists}(\mathfrak{B}) = \bigcup_{B \in \mathfrak{B}} B\)
the \emph{existential} projection of a collection \(\mathfrak{B}\) of
information states, applied to the first two terms, whose elements are information states --- the informational hierarchy
reads
\[
\begin{aligned}
\mathrm{proj}_{\exists}\big(\mathrm{IRViab}_{\mathcal{J}}(\mathcal{V})\big)&\;\subseteq\; \mathrm{proj}_{\exists}\big(\mathrm{ERViab}_{\mathcal{I}}(\mathcal{V})\big) \\
&\;\subseteq\; \mathrm{RViab}(\mathcal{V}) \;\subseteq\; \mathrm{Viab}(\mathcal{V}),
\end{aligned}
\]
where \(\mathrm{IRViab}_{\mathcal{J}}(\mathcal{V})\) is the institutionally restricted kernel of Section 6.4. The physical implication of epistemic viability is the
inclusion at the level of beliefs:
\(B_0 \in \mathrm{ERViab}_{\mathcal{I}}(\mathcal{V}) \Rightarrow B_0 \subseteq \mathrm{RViab}(\mathcal{V})\).
Each inclusion can be strict, with a distinct cause: robust contraction
arises from disturbances; epistemic contraction from
indistinguishability; institutional contraction from restricted
authority, enforcement, and allocation. The present paper characterizes
the \emph{epistemic} contraction: the mechanisms by which
indistinguishability empties or shrinks the kernel.

\subsubsection{2.2 The safe-control
correspondence}\label{calc-the-safe-control-correspondence}

For a constraint set \(\mathcal{V}\) described by finitely many \(C^1\)
constraint functions,
\(\mathcal{V} = \{ x : q_j(x) \ge 0,\ j = 1..m \}\), define the
\textbf{safe-control correspondence} at \(x \in \mathcal{V}\):
\[
\begin{aligned}
\mathcal{R}_{\mathcal{V}}(x) \;=\; \bigl\{ u \in U(x) : {}& \forall d \in D(x),\; \forall j \text{ with } q_j(x) = 0,\\
& \nabla q_j(x) \cdot f(x, u, d) \ge 0 \bigr\},
\end{aligned}
\]
the set of controls whose drift is inward (subtangential) on every
active constraint face, uniformly over the disturbance. For \(x\) in the
interior of \(\mathcal{V}\) (no active constraints),
\(\mathcal{R}_{\mathcal{V}}(x) = U(x)\). Two levels of the tangency condition are distinguished throughout. (i) \emph{Local reading.} The correspondence is applied to the constraint set itself: by Nagumo's theorem (Nagumo, 1942) in its robust form, if \(\mathcal{V}\) is closed and
\(\mathcal{R}_{\mathcal{V}}(x) \neq \varnothing\) on
\(\partial\mathcal{V}\) with a measurable selection
\(u(x) \in \mathcal{R}_{\mathcal{V}}(x)\), then \(\mathcal{V}\) is
itself robustly invariant,
i.e.~\(\mathrm{RViab}(\mathcal{V}) = \mathcal{V}\) (Aubin, 1991;
Frankowska, 1989). (ii) \emph{Kernel reading.} When \(\mathcal{V}\) is not invariant, viability of a point \(x_0 \in \mathcal{V}\) requires instead a selection \(u(x) \in \mathcal{R}_G(x)\) along the trajectory, where \(G = \mathrm{RViab}(\mathcal{V})\) and \(\mathcal{R}_G\) is the same correspondence computed on the kernel: points of
\(\mathcal{V} \setminus G\) may admit controls that are instantaneously
safe while every trajectory must exit later. The obstruction
theorems below are \emph{local certificates against
\(\mathcal{R}_{\mathcal{V}}\)}: exhibiting a state at which even the
weaker pointwise condition fails certifies exclusion from the kernel,
and this is the reading in which the certificates are used. The
correspondence is the bridge between the pointwise (tangency) and the
global (kernel) description, and it is the object the observation
structure acts upon: an observation-based policy can select controls
only as a function of the observation record.

\subsubsection{2.3 Observation structures, information states, and
epistemic
kernels}\label{calc-observation-structures-information-states-and-epistemic-kernels}

\textbf{Conventions.} Throughout, the plant evolves in continuous time
and a policy acts at review times, holding the selected action until the
next review; the observation record is generated by the observation map
\(O\) (the error, hidden-parameter, and delay extensions are noted
below). Between reviews the record \(y(t)=O(x(t))\) is observed
continuously but cannot change the held action; the certificates of
Sections 3.2--3.3 and 3.6 use the record only through what it discriminates
(observation-equivalence), so they remain valid under any refinement of
the within-interval record that does not separate compatible branches.
A policy is observation-based if its action depends on the
record only through the information set \(B_t\); under the delay-free
set-membership semantics used here the information set is a sufficient
statistic, so observation-based and record-based policies coincide (with
delays, hidden parameters, or history-dependent constraints the belief
must be augmented and sufficiency is not automatic; Section 6.5). The disturbance
is chosen adversarially after the control (the Isaacs order of
Theorems~\ref{calc-thm:exit} and \ref{calc-thm:common-action}), an action outside
\(U(x)\) at a compatible state \(x\) is itself a failure (Section 2.4),
and constraint violation is the strict inequality \(q<0\).

An \textbf{observation structure} \(\mathcal{I} = (Y, O)\) is a
measurable space \(Y\) and an observation map \(O : X \to Y\). The
policy observes \(y(t) = O(x(t))\). At time \(t\), the
\textbf{information set} (belief) of a regulator who knows the dynamics,
the disturbance class, its own applied controls \((u(s))_{s<t}\), and the record \((y(s))_{s \le t}\) is
\[
\begin{aligned}
B_t \;=\; \bigl\{ \xi(t) : {}& \xi(0) \in B_0,\\
& \dot\xi(s) = f(\xi(s), u(s), d(s)),\ d(s) \in D(\xi(s)),\\
& O(\xi(s)) = y(s) \text{ for all observed } s \le t \bigr\},
\end{aligned}
\]
where the applied control history \(u(\cdot)\) in the dynamics is
the regulator's own --- common to every candidate trajectory --- and
\(\xi(\cdot)\) ranges over candidate trajectories: the set of current
states compatible with the record and the applied controls. We assume the standard set-membership semantics (Kurzhanski and V\'alyi, 1997), so that \(B_t\) contains every state
consistent with the observations and the applied controls. (With
observation error, the equality becomes membership
\(y(s) \in H(\xi(s), v(s))\); with hidden parameters the belief lives in
\(X \times \Theta\); with delays it lives in a history space --- the extensions are noted as open in Section 6.5.) A policy is
\textbf{observation-based} if \(u(t)\) depends on the record only
through \(B_t\); this is a special case of record-based policies, and
the two coincide when the belief is a sufficient statistic for the
decision problem, which holds in the delay-free set-membership semantics used
throughout. We write \(U^B(B) = \bigcap_{x \in B} U(x)\) for the
controls available at every compatible state --- the only controls an
observation-based policy may select without risking inadmissibility.

\begin{definition}[robust epistemic kernel]\label{calc-def:kernel}

\(B_0 \in \mathrm{ERViab}_{\mathcal{I}}(\mathcal{V})\) if there exists a
nonanticipative observation-based policy \(\pi\) such that, for every
compatible initial state \(x_0 \in B_0\) and every admissible
disturbance realization \(d(\cdot)\), the trajectory of
\(\dot x = f(x, u(t), d(t))\) generated from \(x_0\) by \(\pi\)'s
actions under \(d(\cdot)\) --- the record being the one that
\((\pi, d(\cdot))\) generates --- remains in \(\mathcal{V}\) for all
time. (This is the standard exists-strategy-for-all-realizations form. The quantifier order is essential: the record is generated by the realization, so conditioning the admissible realizations on the record would be circular.) Admissible initial
beliefs are the compact information sets \(B_0 \subseteq X\) generated
by the observation structure --- the observation fibres \(O^{-1}(y)\),
\(y \in O(X)\), in the constructions of Sections 3.2--3.3 and 3.6. The existential (non-robust) counterpart \(\mathrm{EViab}_{\mathcal{I}}(\mathcal{V})\) is the non-robust contrast class.

\end{definition}

We use \(\mathrm{ERViab}\) throughout, since the disturbance classes of sustainability problems are adverse by construction. The theorems are stated for this robust notion; they do not transfer a fortiori to the non-robust contrast class \(\mathrm{EViab}_{\mathcal{I}}(\mathcal{V})\), in which there exist an observation-based policy and \emph{some} admissible disturbance realization under which every compatible trajectory remains in \(\mathcal{V}\) for all time. Since the robust kernel is contained in the non-robust one, an observation-based policy may exist there that no robust policy survives.

The thesis is now precise:
\textbf{the epistemic kernel is the greatest recursively viable collection of information states in the sense of the estimation-space reduction cited in Section 5; common-action compatibility is a necessary one-step condition, but not by itself a
complete characterization --- and the obstruction calculus certifies,
mechanism by mechanism, the information states outside it.}

\subsubsection{2.4 Notation}\label{calc-notation}

We collect the principal symbols used throughout the paper and fix the two structural symbols: the \textbf{observation structure} is always
\(\mathcal{I}\) (calligraphic I) and the \textbf{institution} of Section
6.4 is always \(\mathcal{J}\) (calligraphic J); no other letter serves either role. The state space is \(X \subseteq \mathbb{R}^n\); the
control set is \(U\) (set-valued \(U(x)\) at state \(x\)); the
disturbance set is \(D\) (set-valued \(D(x)\)). The dynamics are
\(\dot x = f(x,u,d)\). The closed constraint set is
\(\mathcal{V} \subseteq X\), described by finitely many \(C^1\)
constraint functions \(q_j\) (\(j = 1,\ldots,m\)). The safe-control
correspondence at \(x \in \mathcal{V}\) is
\(\mathcal{R}_{\mathcal{V}}(x)\). The observation structure is
\(\mathcal{I} = (Y, O)\) with observation map \(O : X \to Y\); the
information set (belief) at time \(t\) is \(B_t\); the initial belief is
\(B_0\) --- admissible initial beliefs are the compact information sets
the observation structure generates, the observation fibres
\(O^{-1}(y)\), \(y \in O(X)\), in the constructions. The common
admissible action set at belief \(B\) is
\(U^B(B) = \bigcap_{x \in B} U(x)\) --- the only controls an
observation-based policy may select without risking inadmissibility,
under the convention, used from Proposition~\ref{calc-prop:emptiness} on, that an action outside
\(U(x)\) at a compatible state is itself failure; the common safe-action
set is
\(\mathcal{R}_{\mathcal{V}}^B(B) = \bigcap_{x \in B} \mathcal{R}_{\mathcal{V}}(x)\).
The tube-safe action set at review length \(\Delta\) is
\(\mathcal{A}_{\mathrm{tube}}(B, \Delta) = \{ a : \mathrm{Reach}_{[0,\Delta]}(B, a) \subseteq \mathcal{V} \}\),
where \(\mathrm{Reach}_{[0,\Delta]}(B, a)\) is the \emph{robust} ---
all-disturbance, all-solution --- reachable set of the system from \(B\) under the
held action \(a\) over \([0,\Delta]\). The first informative observation
time is \(T_{\mathrm{obs}}\): the first time the observation record
discriminates between compatible branches --- before it, the branches of
Theorem~\ref{calc-thm:delayed}'s (H3.1) produce identical records. The adverse-selection
correspondences of the exit certificates are
\(D_{\varepsilon}(x,u) = \{ d \in D(x) : D^+ q(x; f(x,u,d)) \le -\varepsilon \}\)
(Theorem~\ref{calc-thm:exit}) and
\(D_{\eta}(x) = \{ d \in D(x) : \nabla q_j(x) \cdot f(x, a, d) \le -\eta/2 \}\)
(Theorem~\ref{calc-thm:common-action}'s proof, local). The upper
right Dini derivative of \(q\) in direction \(v\) is
\(D^+ q(x; v) = \limsup_{h \downarrow 0} [q(x + hv) - q(x)]/h\); for
\(q\) of class \(C^1\) --- the class assumed in the exit theorems ---
\(D^+ q(x; v) = \nabla q(x) \cdot v\).

The kernels of the hierarchy are
\(\mathrm{Viab}(\mathcal{V}; U, \pi_{\mathrm{perf}})\) (perfect
information), \(\mathrm{RViab}(\mathcal{V})\) (robust), and
\(\mathrm{ERViab}_{\mathcal{I}}(\mathcal{V})\) (the robust epistemic
kernel --- Definition~\ref{calc-def:kernel}, the object of this paper), with
\(\mathrm{proj}_{\exists}(\mathfrak{B}) = \bigcup_{B \in \mathfrak{B}} B\)
the existential projection of a collection \(\mathfrak{B}\) of information states. Consistency check: when the observation \(O\) is injective the record determines the state, so the epistemic and robust kernels agree, \(\mathrm{ERViab}_{\mathcal{I}}(\mathcal{V}) = \mathrm{RViab}(\mathcal{V})\), the informational certificates of Sections 3.2--3.3 and 3.6 reduce to the no-merge case, and only the dynamic obstruction of Section 3.5 can certify exclusion --- the definitions degenerate to classical robust viability exactly when observation is perfect. The non-robust counterpart
\(\mathrm{EViab}_{\mathcal{I}}(\mathcal{V})\) is the non-robust contrast class (Section 2.3). The institutionally restricted kernel is \(\mathrm{IRViab}_{\mathcal{J}}(\mathcal{V})\) (Section 6.4). The policy classes
are \(\pi_{\mathrm{perf}}\) (state feedback), \(\Pi_{\mathrm{output}}\),
\(\Pi_{\mathrm{state}}\), and \(\Pi_{\mathrm{CE}}\) --- the
certainty-equivalence class defined in Remark~\ref{calc-rem:ce-trap}: fixed
perfect-information regulation laws applied uncorrected to the
observation. The policy-specific safety set of Section 4.1 is
\(\mathrm{Safe}^{\Pi}_T\); the certainly-safe set of Corollary~\ref{calc-cor:certainly-safe} is
\(\mathcal{Y}_{\mathrm{safe}}\).

A few letters carry local meanings, as declared here: the strip width \(a\) of
Theorem~\ref{calc-thm:exit}'s \(\mathcal{S}_a\) is not the held action
\(a = \pi(B)\) of Theorem~\ref{calc-thm:common-action}'s proof, nor the
pathway vector \(a \ge 0\) of Section 5(d); \(\lambda\) is the Farkas
multiplier of Theorem~\ref{calc-thm:common-action}'s checkability
certificate and the observer decay rate of Section 5(c); and
\(\alpha\) is the comparison function of Remark~\ref{calc-rem:comparison}
and a multiplier of Section 5(d)'s substitution alternative. The
letter \(K\) denotes the safe set of Section 4's certification
problems, while in Section 5(c) it denotes the kernel
itself, with the buffered set \(K_\zeta\) and eroded kernels
\(K^{-c\zeta}\) of Section 5(c). The observation-error map of
Section 2.3 is \(H(\xi, v)\), whereas the harvest vector of
the Supplementary Material is the lowercase \(h\); and the finite-horizon recursion
set \(\mathcal{W}_N\) of Section 3.1 is calligraphic, distinct from
the Lyapunov function \(W = S_1 + S_2\) of the Supplementary Material.

\medskip
\noindent\textbf{Shared symbols across the companion papers.} The
companion papers (worked systems, computational certification,
probabilistic sufficiency, exact belief computation, minimax dual
certificates, applied regime viability, and the applied forecast ladder)
deliberately share one formal vocabulary; the letters below carry
paper-local meanings here that differ there:
\begin{center}\small
\begin{tabular}{@{}p{0.38\linewidth}p{0.52\linewidth}@{}}
\toprule
\textbf{Here} & \textbf{In the companions} \\
\midrule
\(K\): the safe set of Section 4's certification problems; the kernel of Section 5(c) & window depth and contact set \(\mathcal{K}(B)\) (minimax duals); the full-information kernel \(K^{T}_{\mathrm{full}}\) (certification); stock carrying capacities (forecast ladder; minimax benchmark) \\
\(\lambda\): the Farkas multiplier of Theorem 3's certificate; the observer decay rate of Section 5(c) & certificate weight vectors (certification; minimax duals) \\
\(\Gamma_{h}\): the recourse certificate of Section 5(d) & the certificate margin \(\Gamma_{\mathcal{I}}\) (certification); the alpha-vector witness sets \(\Gamma^{\Pi}_{k}\) (sufficiency) \\
\(\tau\): delay and observation times & the worst-case exit time \(\tau^{*}\) (worked systems) \\
\(\sigma\): the recourse construction's trajectory parameter & the kernel time variable (certification; minimax duals); the floored exit count \(\sigma^{*}\) (worked systems); a noise scale (forecast ladder) \\
\(\mu\): a Farkas multiplier & an adversarial probability measure (certification; minimax duals) \\
\bottomrule
\end{tabular}
\end{center}


The four correspondences of this section are one skeleton. An
observation-based regulator selects a single response for an entire
information set --- a held action, a feedback segment, or a
safe/unsafe verdict --- and observation-based viability is the existence
of a response that is simultaneously admissible, safe, and recursively
viable for every compatible state and disturbance history. Each
obstruction certificate of Sections 3 and 4 is a checkable sufficient
condition under which one of these response correspondences is empty:
\(U^B(B)\) (admissibility), \(\mathcal{R}_{\mathcal{V}}^B(B)\)
(instantaneous safety), \(\mathcal{A}_{\mathrm{tube}}(B,\Delta)\)
(review-interval safety), or the recursive predecessor of Section 3.1.
The mechanisms differ in which correspondence fails, and the
obstructions are nested (Proposition~\ref{calc-prop:ladder}). The organizing
claim of the calculus is the following, stated once and used implicitly
by every certificate of Sections 3 and 4.

\begin{proposition}[selector principle]\label{calc-prop:selector}
If \(B\in\mathrm{ERViab}_{\mathcal{I}}(\mathcal{V})\) and \(\pi\) is a
policy witnessing viability on \(B\), then at every review time the
action \(a=\pi(B)\) selected by \(\pi\) is a response that is
(i) \emph{admissible}, \(a\in U^B(B)\); (ii) \emph{safe}, keeping every
compatible trajectory in \(\mathcal{V}\) until the next observation ---
\(\mathrm{Reach}_{[0,\Delta]}(B,a)\subseteq\mathcal{V}\) at review
length \(\Delta\); and (iii) \emph{recursively viable}, in that every
post-observation belief reached at the next review is again in
\(\mathrm{ERViab}_{\mathcal{I}}(\mathcal{V})\) under the continuation
of \(\pi\). A belief is therefore viable only if each of the response
correspondences of this section --- \(U^B(B)\),
\(\mathcal{R}_{\mathcal{V}}^B(B)\), \(\mathcal{A}_{\mathrm{tube}}(B,\Delta)\),
and the recursive predecessor of Section 3.1 --- is nonempty along
every reachable branch; hence a certificate that any one of them is
empty at a reachable belief certifies nonviability. In finite systems
the converse is exact (Theorem~\ref{calc-thm:finite-horizon}); in continuous
time the converse is the backward-recursion completion noted in
Section 6.5.
\end{proposition}



\subsection{3. The Obstruction Calculus}\label{calc-the-obstruction-calculus}

\subsubsection{3.1 Finite-horizon completeness (finite systems)}\label{calc-finite-horizon-completeness}

In the finite setting the calculus is complete: backward recursion over beliefs --- the set-membership analogue of the belief-state recursion of partially observable Markov processes (Smallwood and Sondik, 1973) --- is sound and complete on any finite horizon, and nonviability is certified by a finite obstruction tree. The certificates of Sections 3.2--3.3 and 3.5--3.6 are sufficient conditions, not a complete characterization of the epistemic kernel.

Fix finite state, action, disturbance, and observation spaces \(X\),
\(A\), \(D\), \(Y\), a transition map \(x^+ = F(x,a,d)\), an observation
map \(y = O(x,a,d,x^+)\), a safe set \(\mathcal V\subseteq X\), and
admissible-action sets \(U(x)\subseteq A\). For a belief \(B\subseteq X\),
an action \(a\), and an observation \(y\), the post-belief is
\[
\begin{gathered}
\mathrm{Post}(B,a,y) = \bigl\{ x^+\in X : \exists\, x\in B,\ d\in D(x), \\
x^+ = F(x,a,d),\ y = O(x,a,d,x^+) \bigr\},
\end{gathered}
\]
and \(y\) is possible under \((B,a)\) if \(\mathrm{Post}(B,a,y)\neq\varnothing\).
With \(U^B(B)=\bigcap_{x\in B}U(x)\), define the one-step predecessor of a
collection \(\mathcal C\) of beliefs by
\[
\begin{gathered}
\mathrm{Pre}(\mathcal C) = \bigl\{ B\subseteq\mathcal V : \exists\, a\in U^B(B)\ \text{such that}\\
\mathrm{Post}(B,a,y)\in\mathcal C\ \text{for every possible } y \bigr\},
\end{gathered}
\]
and the finite-horizon kernels by \(\mathcal W_0 = \{B : B\subseteq\mathcal V\}\),
\(\mathcal W_{k+1} = \mathrm{Pre}(\mathcal W_k)\).
The \(N\)-step common safe-action set at a belief \(B\) is
\(\mathcal{A}_N(B) = \bigl\{ a \in U^B(B) : \mathrm{Post}(B,a,y) \in \mathcal W_{N-1} \text{ for every possible } y \bigr\}\),
so that \(B \in \mathcal W_N \iff \mathcal{A}_N(B) \neq \varnothing\). The recursion is exact but of worst-case complexity exponential in
\(|X|\) --- the belief space contains up to \(2^{|X|}\) beliefs --- so
finite-checkability is per belief, not a polynomial-time global procedure.

\begin{theorem}[finite-horizon soundness and completeness]\label{calc-thm:finite-horizon}
An initial belief \(B_0\) admits an observation-based policy that keeps
every compatible trajectory in \(\mathcal V\) for \(N\) steps if and only
if \(B_0\in\mathcal W_N\).
\end{theorem}

\emph{Proof sketch.} Both directions are inductions on the horizon. Necessity: a policy safe for \(N\) steps picks \(a \in U^B(B_0)\), and every possible successor belief is safe for \(N - 1\) steps, so \(B_0 \in \mathrm{Pre}(\mathcal{W}_{N-1})\). Sufficiency: choose the witness action at each step, and every reachable state remains in \(\mathcal{V}\) for \(N\) steps. If \(B_0 \notin \mathcal{W}_N\), the failure is certified by a finite obstruction tree alternating regulator and nature turns --- the complete construction, with the one-step abstraction of Example~\ref{calc-ex:hidden-mode} as its smallest tree, is given in the Supplementary Material (S1; the tree is drawn in Figure~S2).

\subsubsection{3.2 The instantaneous common-action
obstruction}\label{calc-the-instantaneous-common-action-obstruction}

The instantaneous obstruction is the static core of the calculus: it certifies nonviability at a single information state, without a timing or horizon argument.

\begin{theorem}[common-action obstruction]\label{calc-thm:common-action}

Let \(B \subseteq \mathcal{V}\) be an information set of the specified observation structure (the safe-control correspondence
\(\mathcal{R}_{\mathcal{V}}\) is defined on \(\mathcal{V}\), so the
belief is taken inside it). Suppose the following hypotheses hold.

(H2.1) The selected action is held until the next observation (sampled
review: no informative observation arrives before the next action must
be selected, and the policy holds the selected action over the review
interval --- the institutional reading; without a positive dwell time
the instantaneous statement below must be replaced by the tube-safety statement below).

(H2.2) Either the \textbf{common admissible action set} is empty,
\[U^B(B) \;=\; \bigcap_{x \in B} U(x) \;=\; \varnothing, \qquad \text{(admissibility obstruction)}\]
or the \textbf{common safe-action set} is empty while common
admissibility holds,
\begin{gather}
\mathcal{R}_{\mathcal{V}}^B(B) \;:=\; \bigcap_{x \in B} \mathcal{R}_{\mathcal{V}}(x) \;=\; \varnothing, \tag{1}\\
\text{with } U^B(B) \neq \varnothing \quad \text{(safety obstruction)} \notag
\end{gather}

(H2.3) \textbf{Closed-loop existence of the local adverse selection
(safety case).} The disturbance correspondence is lower semicontinuous,
or locally constant near the boundary point at which the obstruction
fires, so that the local adverse selection of the proof is realized
along the compatible trajectory; alternatively, the certificate is read
against the convexified (relaxed) inclusion, which preserves the
signed drift (for \(C^1\) \(q\) the directional derivative is affine in the velocity).

Then, under hypotheses (H2.1)--(H2.3),
\(B \notin \mathrm{ERViab}_{\mathcal{I}}(\mathcal{V})\): every
compatible state may be individually robustly viable, and the belief is
nevertheless nonviable, because the state-specific safe actions are
incompatible.

\end{theorem}

\emph{Proof sketch.} If the selected action is inadmissible at some compatible state, the policy fails there outright. Otherwise \(a \in U^B(B)\), and the emptiness of the common safe-action set is caused by the boundary: some \(\bar x \in B \cap \partial\mathcal{V}\) has \(a \notin \mathcal{R}_{\mathcal{V}}(\bar x)\), giving an active constraint \(j\) and a disturbance \(\bar d\) with \(\nabla q_j(\bar x) \cdot f(\bar x, a, \bar d) < 0\). A measurable local adverse selection (closed graph of \(D\), continuity of the drift; realized by (H2.3)) then drives \(q_j\) strictly downward from \(q_j(\bar x) = 0\) within the first review period, before any observation can arrive.

For polyhedral data the emptiness certificate is a Farkas pair --- a finite, independently checkable certificate of epistemic infeasibility --- and for infinite beliefs the Helly-type witness of Proposition~\ref{calc-prop:helly} supplies a finite-witness reduction under the convexity hypotheses of Section 3.4; the complete discussion is given in the Supplementary Material (S1).

\begin{remark}[numerical robustness of the polyhedral certificate]\label{calc-rem:robust-farkas}
Linearity makes the Farkas certificate robust in two practical senses. First,
bounded data perturbation degrades the margin gracefully: if
\(\mu \ge 0\) certifies \(\mu^{\top} A = 0\) and \(\mu^{\top} b \le -\mu_0 < 0\),
then for perturbed right-hand sides \(b + \Delta b\) with
\(\|\Delta b\|_\infty \le \epsilon\) the same multipliers give
\(\mu^{\top}(b + \Delta b) \le -\mu_0 + \|\mu\|_1 \epsilon\), still negative
while \(\epsilon < \mu_0/\|\mu\|_1\); the margin \(\mu_0\) is exactly the
certificate's numerical robustness, and normalized multipliers
(\(\|\mu\|_1 = 1\), as in Section~3.4's worked certificate) maximize the
tolerable perturbation. Second, verified approximate data suffice: if the
multipliers are computed with residual \(r := \hat\mu^{\top} \hat A\) and
objective \(c := \hat\mu^{\top} \hat b\), then feasibility of some
\(u \in U\) would imply \(-h_U(-r) \le r^{\top} u \le c\), where
\(h_U(g) := \sup_{v \in U} g^{\top} v\) is the support function of the action
set; hence the test \(c + h_U(-r) < 0\) --- with all integration, rounding,
and enclosure errors charged into \(c\) and \(r\) before it is applied ---
certifies infeasibility without exact stationarity. Both extensions share the
orientation rule for nonviability certificates: uncertain row data must be
relaxed toward a larger feasible controller set, never toward a smaller one;
tightening \(\hat b\) on the strength of an unverified worst case can
manufacture a false obstruction.
\end{remark}

\begin{figure}[htbp]
\centering
\includegraphics[width=0.82\columnwidth]{figs_p2/fig_p2_common_action.png}
\caption{Incompatible safe controls. The states \(x_1\) and \(x_2\)
share an observation, but their safe-action sets \([0,0.4]\) and
\([0.6,1]\) are disjoint, so no common safe action exists
(Theorem~\ref{calc-thm:common-action}).}
\label{calc-fig:common-action}
\end{figure}

\textbf{The tube-safety form.} The instantaneous condition (1) is the
\(\Delta \downarrow 0\) boundary approximation of the exact one-step
obstruction: with the action held over a review interval of length
\(\Delta\), the exact condition is the emptiness of the tube-safe action
set
\[\mathcal{A}_{\mathrm{tube}}(B, \Delta) \;=\; \Big\{ a : \mathrm{Reach}_{[0,\Delta]}(B, a) \subseteq \mathcal{V} \Big\},\]
with \(\mathrm{Reach}_{[0,\Delta]}(B, a)\) the robust (all-disturbance) reachable set (the reading specified in Section 2.4),
and \(\mathcal{A}_{\mathrm{tube}}(B, \Delta) = \varnothing\) certifies
\(B \notin \mathrm{ERViab}_{\mathcal{I}}(\mathcal{V})\) for every review
interval of length \(\Delta\) and every target collection. The theorems
above and below use the instantaneous and timing readings; the tube form
is the unifying one-step object.

\begin{proposition}[uniform margin gives a tube obstruction at every review length]\label{calc-prop:uniform-margin}
Let \(B\subseteq\mathcal V\) be a compact information set with
\(U^B(B)\neq\varnothing\), \(f\) and \(\nabla q_j\) continuous, and
suppose there is \(\eta>0\) such that for every \(a\in U^B(B)\) there
exist a boundary state \(x_a\in B\cap\partial\mathcal V\), an active
constraint \(j\) with \(q_j(x_a)=0\), and a disturbance \(d_a\in D(x_a)\)
satisfying
\[\nabla q_j(x_a)\cdot f(x_a,a,d_a) \le -\eta.\]
Under the closed-loop existence convention of
Theorem~\ref{calc-thm:common-action}, the tube-safe action set is then empty
at every review length,
\[\mathcal{A}_{\mathrm{tube}}(B,\Delta)=\varnothing \qquad \forall\,\Delta>0,\]
and hence \(B\notin\mathrm{ERViab}_{\mathcal I}(\mathcal V)\) under
sampled review of any length.
\end{proposition}

\emph{Proof sketch.} For fixed \(a \in U^B(B)\), continuity of the drift and the closed graph of \(D\) extend the margin to \(-\eta/2\) on a neighbourhood of \((x_a, a, d_a)\), and a measurable local selection realizes it; the trajectory from \(x_a\) under the held action then satisfies \(q_j(x(t)) \le -(\eta/2) t < 0\) for small \(t\), so \(a \notin \mathcal{A}_{\mathrm{tube}}(B, \Delta)\) for every \(\Delta > 0\). As \(a\) was arbitrary, the tube-safe action set is empty.

\begin{proposition}[obstruction ladder]\label{calc-prop:ladder}
Under the standing regularity assumptions and the closed-loop existence
convention of Theorem~\ref{calc-thm:common-action}, the common response sets
of Section 2.4 are nested:
\[\mathcal{A}_{\mathrm{tube}}(B,\Delta) \;\subseteq\; \mathcal{R}_{\mathcal{V}}^B(B) \;\subseteq\; U^B(B) \qquad \forall\,\Delta>0,\]
and emptiness descends the ladder:
\[
\begin{aligned}
U^B(B)=\varnothing \;&\Longrightarrow\; \mathcal{R}_{\mathcal{V}}^B(B)=\varnothing \;\Longrightarrow\; \mathcal{A}_{\mathrm{tube}}(B,\Delta)=\varnothing \\
&\Longrightarrow\; B\notin\mathrm{ERViab}_{\mathcal{I}}(\mathcal{V}).
\end{aligned}
\]
\end{proposition}

\emph{Proof sketch.} The inclusion \(\mathcal{R}_{\mathcal{V}}(x) \subseteq U(x)\) gives the second nesting. For the first, an action outside \(\mathcal{R}_{\mathcal{V}}^B(B)\) violates \(\nabla q_j(x) \cdot f(x, a, d) \ge 0\) at some active boundary state, and the corresponding trajectory leaves \(\mathcal{V}\) in arbitrarily small time, so the action is not tube-safe.

The obstruction ladder's four rungs --- admissibility, common safe-action, tube, and recursive --- are itemized and shown strict in general in the Supplementary Material (S1, Figure~S1); the dynamic exit certificate of Theorem~\ref{calc-thm:exit} sits below the ladder, emptying a response set at some compatible state even under full information.

\begin{example}[hidden-mode conflict]\label{calc-ex:hidden-mode} Let an unobserved parameter satisfy
\(\theta \in \{-1, +1\}\), with \(\dot z = \theta u\),
\(u \in \{-1, +1\}\), and constraint \(z \ge 0\). At \(z = 0\): if
\(\theta = +1\), only \(u = +1\) is safe; if \(\theta = -1\), only
\(u = -1\) is safe. Both compatible states are individually robustly
viable under full information (each admits its safe control), but the
information set \(B = \{(0, +1), (0, -1)\}\) satisfies
\(\mathcal{R}_{\mathcal{V}}^B(B) = \{+1\} \cap \{-1\} = \varnothing\),
so \(B \notin \mathrm{ERViab}_{\mathcal{I}}(\mathcal{V})\). This is a
purely informational failure: no stochasticity and no estimation-quality
argument is involved. The example is the two-action skeleton of every
``the stock is either recovering or collapsing, and the safe policy
differs between the two'' situation.

\end{example}

\subsubsection{3.3 The delayed-information
obstruction}\label{calc-the-delayed-information-obstruction}

\begin{theorem}[delayed-information obstruction]\label{calc-thm:delayed}

Let \(B_0\) be an initial information set. Suppose the following
hypotheses hold.

(H3.1) No informative observation arrives before time
\(T_{\mathrm{obs}} > 0\) --- formally, there are branches that are
\textbf{observation-equivalent on \([0, T_{\mathrm{obs}})\)}: they
produce identical observation records throughout that interval, so the
policy cannot distinguish them before \(T_{\mathrm{obs}}\); between
observations the information set evolves by the set-membership semantics
of Section 2.3.

(H3.2) There exist a constraint function \(q\) of class \(C^1\) (the class of Theorem~\ref{calc-thm:exit}'s proof, Section 3.5) with \(\mathcal{V} \subseteq \{ q \ge 0 \}\)
and a constant \(\varepsilon > 0\) such that, \textbf{for every implementable blind-window control} --- every
measurable \(u(\cdot)\) on \([0, T_{\mathrm{obs}})\) that is admissible at
every compatible state throughout the blind window, \(u(t) \in U^B(B_t)\)
with \(B_t\) the set-membership belief propagated under \(u(\cdot)\) ---
which is exactly the class of controls that an observation-based policy
can realize before the first informative observation, where its actions
are a fixed function of time --- there are a compatible initial state \(x^* \in B_0\) attaining
\(q(x^*) = \inf_{x \in B_0} q(x)\) (with \(B_0\) compact and \(q\) lower
semicontinuous; otherwise replace the infimum by \(\inf q + \delta\) for
an arbitrary \(\delta\)-minimizer, let \(\delta \downarrow 0\), and read
the timing bound with \(\inf q + \delta\)) and an admissible disturbance
realization such that, along the realized trajectory from \(x^*\) under
that open-loop control, the record remains compatible and
\[D^+ q(x(t); f(x(t), u(t), d(t))) \;\le\; -\varepsilon \qquad \text{while } q(x(t)) > 0, \tag{2}\]

(H3.3) The timing bound holds:
\[T_{\mathrm{obs}} \;>\; \frac{\inf_{x \in B_0} q(x)}{\varepsilon}. \tag{3}\]

Then, under hypotheses (H3.1)--(H3.3),
\(B_0 \notin \mathrm{ERViab}_{\mathcal{I}}(\mathcal{V})\): information
may be accurate but arrive too late.

\end{theorem}

\emph{Proof sketch.} By (H3.1) the policy's actions on \([0, T_{\mathrm{obs}})\) form one fixed open-loop control; if it is ever inadmissible at a compatible state the policy fails outright, otherwise (H3.2), applied to this implementable control, supplies a compatible state \(x^*\) with \(q(x^*) = \inf_{B_0} q\) and an admissible disturbance realizing the drift (2); integrating (2) gives \(q(x(t)) \le \inf_{B_0} q - \varepsilon t\), so the violation \(q < 0\) occurs by time \(\inf_{B_0} q / \varepsilon\), which (3) places strictly before \(T_{\mathrm{obs}}\) --- before any informative observation, so the policy cannot react in time.

Hypothesis (2) is a condition on the data of the problem, not on the policies, and condition (3) is the timing bound: inverted, it is the design requirement that the first informative observation precede the enforced exit time. The complete discussion of both hypotheses, and the formal-zero-margin comparison with Theorem~\ref{calc-thm:common-action}, are given in the Supplementary Material (S1).

\begin{remark}[threshold form of the timing obstruction]\label{calc-rem:sigma}
The delayed-information obstruction has a sharp threshold. Let
\(\sigma^*(B_0)\) be the guaranteed blind-window survival time: the
largest duration such that some pre-observation policy keeps every
observation-equivalent branch in \(\mathcal{V}\) until then,
\[\sigma^*(B_0) \;=\; \sup_{u(\cdot)} \inf_{x_0\in B_0,\, d(\cdot)} \tau(x_0,u,d),\]
with the infimum over branches that remain observation-equivalent on
\([0,T_{\mathrm{obs}})\) and \(\tau\) the first violation time under the
convention of Theorem~\ref{calc-thm:delayed}. Then
\[\sigma^*(B_0) < T_{\mathrm{obs}} \;\Longrightarrow\; B_0 \notin \mathrm{ERViab}_{\mathcal{I}}(\mathcal{V}),\]
because before \(T_{\mathrm{obs}}\) a policy cannot condition its action
on which branch realizes. Condition (3) certifies this threshold: under
(H3.1)--(H3.2) every branch exits by time
\(\inf_{x\in B_0}q(x)/\varepsilon\), hence
\(\sigma^*(B_0)\le\inf_{x\in B_0}q(x)/\varepsilon\), and (3) states that
this upper bound lies below \(T_{\mathrm{obs}}\); the threshold itself
may be strictly smaller, in which case nonviability holds even when (3)
fails. The threshold is the blind-window response correspondence of the
timing layer: its emptiness is the timing certificate, the
delayed-information analogue of the tube and common-action
correspondences.
\end{remark}

\textbf{Checkability of the timing certificate.} The drift hypothesis (3) is finitely checkable whenever the blind-window control class is finite: if \(B_{0}\) is finite and the admissible blind-window controls form a finite family (hold-until-\(T_{\mathrm{obs}}\) controls with \(u\) in a finite set, or piecewise-constant controls on a finite partition), the certificate reduces to the finite check that every candidate control enforces exit before \(T_{\mathrm{obs}}\), each worst-case exit time computed by branchwise integration or by the recursion of Section 3.1. The hidden-regime instance of Section 9 is this check in closed form: the worst branch declines at unit rate, giving \(\sigma^{*}(B_{0}) = z_{0} - 1\), and the certificate fires exactly when \(T_{\mathrm{obs}} > z_{0} - 1\). For continuous control classes the check remains a template, part of
Open Problem~\ref{calc-op:dynamic}; over polytope-declared classes, however, the
check is exact and finite (Theorem~\ref{calc-thm:lp-instant}), and the class
declaration itself is load-bearing (Example~\ref{calc-ex:relax-gap}).

\begin{theorem}[linear-programming instantiation of the timing certificate]\label{calc-thm:lp-instant}
Let the blind window have \(K\) discrete steps, finitely many branches \(j \in J\) with
robust affine dynamics \(x^{+} = A_{j} x + B_{j} u + E_{j} w\), \(w \in W_{j}\) a polytope,
known initial branch states \(x_{j,0}\), polyhedral safe set \(\mathcal{V} = \{x : Fx \le g\}\),
and declared blind-window class \(\Pi_{B} = \{(u_{0},\dots,u_{K-1}) : u_{k} \in U\}\) for a
polytope \(U\). Then: (i) \(\sigma^{*}(B_{0}; \Pi_{B})\) is computed exactly by at most
\(\lceil \log_{2} K \rceil + 1\) linear programs of size polynomial in \(K\), \(|J|\), and the
description lengths of \(U\), \(W_{j}\), \(\mathcal{V}\) --- unrolling the dynamics makes each
robust floor row one linear inequality in the control sequence, the disturbance entering only
through precomputed support constants --- and with rational data every step is exact.
(ii) Infeasibility at horizon \(K\) is certified by Farkas multipliers: nonnegative weights on
the rows whose aggregate control projection vanishes and whose aggregate constant is negative
--- the timing certificate in adjoint-row form. (iii) The program's feasible set is the
declared class itself, so the computation is sound and complete over \(\Pi_{B}\): no relaxation
gap arises. (iv) The declaration is load-bearing: the same computation over a polytope
\(U \supseteq U_{\mathrm{fin}}\) holding a finite implementable class computes the relaxation,
which remains sound but can be arbitrarily incomplete --- on the hidden-regime data with
\(U_{\mathrm{fin}} = \{-1, +1\}\) one has \(\sigma^{*} = z_{0} - 1\), while \(U = [-1,1]\)
admits \(u \equiv 0\), which holds both branches on the floor forever, so
\(\sigma^{*}_{\mathrm{rel}} = \infty\) and the relaxed certificate never fires.
\end{theorem}

\emph{Proof sketch.} Unrolling the affine dynamics makes each branch state
\(x_{j,k}\) affine in \((u_{0},\dots,u_{k-1})\), so every robust floor row is one linear
inequality once \(\max_{w \in W_{j}} F E_{j} A^{\ell} w\) is precomputed; feasibility of the
horizon-\(K'\) system is monotone in \(K'\) (prefix truncation), giving (i) by bisection, and
Farkas' lemma gives (ii). Claim (iii) is definitional: the program's variables are the declared
controls themselves. Claim (iv) is the computation displayed in
Example~\ref{calc-ex:relax-gap}; the complete proof, including the duality computation, is in
Supplementary S1.

\begin{example}[the relaxation gap is unbounded]\label{calc-ex:relax-gap}
Two branches share \(z^{+} = \tfrac{9}{10} z \pm u\) from a common \(z_{0}\), floor
\(\mathcal{V} = \{z \ge 1\}\), declared class \(U = [-1,1]\) per step. The branch sum
\(z^{+}_{k} + z^{-}_{k} = 2 (\tfrac{9}{10})^{k} z_{0}\) is invariant under the control, so both
floors force \(z_{0} \ge (\tfrac{10}{9})^{k}\) at every step \(k\): the window-\(K\) program is
feasible exactly when \(z_{0} \ge (\tfrac{10}{9})^{K}\). The level-\(k\) Farkas certificate is
the pair of floor rows summed --- their control coefficients cancel pairwise --- leaving
\(2(\tfrac{9}{10})^{k} z_{0} \ge 2\); and \(u \equiv 0\) attains feasibility on the boundary
\(z_{0} = (\tfrac{10}{9})^{K}\), keeping the branches balanced at
\((\tfrac{10}{9})^{K-k} \ge 1\). The program is thus a complete characterization of the declared
class. If instead the declared class is the finite set \(\{-1, +1\}\), the guaranteed survival
time is \(\sigma^{*} = z_{0} - 1\): strictly smaller. The relaxation enlarges the class and, on
the hidden-regime data, extinguishes the certificate altogether.
\end{example}

\begin{remark}[adjoint-row correspondence]\label{calc-rem:adjoint-duality}
The Farkas multipliers of (ii) are the finite analogue of the adjoint safety rows of the
continuous-to-finite bridge: a weight \(\lambda_{j,k,i}\) on the \(i\)-th floor row of branch
\(j\) at step \(k\) plays the role of an adjoint variable on a continuous-time safety row, the
vanishing aggregate control projection is the pooled adjoint condition, and the negative
aggregate constant is the certificate margin. The control coefficients of
Example~\ref{calc-ex:relax-gap} cancel pairwise within each level, which is the discrete image of
the pooled adjoint rows' cancellation. The bridge to continuous time --- error-controlled
sampling of these programs --- is developed in a companion work.
\end{remark}

\begin{figure}[htbp]
\centering
\includegraphics[width=0.82\columnwidth]{figs_p2/fig_p2_timing.png}
\caption{Delayed information. The constraint margin \(q\) reaches zero at
\(t^*=q_0/\varepsilon\), before the first informative observation at
\(T_{\mathrm{obs}}\): the information is accurate but arrives too late
(Theorem~\ref{calc-thm:delayed}).}
\label{calc-fig:timing}
\end{figure}

\subsubsection{3.4 The sparse finite
witness}\label{calc-the-sparse-finite-witness}

\begin{proposition}[sparse common-action witness]\label{calc-prop:helly}
Let \(B\subseteq\mathcal V\) be an information set, \(U(x)\subseteq U\)
with \(U\subseteq\mathbb R^m\) compact and convex and \(U(x)\) convex and
closed, let the dynamics be affine in the control,
\(f(x,u,d)=f_0(x,d)+f_u(x,d)\,u\), with \(D(x)\) compact, and let
\(\mathcal V\) have \(C^1\) constraint functions \(q_j\). Then each
safe-control set \(\mathcal R_{\mathcal V}(x)\) is a compact convex
subset of \(U\), and the safety-case common-action obstruction
\(\bigcap_{x\in B}\mathcal R_{\mathcal V}(x)=\varnothing\) holds if and
only if it is witnessed by at most \(m+1\) compatible states: there
exist \(x_1,\dots,x_{m+1}\in B\) with
\(\bigcap_{i=1}^{m+1}\mathcal R_{\mathcal V}(x_i)=\varnothing\). The
admissibility analogue (with \(\mathcal R_{\mathcal V}(x)\) replaced by
\(U(x)\)) holds under the same convexity assumptions.
\end{proposition}

\textbf{Proof.} For fixed \((x,j,d)\) the inequality
\(\nabla q_j(x)\cdot f(x,u,d)\ge 0\) is affine in \(u\), so
\(\mathcal R_{\mathcal V}(x)\), an intersection of closed halfspaces
(over active \(j\) and \(d\in D(x)\)) with the closed convex \(U(x)\),
is closed and convex; contained in the compact \(U\), it is compact. If
the intersection over all \(x\in B\) is empty, then by compactness some
finite subfamily of \(\{\mathcal R_{\mathcal V}(x)\}_{x\in B}\) already
has empty intersection, and by the finite Helly theorem at most \(m+1\)
members suffice; the converse direction is immediate. \hfill\(\square\)

\emph{Scope.} Both conclusions rest on convexity: the finite Helly theorem is what bounds the witness by \(m+1\), and beyond convex safe-control sets neither the bound nor the finite checkability of the certificate survives in general.

\textbf{A worked certificate (two floors, one scalar action).} Two
floors demand incompatible extraction rates on a single control
\(u\in[0,1]\). At the compatible state \(x_1\) (floor 1 active) the safe
controls are \(u\le 0.4\); at \(x_2\) (floor 2 active) they are
\(u\ge 0.6\). The observation merges \(x_1\) and \(x_2\), so a common
safe action must satisfy both. Writing the constraints as \(Au\le b\)
with \(A=(1,-1)^{\top}\) and \(b=(0.4,-0.6)^{\top}\), the stacked system
\(u\le 0.4\), \(-u\le -0.6\) is infeasible; the Farkas certificate is
\(\lambda=(1/2,1/2)\) (normalized, \(\|\lambda\|_1=1\)), for which
\(\lambda^{\top}A=0\) and \(\lambda^{\top}b=-0.1<0\). Under the normalization
\(\|\lambda\|_1=1\), the magnitude \(|\lambda^{\top}b|=0.1\) is the
infeasibility margin of the stacked system, shrinking to zero exactly as
the two safe-action sets come to touch. The multiplier
exhibits the contradiction as a convex combination of the two floor
constraints: no common extraction rate respects both floors, and any
observation that separates \(x_1\) from \(x_2\) restores one-step
viability (Figure~\ref{calc-fig:common-action}). With \(m=1\), the two
incompatible states are exactly the \(m+1\)-state witness of
Proposition~\ref{calc-prop:helly}.

\textbf{Scope of the convexity hypotheses.} Convexity is essential to the witness: without it the bound fails even in a single control dimension. Take \(U=\mathbb R\) (\(m=1\)) and three nonconvex safe-action sets \(\mathcal R_1=\{0,1\}\), \(\mathcal R_2=\{0,2\}\), \(\mathcal R_3=\{1,2\}\). Each pair intersects (\(\mathcal R_1\cap\mathcal R_2=\{0\}\), \(\mathcal R_1\cap\mathcal R_3=\{1\}\), \(\mathcal R_2\cap\mathcal R_3=\{2\}\)), yet the three-way intersection is empty, so no \((m+1)=2\)-state witness exists. The convexity hypotheses of Proposition~\ref{calc-prop:helly} exclude exactly this failure mode.

\textbf{Scope of the count: instantaneous actions, not window-policies.}
The \(m+1\) bound counts independent instantaneous decisions: it applies to a
common action \(u \in U \subseteq \mathbb{R}^m\), and it fails for common blind
control \emph{functions}. Partition a blind window into \(Q\) blocks
\(I_1, \dots, I_Q\) and let \(q + 1\) branches be indistinguishable throughout,
with scalar input (\(m = 1\)): branch \(j \le q\) evolves as
\(\dot x_j = \mathbf{1}_{I_j}(t)(u - 1)\) from \(x_j(0) = 0\) under the floor
\(x_j \ge 0\), which forces \(u \equiv 1\) on every block (each
\(w_j = \int_{I_j} u \, dt = 1\)), while branch \(q+1\) evolves as
\(\dot x_{q+1} = -u\) from \(x_{q+1}(0) = q - \varepsilon\) under the same
floor, forcing \(\sum_j w_j \le q - \varepsilon\). Summing the forced
withdrawals gives \(0 \le -\varepsilon\): the full system is infeasible, while
removing any single branch restores feasibility (\(u \equiv 1\), or \(u = 0\)
on the dropped branch's block and \(1\) elsewhere), so every obstruction
involves all \(q+1\) branches --- far more than \(m + 1 = 2\). The adjoint
kernels here are the block indicators, of rank \(q\): for window-policies the
governing count is the information--time rank, and sparsity claims for
review-interval and delayed-sensing certificates must be stated in the block
dimension \(mQ\), not the instantaneous dimension \(m\) (cf.\ the block
structure of the delayed-information check in Section~8).

\subsubsection{3.5 The finite-time exit
certificate}\label{calc-the-finite-time-exit-certificate}

The base obstruction operates before any observation-theoretic argument: it
certifies membership in the full-information ``shadow'' --- the
complement of the viability kernel (Aubin, 2001; Aubin and
Catt\'e, 2002).
It is a drift condition under which \emph{no} control --- of any
information structure --- can be viable. Put in management terms, when the dynamics themselves force exit before any control can react, no monitoring design can rescue viability, and the certification problem is closed without observation-theoretic reasoning.

\begin{theorem}[finite-time exit certificate]\label{calc-thm:exit}

Let \(q : X \to \mathbb{R}\) be a constraint function of class \(C^1\) --- the class assumed throughout the exit theorems --- with
\(\mathcal{V} \subseteq \{ q \ge 0 \}\), and suppose the following
hypotheses hold.

(H5.1) There exist constants \(a > 0\) and \(\varepsilon > 0\) such that
on the strip \(\mathcal{S}_a = \{ x : 0 \le q(x) \le a \}\),
\[\sup_{u \in U(x)} \inf_{d \in D(x)} D^+ q(x; f(x,u,d)) \;\le\; -\varepsilon \qquad \forall x \in \mathcal{S}_a, \tag{4}\]
where \(D^+ q(x; v) = \limsup_{h \downarrow 0} [q(x + hv) - q(x)]/h\) is
the upper right Dini derivative of \(q\) in direction \(v\) (for \(C^1\)
\(q\), \(D^+ q(x; v) = \nabla q(x) \cdot v\) --- the form the proof
uses).

(H5.2) \textbf{Closed-loop existence of the adverse realization.} The
disturbance correspondence is lower semicontinuous, or constant, on the
strip, so that the state-dependent adverse selection constructed in the
proof is realized by an admissible control--disturbance pair for which
the Carath\'eodory existence theorem applies. (Alternatively, without (H5.2), the certificate is read against the convexified (relaxed) inclusion
\(\dot x \in \mathrm{co}\, f(x, u(t), D_{\varepsilon}(x, u(t)))\), which
admits solutions under Marchaud-type conditions and preserves the drift
inequality because, for \(C^1\) \(q\), the directional derivative
\(\nabla q(x)\cdot v\) is affine in the velocity; closure and
relaxation are the standard tools for this step.)

Then, under hypotheses (H5.1)--(H5.2), for every
admissible control and every initial state \(x_0 \in \mathcal{S}_a\)
there exists an admissible disturbance realization under which the
trajectory attains \(q(x(t)) < 0\) --- hence leaves
\(\{ q \ge 0 \} \supseteq \mathcal{V}\) --- by every time strictly
larger than \(q(x_0)/\varepsilon\):
\(\inf\{ t : q(x(t)) < 0 \} \le q(x_0)/\varepsilon\), in particular
within time at most \(a/\varepsilon\) (since \(q(x_0) \le a\)).
(Convention, used in both exit theorems: \emph{violation} is the strict
inequality \(q < 0\); at \(q = 0\) the state may still lie on
\(\partial \mathcal{V} \subseteq \mathcal{V}\).)

\end{theorem}

\emph{Proof sketch.} Fix an admissible control. By (4), the adverse-selection correspondence \(D_\varepsilon(x,u)\) has nonempty values on the strip; its closed graph and compact values (Section 2.1) yield a measurable selection \(d(\cdot)\), realized in closed loop by (H5.2). The fundamental theorem of calculus integrates (4) along the resulting trajectory to \(q(x(t)) \le q(x_0) - \varepsilon t\), so \(q\) attains the strict violation \(q < 0\) by time \(q(x_0)/\varepsilon \le a/\varepsilon\).

\begin{remark}[comparison-function form]\label{calc-rem:comparison}
Hypothesis (H5.1) uses a constant drift margin \(\varepsilon\). The
certificate extends to state-dependent decline: if there is a continuous
\(\alpha:(0,a]\to(0,\infty)\) with
\[\sup_{u\in U(x)}\inf_{d\in D(x)} D^+q(x; f(x,u,d)) \le -\alpha(q(x)) \qquad \forall x\in\mathcal{S}_a,\]
then the exit-time bound sharpens to
\[t_{\mathrm{exit}} \le \int_0^{q(x_0)} \frac{ds}{\alpha(s)},\]
whenever the integral is finite. The proof replaces the linear comparison
\(q(x(t))\le q(x_0)-\varepsilon t\) with the solution \(r\) of
\(\dot r=-\alpha(r)\), \(r(0)=q(x_0)\), by the standard comparison
argument. If \(\alpha(0)=0\) and the integral diverges, the conclusion
weakens to asymptotic approach rather than finite-time exit; the
constant-margin case \(\alpha\equiv\varepsilon\) is the sharp finite-time
instance.
\end{remark}

Two readings of (4) matter. First, it is an \textbf{Isaacs-type
condition}: the disturbance chooses after the control, so the
certificate asserts that the disturbance can \emph{enforce} exit against
any policy. Second, (4) is an outward-drift certificate \emph{dual in
spirit, but not logically complementary}, to an inward barrier
condition. The pointwise robust barrier reading is
\(\sup_u \inf_d D^+ q \ge 0\) --- some control keeps the drift inward
against every disturbance --- and the two conditions leave an unresolved
intermediate region in which neither holds (Section 6.2). Theorem~\ref{calc-thm:exit} is
the ``unsafety'' side of that duality, and it is unconditional: it needs
no observation argument because it defeats all controls, including
clairvoyant ones.

\subsubsection{3.6 Epistemic emptiness by admissibility: a minimal
construction}\label{calc-epistemic-emptiness-by-admissibility-a-minimal-construction}

The admissibility obstruction applies to systems in which every state is
individually viable under full information and the kernel empties only under the observation structure.
The minimal construction of this section empties the kernel through
\emph{admissibility} --- a state-dependent control set whose common
action is forced and unsafe at the boundary; the common-action
obstruction of Section 3.2 and Example~\ref{calc-ex:hidden-mode}'s hidden-parameter
conflict are the informational cases proper, merging states with
incompatible safe controls. In ecological terms, each stock is manageable on its own; the failure arises because the indicator cannot tell the manager which stock is present.

\begin{proposition}[epistemic emptiness by admissibility --- a minimal construction]\label{calc-prop:emptiness}

There exists a system with
\(\mathrm{Viab}(\mathcal{V}; U, \pi_{\mathrm{perf}}) = \mathcal{V} \neq \varnothing\)
such that, for a non-injective observation map \(O\), every admissible
initial information state --- the observation fibres
\(B_0 = O^{-1}(O(S_0))\), \(S_0 \in \mathcal{V}\) --- lies outside
\(\mathrm{ERViab}_{\mathcal{I}}(\mathcal{V})\): the epistemic kernel
over the fibre-induced initial beliefs is empty. The mechanism is \emph{admissibility} rather than a hidden mode: the
state-dependent control set of the construction leaves a single common
action at the merged belief --- one that is outward-drifting at the
boundary state --- so the emptiness is the common-action obstruction
(Theorem~\ref{calc-thm:common-action}'s safety case) at minimal size, produced by the control-set primitive rather than by information; the
hidden-mode instance is Example~\ref{calc-ex:hidden-mode}.

\end{proposition}

\emph{Proof sketch (construction).} Take \(X = [1,2]\), \(\dot S = u - r(S)\) with \(r\) continuous and injective, \(U(S) = \{0, r(S)\}\), \(\mathcal{V} = [1,2]\), and the constant observation \(O \equiv 0\). Under full information \(u = r(S)\) gives \(\dot S = 0\), so every state is viable. Under \(O \equiv 0\) the common control set over the fibre is \(\{0\}\) --- injectivity of \(r\) rules out a common \(r(S)\) --- so the policy must hold \(u = 0\), and \(\dot S = -r(S) < 0\) drives every compatible trajectory out of \(\mathcal{V}\) in finite time.

The construction is the minimal instance of Theorem~\ref{calc-thm:common-action}'s admissibility mechanism --- with a constant control set in the same plant the fibre is epistemically viable --- so the empty-kernel phenomenon is epistemic rather than dynamic; the complete isolation of the mechanism is given in the Supplementary Material (S1).

\subsubsection{3.7 A recourse certificate for the post-observation
mode}\label{calc-recourse-certificate}

The certificates of the preceding subsections fire before information discriminates: the
common-action obstruction is instantaneous (Section 3.2) and the timing
obstruction exits within the blind window (Section 3.3, Remark
\ref{calc-rem:sigma}). Section 6.5(ii) isolates the complementary mode:
insufficient post-observation recourse. In that mode a common blind-window
control keeps every branch safe, yet no observation-adapted control can
recover every branch once the distinguishing observation arrives; the discrete
instance of Supplementary A.3 exhibits the mode while no certificate fires.
This subsection supplies a continuous-time certificate for the mode. Its
soundness covers every measurable information-adapted policy. Completeness is
not claimed, because the certificate relaxes the post-observation class
upward: an oracle grants each stored branch its own controller whether or not
the sensor reveals that much, and an oracle failure is only a sufficient
condition for real failure.

\begin{proposition}[oracle-recourse obstruction]\label{calc-prop:recourse}
Let the dynamics be affine in the control,
\(f(x,u,d) = f_0(x,d) + B(x)\,u\), with compact convex
\(U \subseteq \mathbb{R}^m\), and let branches
\(x_1, \dots, x_L \in B_0\) be observation-equivalent on \([0, \tau)\) with
horizon \(T > \tau\) (disturbances admitted, \(D\) compact). Fix a label
\(\ell = (\theta, j, t)\) --- branch \(\theta\), constraint \(j\), time
\(t \in (\tau, T]\) --- and a disturbance realization \(d(\cdot)\). Let
\(\bar x_\theta^d(\cdot)\) denote branch \(\theta\)'s zero-control
realization under \(d(\cdot)\), \(X_\theta^d(\cdot, \cdot)\) its
state-transition matrix, and \(B^d(\sigma) := B(\bar x_\theta^d(\sigma))\).
Define the input kernel
\[k_\ell^d(\sigma) = \mathbf{1}_{\{\sigma \le t\}}\,
\big(B^d(\sigma)\big)^{\top}
\big(X_\theta^d(t, \sigma)\big)^{\top}
\nabla q_j\big(\bar x_\theta^d(t)\big),\]
let \(\beta_\ell^d = q_j\big(\bar x_\theta^d(t)\big)\) denote the
zero-control facet margin at the label, and write
\(\phi_U(g) := \min_{v \in U} g^{\top} v\). For weights
\(\lambda_\ell \ge 0\), \(\sum_\ell \lambda_\ell = 1\), define the recourse certificate \[\begin{aligned} \Gamma_{h}(\lambda) = \inf_{d(\cdot)} \Big\{ &\sum_{\ell} \lambda_{\ell} \beta_{\ell}^{d} + \int_0^{\tau} h_U\Big(\sum_{\ell} \lambda_\ell k_\ell^d(\sigma)\Big) d\sigma \\ &\quad + \sum_{\ell} \int_{\tau}^{T} h_U\big(\lambda_\ell k_\ell^d(\sigma)\big) d\sigma \Big\}, \end{aligned}\] with \(h_U(g) := \max_{v \in U} g^{\top} v\) the support function of \(U\); the infimum runs over admissible disturbance realizations (for affine dynamics each bracket is computable in closed form by variation of constants). If \(\Gamma_{h}(\lambda) < 0\), then \(B_0 \notin \mathrm{ERViab}_{\mathcal{I}}(\mathcal{V})\): some stored facet is violated along some branch, for every policy in the stated class and every branch selection.
\end{proposition}

\emph{Proof sketch.} Fix a disturbance realization attaining the infimum and any admissible policy. Variation of constants expands the weighted facet sum at the stored labels into the \(\beta\)-term, the window pairing (one common measurable signal on \([0, \tau)\), by the blind-window reduction of Theorem~\ref{calc-thm:delayed}), and the branch pairings (separate post-reveal signals); each realized pairing is bounded above by the corresponding support value, so the weighted facet sum is at most \(\Gamma_{h}(\lambda)\). A strictly negative weighted sum forces some facet value strictly negative --- the facet is violated along its branch --- and the window signal was arbitrary, so every policy in the class fails. The complete proof, the comparison with the attainment-special \(\phi_{U}\)-variant, and the position of the pooling weights are given in the Supplementary Material (S1). \hfill\(\square\)

\begin{example}[three branches, delayed recourse]\label{calc-ex:three-branch}
A planar braking instance in which the obstruction is purely one of
post-observation recourse. Branches \(j = 1, 2, 3\) run \(\dot p = v\),
\(\dot v = u\) (no disturbance) from
\(p_j(0) = \tfrac{34}{25} n_j\), \(v_j(0) = n_j\), where
\(n_1 = (1, 0)\), \(n_2 = (-\tfrac{3}{5}, \tfrac{4}{5})\),
\(n_3 = (-\tfrac{3}{5}, -\tfrac{4}{5})\); the actuator set is
\(U = \mathrm{conv}\{\pm n_1, \pm n_2, \pm n_3\}\) (the hexagon
\(|u_2| \le \tfrac45\), \(|2u_1 \pm u_2| \le 2\)); safety is
\(n_i^{\top} p \le 2\) for \(i = 1, 2, 3\) with
\(\|v\|_\infty \le \tfrac65\); the realized branch is revealed at time
\(\tau\). Every branch is viable under full information --- brake along
\(-n_j\) for one second; peak critical position
\(\tfrac{34}{25} + \tfrac12 = \tfrac{93}{50} < 2\) --- and \(u \equiv 0\)
keeps every branch strictly safe through any blind window of length up to
\(\tfrac{16}{25}\) \((n_j^{\top} p_j(\tau) = \tfrac{34}{25} + \tau)\):
no instantaneous or blind-window certificate fires. Yet the label
\((j,\ n_j^{\top} p \le 2,\ \tau + 1)\) along branch \(j\) has
zero-control margin \(\beta_j = \tfrac{16}{25} - \tau - 1\) and kernel
\(\mathbf{1}_{\{\sigma \le \tau + 1\}} (\tau + 1 - \sigma) n_j\); the
weights \(\lambda = (\tfrac38, \tfrac5{16}, \tfrac5{16})\) pool the
pre-observation kernel to zero, and the certificate evaluates to \[\Gamma_{h}(\tau) = \Big(\tfrac{16}{25} - \tau - 1\Big) + \sum_{j} \lambda_{j} \int_{\tau}^{\tau+1} (\tau + 1 - \sigma)\, d\sigma = \tfrac{7}{50} - \tau,\] the pooled pre-observation kernel vanishing under \(\lambda\) and each post-reveal branch's pairing bounded by its support value \(\big(\sum_{j} \lambda_{j}/2 = \tfrac{1}{2}\big)\). For \(\tau > \tfrac{7}{50}\) the certificate fires and the triple is nonviable --- the blind window
accumulates braking liability that the post-observation authority cannot
discharge on every branch at once --- while the hold-then-brake policy
(\(u = 0\), then \(-n_j\) for one second) is safe exactly to
\(\tau = \tfrac{7}{50} = 0.14\): the threshold is exact. The obstruction
is genuinely three-way: at \(\tau = \tfrac15\), every pair of branches is
jointly viable under a common blind control (worst critical positions
\(\tfrac{2469}{1250} = 1.9752\) and \(\tfrac{2419}{1250} = 1.9352\), both
below \(2\)), so six of the seven nonempty priors over these branches are
viable at a delay where the triple is not. The discrete analogue in which no
certificate fires is Supplementary A.3; here the certificate fires, \(\Gamma_{h}(\tfrac15) = -\tfrac3{50} < 0\).
\end{example}

The certificate is deliberately asymmetric, like the calculus it joins: it is
sound against every measurable policy of the information class, its stored
labels are finite and independently checkable (exact arithmetic throughout the
instance above), and its incompleteness is explicit --- a complete
continuous-time characterization of the recourse mode remains open
(Section 6.5(ii)).

\subsubsection{3.7 A worked case: an exact obstruction on a continuum}
\label{calc-worked-case}

The calculus is motivated by the regime in which the epistemic kernel cannot be computed.
It is therefore worth exhibiting an instance in which a certificate returns an exact verdict
on a set of states that no finite procedure could enumerate. This subsection gives one,
verified in exact rational arithmetic.

\paragraph{The system.} A state is a pair of stocks \((x_1,x_2)\ge 0\) carrying a demand
floor on each coordinate. Let \(Y=x_1+x_2\) be the aggregate, and let the \emph{\(Y\)-fibre}
be the set of all states with that aggregate --- a continuum of states, since
\((x_1,x_2)\) ranges over the whole segment \(x_1+x_2=Y\), \(x_1,x_2\ge 0\). The control set
is \(U=\{u\ge 0:\ u_1+u_2\ge 2\}\), the allocations meeting total demand. Feasibility on the
fibre is constrained by \(Y\)-dependent caps
\[
\mathrm{cap}_1(Y)=\tfrac32-\tfrac{Y-2}{10},\qquad
\mathrm{cap}_2(Y)=\tfrac{59}{50}-\tfrac{Y-2}{10},
\]
an allocation \(u\in U\) being feasible on the fibre exactly when \(u_i\le\mathrm{cap}_i(Y)\)
for each \(i\). Fibre feasibility --- not the transition dynamics --- is the audited object
here, and the scope of the claim below is correspondingly static.

\paragraph{The certificate.} Viability of the fibre turns on a single rational inequality.
The cap sum is the affine function \(\tfrac{67}{25}-\tfrac{Y-2}{5}\), which equals \(2\) at
exactly one aggregate:
\[
Y^{*}=\tfrac{27}{5}.
\]
For \(Y\le Y^{*}\) the fibre is feasible; for \(Y>Y^{*}\) it is not, and the obstruction is
certified by the \(\ell_1\)-normalised dual measure \(\lambda=(\tfrac12,\tfrac12)\) on the two
cap rows, which for every \(u\in U\) gives
\[
\lambda^{\top}(\mathrm{cap}-u)\;\le\;\tfrac{\mathrm{cap}_1(Y)+\mathrm{cap}_2(Y)}{2}-1
\;=\;-\,\frac{Y-Y^{*}}{10}\;<\;0 .
\]
The certificate is therefore two rationals and one inequality, and its margin is exact:
\[
\mathrm{margin}(Y)=\frac{Y-\tfrac{27}{5}}{10}.
\]

\paragraph{The two audited fibres.} At \(Y=5\) the cap sum is \(\tfrac{52}{25}>2\) and the
fibre is feasible, with witness \(u=(\tfrac65,\tfrac45)\): it meets demand
(\(\tfrac65+\tfrac45=2\)) and respects both caps. At \(Y=6\) the cap sum is
\(\tfrac{47}{25}<2\), the fibre is obstructed, and \(\lambda\) certifies it with margin
exactly \(\tfrac{3}{50}\). Both statements, and the crossover, were recomputed independently
in exact rational arithmetic; the margin identity
\(\mathrm{margin}(Y)=(Y-\tfrac{27}{5})/10\) was checked at two hundred aggregates on the
half-line \(Y\ge Y^{*}\) with no violation.

\paragraph{What this establishes.} The \(Y=6\) fibre is a continuum of states, and the
question "is there an allocation meeting both floors from every state in this fibre?" is not
a question a finite enumeration could settle: there is no finite object to enumerate. The
certificate settles it exactly, with a rational witness and a rational margin, and by the
same argument settles it for \emph{every} \(Y\ge\tfrac{27}{5}\) simultaneously. That is the
regime the obstruction calculus is built for, and this is the instance that shows it
operating there.

\paragraph{Scope, stated plainly.} What is certified is a \emph{static} obstruction on the
observation fibre, not the full dynamic epistemic kernel: the audited object is fibre
feasibility, and no claim is made here about the kernel of the transition system induced by
these caps. The distinction matters and should not be blurred. The instance shows a
certificate returning an exact verdict on a continuum, which is the property the calculus
needs; it does not show, and does not claim to show, a certificate deciding a dynamic kernel
that is in principle undecidable. For the limits of what certificate pairs can decide in the
dynamic two-phase model, see Proposition~\ref{calc-prop:window-nogo}.

\subsubsection{3.8 Mechanized verification of the discrete core}
\label{calc-mech}

The results above are proved by hand. The discrete core of the calculus is
additionally machine-checked, and this subsection records exactly what that covers and
what it does not.

\paragraph{What was checked.} The finite-horizon theory, the kernel recursion, the
common-action obstruction, the certifier, and the monotonicity statements are formalized in Lean~4 in \texttt{Formalizations/P1\_Obstruction}, which imports only the
project prelude. The project builds cleanly: all \(60\) build jobs pass. The formalization comprises \(58\) declarations (\(33\) theorems and \(25\) definitions, excluding one \texttt{structure} and three \texttt{inductive} declarations), and the mapping onto the results above is
direct:

\begin{center}
\begin{tabular}{@{}ll@{}}
\toprule
Lean declaration & Result here\\
\midrule
\texttt{finite\_horizon\_sound}, \texttt{finite\_horizon\_complete} & Theorem~\ref{calc-thm:finite-horizon}\\
\texttt{Wk\_antitone}, \texttt{Wk\_descending} & Definition~\ref{calc-def:kernel}\\
\texttt{commonSafe}, \texttt{Blocked}, \texttt{blocked\_iff} & Theorem~\ref{calc-thm:common-action}\\
\texttt{preOp\_mono} & Proposition~\ref{calc-prop:monotone}\\
\bottomrule
\end{tabular}
\end{center}

\paragraph{No admitted gaps.} In Lean, an unfinished proof appears as the axiom
\texttt{sorryAx} in the axiom footprint of every declaration that depends on it. That
footprint was inspected for the declarations above. Several depend on no axioms at all
(\texttt{finite\_horizon\_sound}, \texttt{tree\_sound}, \texttt{Wk\_antitone},
\texttt{Wk\_descending}, \texttt{preOp\_mono}, \texttt{commonSafe}); the remainder depend
only on \texttt{Classical.choice}, \texttt{propext} and \texttt{Quot.sound}, the standard
axioms of Lean's classical logic. \texttt{sorryAx} appears nowhere.
\noindent\textbf{Verification provenance.} The project is pinned to Lean~4 in
\texttt{lean-toolchain} (currently \texttt{v4.34.1}). The mechanized checks reported in this
section were run under \texttt{v4.14.0}. On 2026-09-30 the source-level invariants were
re-verified against the current tree: a comment-stripped scan of all \texttt{.lean} files in
the project finds no \texttt{sorry}, \texttt{admit}, \texttt{axiom} or \texttt{constant}
declaration anywhere, and the declaration counts and named declarations quoted below were
re-counted and confirmed present. \textbf{The build has not been re-run since the toolchain
pin moved}, so the build-job figures below date from the \texttt{v4.14.0} run and should be
re-confirmed before submission.
 A source-level scan
with comments stripped finds no \texttt{sorry}, \texttt{admit} or \texttt{axiom}
declaration anywhere in the formalization layer; the only modifier present is
\texttt{noncomputable}, which marks a definition as using classical choice and is not a
gap.

\paragraph{Two scope limits, stated plainly.}

First, the formalization covers the \emph{discrete} core. The paper's continuous-time
statements --- the Dini-derivative arguments, the exit certificates, the timing bounds
--- live in an analysis setting that a dependency-free layer does not reach, and they are
\emph{not} formalized. The claim is that the discrete core is machine-checked, not that
the whole paper is.

Second, completeness is formalized with respect to policy \emph{trees}, not stationary
policies: \texttt{finite\_horizon\_complete} returns a \texttt{PolicyTree}, a
history-dependent object. Where the paper asserts that a stationary policy suffices,
that assertion rests on the hand proof, not on the formalization.

\paragraph{Fidelity to the text.} The formalization's source of record is version 53 of
this manuscript; the present version descends from version 57. Comparing the two, no
labelled result present in the earlier version has been dropped, and every result named
above is present in both. The formalization therefore applies to the results as they
stand here.

\subsection{4. Certification Limits}\label{calc-certification-limits}

\subsubsection{4.1 Exact certification and the fibre
criterion}\label{calc-exact-certification-and-the-fibre-criterion}

The obstruction theorems of Section 3 concern \emph{policies}.
Sustainability governance also uses \emph{certificates}: verdicts
--- safe or unsafe --- computed from observations alone, without
reference to a policy. The question is when such verdicts can be exact.
For a manager operating under partial ecological information, the question is whether an indicator reading alone can ever settle whether the system is on the safe side of a specified floor.

\begin{definition}[exact certifier]\label{calc-def:certifier}

For an admissible domain
\(Z \subseteq X\) and a safe set \(K \subseteq Z\), an \textbf{exact
certifier} based on the observation map \(O : Z \to Y\) is a function
\(C : Y \to \{0,1\}\) such that
\[C(O(z)) = 1 \iff z \in K \qquad \text{for every } z \in Z.\]

\end{definition}

\begin{proposition}[observation-fibre criterion]\label{calc-prop:fibre}

An exact certifier
based on \(O\) exists if and only if membership in \(K\) is constant on
every observation fibre:
\[O(z_1) = O(z_2) \;\Longrightarrow\; \big[ z_1 \in K \iff z_2 \in K \big] \qquad \forall z_1, z_2 \in Z, \tag{5}\]
equivalently, \(K = O^{-1}(O(K))\) on the admissible domain.

\end{proposition}

\emph{Proof sketch.} If \(C\) exists and \(O(z_1) = O(z_2)\), then \(\mathbf{1}_K(z_1) = C(O(z_1)) = C(O(z_2)) = \mathbf{1}_K(z_2)\), so membership is fibre-constant. Conversely, (5) makes \(C(y) = \mathbf{1}_K(z)\) for any \(z\) with \(O(z) = y\) well defined and exact.

The fibre criterion concerns static classification --- whether
membership in a safe set factors through the observation --- and applies
verbatim to three distinct certification problems, which should not be
conflated: constraint certification (\(K = \mathcal{V}\)), robust
viability certification (\(K = \mathrm{RViab}(\mathcal{V})\)), and
policy-specific safety certification (\(K = \mathrm{Safe}^{\Pi}_T\)).
The same theorem governs all three, but its governance meaning differs.
The criterion is the zero-step, label-valued analogue of the
common-action test: an observation-based verdict assigns a single label
to an entire fibre, exactly as an observation-based action assigns a
single action to an entire belief, so both are selector problems. They
remain distinct --- a fibre may fail exact certification while a policy
is still viable, and exact static certification does not imply dynamic
viability; hence certification is treated as a layer separate from
viability throughout.

\begin{corollary}[safety-crossing fibres and the certainly-safe set]\label{calc-cor:certainly-safe} If two admissible states share an observation and lie on opposite
sides of the \emph{same} component safety constraint --- with the safe
set of that certification problem read as the single floor
\(K = \{ q_j \ge 0 \}\); or, in the general form of the hypothesis, one
of the two states lies in \(K\) and the other outside \(K\) --- no exact
observation-only certificate exists. (Crossing one floor of an
intersection \(K = \bigcap_j \{ q_j \ge 0 \}\) does not by itself put
the pair on opposite sides of the intersection --- the state satisfying
floor \(j\) may violate floor \(i\) --- hence the single-floor scoping.)
The largest set of observations that can soundly be labeled safe without
further information is the certainly-safe set
\[\mathcal{Y}_{\mathrm{safe}} \;=\; \big\{ y \in O(Z) : O^{-1}(y) \subseteq K \big\}.\]

\end{corollary}

\emph{Proof sketch.} A fibre containing one state in \(K\) and one outside violates (5), so Proposition~\ref{calc-prop:fibre} denies the certifier. On \(\mathcal{Y}_{\mathrm{safe}}\) every compatible state is safe, so the constant verdict ``safe'' is sound; outside it some compatible state is unsafe, so no sound ``safe'' verdict exists.

The governance reading: a composite index whose fibres cross a floor boundary cannot support an exact certification of that floor, however carefully it is thresholded; the certainly-safe set is the sound region where the index may certify safety, outside which it must fall silent. The complete discussion, with the fibre figure, is given in the Supplementary Material (S1).

\subsubsection{4.2 The certainty-equivalence
trap}\label{calc-the-certainty-equivalence-trap}

The obstruction of Proposition~\ref{calc-prop:emptiness} relies on a non-injective observation ---
information lost. The same emptying occurs with a fully
\emph{injective} observation for a single policy: the canonical
certainty-equivalence controller that applies a perfect-information law
to the uncorrected observation.

\begin{remark}[the certainty-equivalence trap]\label{calc-rem:ce-trap}

Consider \(\dot S = u - g(S)\) with \(u \in [0, \bar u]\), \(g\) strictly
increasing on \([S_{\min}, S^* + b]\), \(g(0) = 0\), and
\(\mathcal{V} = [S_{\min}, S^*]\). Assume the range condition
\(g(S + b) \le \bar u\) for every \(S \in [S_{\min}, S^*]\), so that the
controller below is admissible. Under perfect information the zeroing
feedback \(u(t) = g(S(t))\) gives \(\dot S = 0\), so every
\(S_0 \in \mathcal{V}\) is viable and
\(\mathrm{Viab}(\mathcal{V}; U, \pi_{\mathrm{perf}}) = \mathcal{V} \neq \varnothing\).
Now take the injective, biased observation \(\hat S = S + b\) with
\(b > 0\), and the canonical certainty-equivalence controller: the
perfect-information zeroing law applied directly to the uncorrected
observation, \(u = g(\hat S) = g(S + b)\). Then
\[\dot S = g(S + b) - g(S) > 0 \qquad \forall S,\]
and since \(g\) is strictly increasing on the compact interval
\([S_{\min}, S^* + b]\), \(\dot S\) is bounded below by a positive
constant there; hence \(S\) strictly increases and exits above \(S^*\) in
finite time from every \(S_0 \in \mathcal{V}\): the kernel of this single
controller is empty. The bias-corrected controller
\(u = g(\hat S - b) = g(S)\) uses the observation map's structure, restores
\(\dot S = 0\), and keeps the kernel nonempty.

The mechanism is a \emph{failure of one policy to use the observation
map's structure}, not a loss of information: for a known invertible
observation the output-feedback class is as powerful as the
state-feedback class. The phenomenon is the viability-theoretic analogue
of Witsenhausen's counterexample in stochastic control (Witsenhausen,
1968): the information pattern --- not the dynamics and not the
constraint --- defeats the uncorrected controller. The remark is
policy-specific; it does not assert that the entire certainty-equivalence
class empties the kernel.

\end{remark}





\subsection{5. The Sufficiency Landscape}\label{calc-the-sufficiency-landscape-cited}

For contrast, we record the sufficiency results against which the obstruction calculus is defined; the proofs are in the cited sources.

\textbf{(a) Veliov's output-feedback condition.} Veliov (1993) gives a sufficient condition for the existence of an output-feedback regulation map under incomplete and inexact measurement, reducing under perfect measurement to the classical viability condition (Haddad, 1981). Veliov's theorem certifies existence; the obstruction calculus certifies nonexistence; a problem satisfying neither requires a stronger observer or a finer observation structure.

\textbf{(b) The estimation-tube programme.} Cardaliaguet, Quincampoix, and Saint-Pierre (2007) pass from imperfect information in the measurement space to perfect information in an estimation space of sets of compatible states, with equality of value functions and a Dini-derivative characterization. Under the set-membership semantics the estimation-space solution propagates the belief under a single control signal, so its controls are common controls; the common-action obstruction (Theorem~\ref{calc-thm:common-action}) is the certificate that no jointly admissible selection exists at the information state in question (Section 6.3).

\textbf{(c) Observer-and-buffer transfer.} If a full-information feedback exists with a strict inward margin and an observer supplies exponentially convergent estimates, output feedback of the estimate preserves safety on a buffered subset (observer-to-viability transfer with safety buffer; Sontag, 1998), and eroded kernels remain invariant under bounded estimation and implementation errors (erosion absorbs the error). Under margins and convergence, output feedback can work; Sections 3--4 state when it cannot.

\textbf{(d) The linear substitution alternative.} For a finite linear model in which a resource-typed system must meet a demand vector through specified substitution pathways, exactly one of the following holds (Farkas, 1902; Gale, 1960): (i) there exists a non-negative pathway vector \(a \ge 0\) satisfying the linear substitution constraints; or (ii) there exist multipliers \(\alpha, \beta, \gamma \ge 0\) with
\[
\begin{aligned}
&\alpha^\top R + \beta^\top E - \gamma^\top Q \ge 0 \quad\text{componentwise},\\
&\gamma^\top s^{\mathrm{req}} > \alpha^\top x + \beta^\top e.
\end{aligned}
\]
The second statement is a separation certificate that the pathways cannot meet demand within the typed resource and capacity bounds --- the feasibility-side complement of the observation-side obstructions of Sections 3--4.
\subsection{6. Discussion}\label{calc-discussion}


\subsubsection{6.1 The shape of the
calculus}\label{calc-the-shape-of-the-calculus}

The six mechanisms form a nonexhaustive taxonomy of information-theoretic failure. Table~\ref{calc-tab:summary} collects the certificates, the response correspondence each one empties, and the design consequence each one licenses.

\begin{table*}[t]
\centering
\footnotesize
\begin{tabular}{@{}l>{\raggedright\arraybackslash}p{3.4cm}>{\raggedright\arraybackslash}p{4.4cm}@{}}
\toprule
Certificate & Emptied response correspondence & Design consequence (Section 6.4) \\
\midrule
finite-time exit (Theorem~\ref{calc-thm:exit}) & the response set of some compatible state, even under full information & none --- dynamics, not information \\
admissibility (Proposition~\ref{calc-prop:emptiness}) & \(U^B(B)=\varnothing\) & enlarge the command set \\
common action and tube (Theorem~\ref{calc-thm:common-action}, Proposition~\ref{calc-prop:uniform-margin}) & \(\mathcal{R}_{\mathcal{V}}^B(B)=\varnothing\), \(\mathcal{A}_{\mathrm{tube}}(B,\Delta)=\varnothing\) & a separating observation; a shorter review interval \\
delayed information (Theorem~\ref{calc-thm:delayed}, Remark~\ref{calc-rem:sigma}) & \(\sigma^*(B_0)<T_{\mathrm{obs}}\) & an earlier informative observation \\
fibre certification (Proposition~\ref{calc-prop:fibre}, Corollary~\ref{calc-cor:certainly-safe}) & \(K\neq O^{-1}(O(K))\) & a refined index; per-floor measurement \\
certainty equivalence (Remark~\ref{calc-rem:ce-trap}) & kernel of the uncorrected controller \(u=g(\hat S)\) is empty & correct the known bias \\
finite-horizon recursion (Theorem~\ref{calc-thm:finite-horizon}) & \(B_0\notin\mathcal{W}_N\) & exact characterization (Section 3.1) \\
\bottomrule
\end{tabular}
\caption{The obstruction certificates, the response correspondence each one empties, and the design consequence each one licenses. The first six are sound sufficient conditions for nonviability; the last is the exact finite-horizon characterization.}
\label{calc-tab:summary}
\end{table*} The exit certificate (Theorem~\ref{calc-thm:exit}) is the
dynamic obstruction: it defeats all controls, so it needs no observation
argument. The epistemic-emptiness construction (Proposition~\ref{calc-prop:emptiness}) exhibits the
mechanism at minimal size --- one scalar state, two admissible controls,
one constant observation --- in the admissibility form of Section 3.6. The common-action obstruction (Theorem~\ref{calc-thm:common-action}) is the static
obstruction: it defeats all policies at a single information state. The
delayed-information obstruction (Theorem~\ref{calc-thm:delayed}) interpolates between them in
time. The fibre criterion (Proposition~\ref{calc-prop:fibre}) is the certification obstruction:
it limits not policies but verdicts. The certainty-equivalence trap
(Remark~\ref{calc-rem:ce-trap}) limits a single uncorrected policy. Finite-horizon completeness
(Theorem~\ref{calc-thm:finite-horizon}) supplies the exact finite-setting
counterpart of these sufficient certificates.

A governance failure of observation can be diagnosed by which
certificate applies. If the exit certificate holds, no monitoring design
helps --- the dynamics are adverse under every control. If only the
common-action obstruction holds, the remedy is a finer observation
structure (an informative measurement distinguishing the incompatible
states). If the delayed-information obstruction holds with a slack
timing bound, the remedy is a shorter review interval or a faster
indicator. If the fibre criterion fails, the remedy is an observation
that separates the sides of the floor rather than a better threshold on
the aggregate.

\subsubsection{6.2 Relation to barrier
certificates}\label{calc-relation-to-barrier-certificates}

Barrier certificates certify safety from a scalar function whose zero level set separates the unsafe region from all trajectories (Prajna and Jadbabaie, 2004; Prajna, Jadbabaie, and Pappas, 2007); the converse, under convex-duality conditions on density functions, is Prajna and Rantzer (2005), and necessary-and-sufficient hybrid characterizations are Maghenem and Sanfelice (2019). Theorem~\ref{calc-thm:exit} is the unsafety complement of that converse, not its observation-theoretic counterpart: a Dini-drift certificate of unsafety needing no observation structure, constructive in the enforcing disturbance and exit time. The remaining certificates (Proposition~\ref{calc-prop:emptiness}, Theorem~\ref{calc-thm:common-action}, Theorem~\ref{calc-thm:delayed}, Proposition~\ref{calc-prop:fibre}) act at the level of the information set itself: barrier constructions presuppose a state --- or a specified state estimate --- for feedback, whereas the common-action and timing obstructions certify the emptiness of the common safe-action set at a given information state, which barrier methods do not address.

\subsubsection{6.3 Relation to the estimation-tube
programme}\label{calc-relation-to-the-estimation-tube-programme}

The estimation-tube programme shows that imperfect measurement can be absorbed into set-valued dynamics on an estimation space with no loss in value (Cardaliaguet, Quincampoix, and Saint-Pierre, 2007); the present theorems delimit its reach rather than contradict it. The reduction preserves values, and its belief-state controls are common controls; the obstruction calculus is the local certificate layer of that programme --- the common-action obstruction certifies that the belief-state predecessor is empty, the timing bound quantifies when it fails for reasons of information arrival, and the fibre criterion states when no observation-based verdict can be exact. Belief-space viability answers \emph{whether}; the obstruction calculus answers \emph{why not, and what to change}. A single worked system carrying the three methods side by side --- barrier certificate, estimation-tube reduction, and obstruction certificate --- is given in the Supplementary Material (S2).

\subsubsection{6.4 Consequences for monitoring and indicator
design}\label{calc-consequences-for-monitoring-and-indicator-design}



Five consequences follow directly from the certificates. Each translates
a formal impossibility into a concrete design specification for
monitoring systems under partial ecological information.

\textbf{Timing.} Theorem~\ref{calc-thm:delayed} gives the minimal monitoring frequency in
closed form: the first informative observation must precede the enforced
exit time \(\inf q / \varepsilon\). A review interval longer than the
worst-case time from the most-constrained compatible state --- the
minimum-\(q\) state, which has the least constraint slack --- to
constraint violation is an information structure under which the
violation is undetectable in time, whatever the review then concludes.
The bound is computable from the drift certificate (2), which a robustness analysis computes in any case.

\textbf{Coarseness.} Theorem~\ref{calc-thm:common-action} states that an observation structure
merging two compatible states with incompatible safe controls is
nonviable \emph{at that information state}, even though both states are
viable in isolation. The design consequence is that the observation must
separate safety classes, not states: an indicator is exposed to Theorem~\ref{calc-thm:common-action}'s obstruction when some observation fibre crosses the safe-control
partition of the state space --- merging states whose safe controls differ is what the theorem certifies, and keeping every fibre within a single class is what removes that exposure; refinement is never harmful (a finer observation merges nothing the coarser one did not), though Theorem~\ref{calc-thm:common-action}
claims no benefit beyond the removal, and other mechanisms may still
certify nonviability on a finer partition.

\textbf{Aggregation.} Proposition~\ref{calc-prop:fibre} and Corollary~\ref{calc-cor:certainly-safe} apply verbatim to
composite indices, which are observation maps from the system state to a
scalar. An index whose fibres cross a floor boundary cannot support
exact floor certification; the certainly-safe set is the domain in which an index-based verdict is sound. Where separately-binding floors matter, exact
certification requires an observation that separates every admissible
pair on opposite sides of a floor (per-floor measurement is one sufficient design), a formal complement to the thesis, argued informally since Martinez-Alier, Munda, and O'Neill (1998), that weakly comparable
values cannot be fused into one metric without loss. The requirement is
sharpest where the floor constrains the productive base rather than its
output: an aggregate of services can read at full level while the base
regenerating them is reduced beneath it, and only an observation
reaching the base itself can detect the discrepancy.

\textbf{Bias.} Remark~\ref{calc-rem:ce-trap} shows that the loss need not be in the
measurement: a biased indicator with an uncorrected
certainty-equivalence policy empties the epistemic kernel of a system
whose perfect-information kernel is nonempty. Monitoring therefore has
two obligations --- bounding the observation error and supplying the
correction --- and the second lies entirely at the governance structure's discretion: the indicator is published, the correction is not applied,
and the kernel empties although the information was sufficient.

\textbf{Institutions.} The hierarchy of Section 2.1 records that
epistemic contraction is one of three causes of kernel shrinkage, the
others being disturbances and institutional restriction (authority,
enforcement, allocation). The epistemic-institutional kernel \(\mathrm{IRViab}_{\mathcal{J}}(\mathcal{V})\) (Section 2.1) combines them: an institution restricted in what it
may observe \emph{and} in what it may command inherits both
contractions. The obstruction calculus characterizes the first; its certificates remain valid, and typically tighten, when the institution's command set is further restricted, since restriction of
the action correspondence can only empty kernels, never create them
(one-sided monotonicity of viability under correspondence restriction).

\begin{proposition}[monotonicity of the epistemic kernel]\label{calc-prop:monotone}
The robust epistemic kernel is monotone in its constituents: enlarging
the action sets, reducing the disturbance sets, refining the information
structure, or enlarging the policy class never shrinks the kernel.
Writing \(\mathrm{ERViab}^{\,\cdot}\) for the kernel with the indicated
constituent, \(U_1\subseteq U_2\), \(D_1\subseteq D_2\),
\(\Pi_1\subseteq\Pi_2\), and ``\(\mathcal{I}_1\) refines
\(\mathcal{I}_2\)'' imply respectively
\[\mathrm{ERViab}^{U_1} \subseteq \mathrm{ERViab}^{U_2}, \qquad
\mathrm{ERViab}^{D_2} \subseteq \mathrm{ERViab}^{D_1},\]
\[\mathrm{ERViab}^{\Pi_1} \subseteq \mathrm{ERViab}^{\Pi_2}, \qquad
\mathrm{ERViab}^{\mathcal{I}_2} \subseteq \mathrm{ERViab}^{\mathcal{I}_1}.\]
\end{proposition}

\emph{Proof sketch.} A policy viable under the smaller action sets, larger disturbance sets, smaller policy class, or coarser information structure remains admissible and nonanticipative under the corresponding enlargement or refinement, while the adversary's options are unchanged or reduced; the same policy witnesses viability.

\subsubsection{6.5 Limitations}\label{calc-limitations}

The certificates are sound but not complete; several extensions are noted as open.

\begin{enumerate}
\def\labelenumi{(\roman{enumi})}
\tightlist
\item
  The continuous-time certificates of Sections 3.2--3.3 and 3.5--3.7 are sufficient conditions for nonviability, not necessary-and-sufficient characterizations of the epistemic kernel (backward recursion is complete in the finite-horizon finite-system setting, Section 3.1); the middle ground (neither a Veliov-type condition nor an obstruction certificate) is open. In
continuous settings completeness would additionally require a
measurable-selection theory beyond the explicit selections the
certificates exhibit; completeness is therefore claimed only where
backward recursion is exact (Section 3.1).
\item
  The timing obstruction of Theorem~\ref{calc-thm:delayed} (and its threshold form, Remark~\ref{calc-rem:sigma}) certifies failure by exit \emph{before} the first informative observation. It is blind to the complementary mode of nonviability --- \emph{insufficient post-observation recourse} --- in which a common blind-window control keeps every compatible branch safe, yet no observation-adapted control can recover every branch once the distinguishing observation arrives; the finite-horizon recursion (Section 3.1) captures that mode exactly, and the oracle-recourse certificate of Section 3.7 certifies it in continuous time --- soundly against every measurable information-adapted policy under an oracle relaxation of the post-observation class, with a worked three-branch instance whose delay threshold is exact; what remains open is a complete continuous-time characterization that does not relax the post-observation class. A discrete instance in which every one-step certificate is silent is given in the Supplementary Material (S3, A.3).
\item
  The information-set semantics is set-membership; probabilistic or belief-space formulations require separate statements, although the common-action obstruction transfers mutatis mutandis to any semantics in which the policy must choose a single action per information state.
\item
  The results are finite-dimensional and time-invariant; the delay-free setting makes the timing bound of Theorem~\ref{calc-thm:delayed} sharp, and delay systems require the retarded extension of the Dini comparison argument (Hale and Verduyn Lunel, 1993).
\item
  Theorem~\ref{calc-thm:exit}'s adversarial-realization step presupposes both the closed-graph regularity of Section 2.1 and the closed-loop existence hypothesis (H5.2) --- \(D\) lower semicontinuous or constant, or the convexified reading; without these the certificate remains a heuristic.
\item
  The institutional consequences of Section 6.4 are design conclusions from the certificates, not empirical findings; their empirical testing belongs to applied studies.
\end{enumerate}




\subsection{7. Complete Characterizations and the Residual Gap}\label{calc-complete-characterizations}

The certificates of Section 3 are sound sufficient conditions for nonviability, and Section 6.5 recorded that they do not exhaust the complement of the epistemic kernel. This section makes that gap precise: it exhibits two settings in which the viability question admits an exact answer --- the one-step (finite discrete) setting, where the common-action certificate is necessary and sufficient, and the static-observation setting, where the problem reduces exactly to open-loop robust viability --- and states the residual dynamic gap as an open problem.

\begin{theorem}[one-step completeness]\label{calc-thm:onestep}
Let \(X, A, D, Y\) be finite, with transition \(x^{+} = F(x,a,d)\), observation \(y = O(x,a,x^{+})\), and safe set \(\mathcal{V} \subseteq X\). Define the one-step safe set
\[
S_{1}(x) = \bigl\{ a \in U(x) : F(x,a,d) \in \mathcal{V} \text{ for all } d \in D(x) \bigr\}.
\]
A belief \(B \subseteq \mathcal{V}\) is one-step viable under \(\mathcal{I}\) --- some observation-based policy achieves \(x_{1} \in \mathcal{V}\) for every \(x_{0} \in B\) and every admissible disturbance --- if and only if \(\bigcap_{x \in B} S_{1}(x) \neq \varnothing\). Hence \(B \notin \mathcal{W}_{1}\) exactly when the common safe-action set is empty, the discrete reading of Theorem~\ref{calc-thm:common-action}.
\end{theorem}

\emph{Proof.} \((\Rightarrow)\) An observation-based policy depends on the record only, and every \(x \in B\) produced the same record, so the policy plays a single \(a \in U^{B}(B)\); safety of the first step forces \(F(x,a,d) \in \mathcal{V}\) for all \(x \in B\) and \(d \in D(x)\), i.e.\ \(a \in \bigcap_{x \in B} S_{1}(x)\). \((\Leftarrow)\) Playing any \(a\) in the intersection keeps every successor in \(\mathcal{V}\). \hfill\(\square\)

Combined with the backward recursion \(\mathcal{W}_{k+1} = \mathrm{Pre}(\mathcal{W}_{k})\) of Theorem~\ref{calc-thm:finite-horizon}, Theorem~\ref{calc-thm:onestep} shows that the pair \emph{(common-action certificate, recursion)} is a complete characterization at every finite horizon in finite systems: the certificate detects first-step emptiness, and the recursion propagates the obstruction to every later step.

\begin{theorem}[static-observation reduction to open-loop viability]\label{calc-thm:static-complete}
Let \(\mathcal{V}\) be compact with \(C^{1}\) constraint functions \(q_{j}\), \(f\) continuous with \(D\) compact-valued and of closed graph, \(U(x) \equiv U\), and suppose the observation is \emph{static}: \(y(t) = O(x(t))\) is constant along every admissible trajectory, so no information accrues over time. Then every observation-based policy coincides with an \emph{open-loop} control --- a measurable function of time only, \(u(\cdot):[0,\infty)\to U\) --- and \(B_{0} \in \mathrm{ERViab}_{\mathcal{I}}(\mathcal{V})\) if and only if there exists an open-loop control under which every trajectory from every \(x_{0} \in B_{0}\) and every admissible disturbance realization stays in \(\mathcal{V}\) for all time. Within the constant-control subclass \(u(\cdot) \equiv u_{0}\), the condition is the robust Nagumo boundary condition
\[
\nabla q_{j}(x) \cdot f(x, u_{0}, d) \;\ge\; 0
\qquad \text{for all } x \in \partial\mathcal{V},\ \text{active } j,\ d \in D(x),
\]
necessary for a constant control to keep \(\mathcal{V}\) invariant and sufficient under the standard regularity assumptions (Aubin, 1991; Frankowska, 1989). A time-varying open-loop control may be viable when no constant control is; the open-loop reduction, not the constant-control condition, is the exact characterization in the static class.
\end{theorem}

\emph{Proof.} A static observation makes the record constant, so a policy's action at time \(t\) can depend on time only --- open-loop; this gives the reduction in both directions (a viable policy yields a viable open-loop control, and every open-loop control is an admissible observation-based policy). The constant-control statement is the robust Nagumo theorem (Aubin, 1991; Frankowska, 1989): \emph{necessary}, because any trajectory leaving \(\mathcal{V}\) under \(u_{0}\) exhibits a boundary point with outward drift, and \emph{sufficient} under the regularity assumptions. \hfill\(\square\)

Theorem~\ref{calc-thm:static-complete} makes the sense of completeness in the static class precise: the obstruction is exactly the nonexistence of a robustly viable open-loop control, and within the constant-control subclass the boundary condition of Theorem~\ref{calc-thm:common-action} is necessary and sufficient. Completeness fails only when information accrues, where the residual modes are post-observation recourse and timing --- the hidden-mode example and the timing bound of Section 3.3.

\medskip
The static-observation theorem localizes where completeness can fail, and the failure
admits a precise structure. Consider the \emph{two-phase model}: a window of \(K\) blind
steps followed by exact branch revelation and branch-adapted continuation, with declared
blind class \(\Pi_{B}\) and full-information kernel \(\mathrm{RViab}(\mathcal{V})\)
available after the reveal.

\begin{proposition}[exact two-phase decomposition]\label{calc-prop:decomposition}
In the two-phase model, \(B_{0}\) is viable if and only if (i) some declared blind control
keeps every branch in \(\mathcal{V}\) through step \(K\), and (ii) some declared blind
control satisfying (i) lands every branch in \(\mathrm{RViab}(\mathcal{V})\) at the reveal.
Consequently, over any declared class on which the window certificates are complete ---
finite classes; polytope-declared classes, by Theorem~\ref{calc-thm:lp-instant} --- the following
are equivalent: both certificates are silent; clause (i) holds; and either \(B_{0}\) is
viable or its nonviability is exactly the post-observation recourse mode, the failure of
(ii) for every window-surviving blind control.
\end{proposition}

\emph{Proof sketch.} A policy for the two-phase model is a blind part followed by an
adapted part. After the reveal the information is complete, so an adapted continuation
exists from a revealed state if and only if that state lies in
\(\mathrm{RViab}(\mathcal{V})\); and the blind part succeeds only branchwise, since a branch
leaving \(\mathcal{V}\) before the reveal, or landing outside the kernel at it, kills
viability upon that branch's realization. This gives the equivalence. Completeness of the
window certificates turns silence into clause (i), and the residual failure of
clause (ii) is the recourse mode; the complete proof is in Supplementary S1.

\begin{proposition}[no window-measurable pair is complete]\label{calc-prop:window-nogo}
Call a certificate pair \emph{window-measurable} when its two verdicts depend only on the
window sub-model: the initial belief, the branch transitions within the window, the floors,
and the observation schedule. Then no window-measurable pair is complete for the two-phase
model: the three-state instance of Supplementary S3/A.3 and its variant agree on all window
data, the certificates are silent in both, and yet one is nonviable (\(x_{4}\) exits under
every post-reveal action) while the other is viable (\(x_{4}\) is maintained).
\end{proposition}

\emph{Proof sketch.} The two instances differ only in \(x_{4}\)'s post-reveal transitions,
so any window-measurable pair returns identical verdicts --- here, silence on both. The
exact recursion decides differently: in the first instance every post-reveal action exits
\(\mathcal{V}\), so \(x_{4} \notin \mathrm{RViab}(\mathcal{V})\) and the belief misses
\(\mathcal{W}_{2}\); in the second \(x_{4}\) is maintained, lies in the kernel, and the
belief is viable. Both computations are machine-verified in exact rational arithmetic
(Supplementary S1, verification record). The residual dynamic gap of
Open Problem~\ref{calc-op:dynamic} is thereby narrowed to a definite target: any complete
certificate pair for the dynamic class must read data beyond the window sub-model.

\begin{openproblem}[dynamic-observation completeness]\label{calc-op:dynamic}
Give a finite, checkable certificate that is necessary and sufficient for membership in \(\mathrm{ERViab}_{\mathcal{I}}(\mathcal{V})\) under dynamic observation beyond the belief-space restatement below, or prove a separation showing that none exists in a natural class. The belief-space restatement is: \(B_{0} \in \mathrm{ERViab}_{\mathcal{I}}(\mathcal{V})\) if and only if \(B_{0}\) belongs to the viability kernel of the set-valued belief dynamics under the observation constraint (Cardaliaguet, Quincampoix, and Saint-Pierre, 2007); this is a complete characterization on the infinite-dimensional space of compact beliefs, and the obstruction calculus of Section 3 is its finite, certifiable layer.
\end{openproblem}

\subsection{8. The Selector Principle, Dual Certificates, and Monitoring Adequacy}\label{calc-selector-principle}

The certificates of Section 3 are stated mechanism by mechanism. This section records the one object they consolidate, the exactness class in which its emptiness is also complete, and the design converse; it adds no new certificate and modifies none.

\paragraph{The selector.} For a declared policy class, information structure, review length \(\Delta\), and horizon \(N\), the recursively viable selector at information state \(B\) is
\[
\begin{split}
\Gamma_{N}(B) = \bigl\{\alpha :&\ \mathrm{Reach}_{[0,\Delta]}(B,\alpha) \subseteq \mathcal{V} \\
&\text{and every successor belief } B' \in \mathcal{W}_{N-1}\bigr\},
\end{split}
\]
with \(\mathcal{W}_{0} = \{B \subseteq \mathcal{V}\}\) and \(\mathcal{W}_{N} = \{B : \Gamma_{N}(B) \neq \varnothing\}\) --- the recursion of Theorem~\ref{calc-thm:finite-horizon} --- whose one-step factors are the ladder of Proposition~\ref{calc-prop:ladder}. The principle: viability within a declared class is nonemptiness of this selector, and each certificate of Section 3 is a checkable witness that one factor of it is empty.

\paragraph{Dual witnesses.} The witnesses, by factor: the common-action factor \(\mathcal{R}^{\mathcal{B}}_{\mathcal{V}}(B) = \varnothing\), witnessed by the Farkas pair of Section 3.2 and, under convexity, by the \(m+1\)-point sparse witness of Proposition~\ref{calc-prop:helly}; the tube factor, witnessed by the uniform-margin review obstruction of Proposition~\ref{calc-prop:uniform-margin}; the timing factor, witnessed by the threshold form \(\sigma^{*}(B_{0}) < T_{\mathrm{obs}}\) of Remark~\ref{calc-rem:sigma} and Theorem~\ref{calc-thm:delayed}; finite-time exit, witnessed by the outward drift of Theorem~\ref{calc-thm:exit} already at singletons; epistemic emptiness, witnessed by the minimal admissibility construction of Proposition~\ref{calc-prop:emptiness}; certification failure, witnessed by a fibre meeting both \(K\) and \(K^{c}\), with the certainly-safe set of Corollary~\ref{calc-cor:certainly-safe} (Proposition~\ref{calc-prop:fibre}); and class restriction, witnessed by the certainty-equivalence trap of Remark~\ref{calc-rem:ce-trap} --- the uncorrected law's strictly positive drift against the corrected law's nonpositive drift.

\paragraph{Exactness class.} Weak duality holds unconditionally: a dual witness certifies nonviability in the declared class. Strong duality --- emptiness if and only if a witness exists --- holds exactly in the classes of Theorem~\ref{calc-thm:onestep} (finite horizon, finite model), Theorem~\ref{calc-thm:static-complete} (static observation), and the delayed class of Remark~\ref{calc-rem:coverage}; beyond them the converse is Open Problem~\ref{calc-op:dynamic}, open. The causes are not unified: dynamic exit empties the selector under full information and is not epistemic; the fibre results are limits on certification, not on policy viability (Section 4); the trap is a class restriction and can fire under injective observation. Unifying the form while keeping the causes distinct is the strongest unifying statement the framework supports, and the companion unification statement carries the same discipline.

\paragraph{Monitoring adequacy (the design converse).} Obstruction answers why nonviability is forced; design asks what must be measured. In the one-step exact class the converse is immediate: an observation structure is adequate on a target exactly when every fibre meeting the target has a nonempty common safe action --- the partition form of the fibre criterion (Proposition~\ref{calc-prop:fibre}); pairwise intersection is necessary and, off convexity, not sufficient (the three two-point codices of the companion monitoring paper). Four rules follow as identities on the audited systems. (i) \emph{Timing:} on the delayed hidden-regime hold class, viability at deadline \(T_{\mathrm{obs}}\) holds exactly when \(T_{\mathrm{obs}} \le \sigma^{*}(B_{0}) = z_{0} - 1\), equivalently \(z_{0} \ge 1 + T_{\mathrm{obs}}\); the boundary case \(z_{0} = 1 + T_{\mathrm{obs}}\) is viable --- the review must land no later than worst-case exit (Section~\ref{calc-case-study}'s grid; the strict form fails at the boundary cell, as the verification script records). (ii) \emph{Coarseness:} a fibre merging states with disjoint safe-action sets obstructs at that information state while every branch remains viable (Theorem~\ref{calc-thm:common-action}); the remedy is a fibre split, not a finer threshold on the same index. (iii) \emph{Aggregation:} an index whose fibres cross a floor cannot support an exact safe/unsafe verdict; report the certainly-safe set or fall silent (Corollary~\ref{calc-cor:certainly-safe}). (iv) \emph{Bias:} the trap of Remark~\ref{calc-rem:ce-trap} empties a nonempty kernel under the uncorrected law; the remedy is structural correction of the observation map, not a better sensor. The consolidated statement, with every audited count cited to its system, is maintained in the supplementary material.

\subsection{9. A Calibrated Case Study and a Coverage Audit}\label{calc-case-study}

We instantiate the three observation-side mechanisms on a concrete resource model and audit the certificates against the exact recursion of Theorem~\ref{calc-thm:finite-horizon}. The model is a single-species stock \(S \in [0,1]\) (fraction of carrying capacity) with discrete dynamics
\[
\begin{gathered}
S^{+} = \mathrm{clip}\bigl(S + S(1-S) - u + w\bigr),\\
u \in [0, 0.3],\quad w \in \{-0.02, 0, +0.02\},
\end{gathered}
\]
harvest \(u\), recruitment noise \(w\) read adversarially, and a sustainability floor \(S_{\min} = 0.3\) (Clark, 1990).

\textbf{The fibre criterion.} The regulator observes a coarse two-bin index, low if \(S < 0.6\) and high if \(S \ge 0.6\). The safe set \(\mathcal{V} = [0.3, 1]\) is not a union of index bins: the low bin contains both \(0.1\) (below the floor) and \(0.4\) (above it). By Proposition~\ref{calc-prop:fibre} no exact observation-only certifier exists, and by Corollary~\ref{calc-cor:certainly-safe} the certainly-safe set is the single bin \(\mathcal{Y}_{\mathrm{safe}} = \{\text{high}\}\). A verdict of ``safe'' on the strength of a low reading is therefore unsound; the index can certify safety only when it reads high.

\textbf{Two-dimensional aggregation.} The fibre criterion also governs a multidimensional stock. Take a two-species state \((S_{1}, S_{2}) \in [0,1]^{2}\) with the floor \(S_{1} \ge 0.4\), and suppose the regulator observes only the aggregate \(I = S_{1} + S_{2}\). The states \((0.5, 0.5)\) and \((0.2, 0.8)\) are both compatible with the reading \(I = 1.0\), yet the first satisfies the floor and the second violates it, so by Proposition~\ref{calc-prop:fibre} no exact observation-only certifier exists; by Corollary~\ref{calc-cor:certainly-safe} the certainly-safe readings are exactly \(I \ge 1.4\) (the smallest \(S_{1}\) compatible with \(I\) is \(I - 1\), so \(I - 1 \ge 0.4\)). An aggregate reading below \(1.4\) must therefore fall silent, however carefully thresholded --- the two-dimensional form of the aggregation consequence of Section 6.4.

\textbf{The delayed-information timing bound.} The timing obstruction is sharpest in the hidden-regime system of Section 3.3: regimes \(\theta \in \{-1, +1\}\), dynamics \(z^{+} = z + \theta u\), controls \(u \in \{-1, +1\}\), floor \(z \ge 1\), with \(\theta\) revealed only after a blind window of \(T_{\mathrm{obs}}\) steps. From \(z_{0} = 2\) the drift of \(q(z) = z - 1\) on the declining branch is \(-1\) under every admissible action, so the threshold of Theorem~\ref{calc-thm:delayed} and Remark~\ref{calc-rem:sigma} is \(\sigma^{*}(B_{0}) = z_{0} - 1 = 1\): nonviability is certified for every \(T_{\mathrm{obs}} > 1\), while \(T_{\mathrm{obs}} = 1\) is viable --- the regime is revealed in time, and \(u = \theta\) then restores growth. In the fishery above the same certificate is silent in the normal regime, because \(u_{0} = 0\) keeps every stock above the floor; the bound fires only when no blind action avoids decline, which is precisely its scope.

\textbf{The certainty-equivalence trap.} With a biased index \(\hat{S} = S + 0.1\), \(g(S) = S^{2}\), and \(\mathcal{V} = [1,2]\), the certainty-equivalence policy \(u = g(\hat{S})\) produces \(\dot{S} = g(S+0.1) - g(S) = 0.2S + 0.01 > 0\), which reaches the upper bound from \(S = 1\) in \(\approx 3.3\) time units, while the bias-corrected policy \(u = g(\hat{S} - 0.1)\) keeps \(\dot{S} = 0\) and the kernel nonempty (Remark~\ref{calc-rem:ce-trap}).

\textbf{A reproducible delayed hidden-regime model.} The audit runs on the exact instance specified here. The state is \((z, \theta)\) with \(\theta \in \{-1, +1\}\) a hidden regime and \(z \ge 0\); the action \(u \in \{-1, +1\}\) is common to both regimes; the transition is \(z^{+} = z + \theta u\); the safe set is \(\mathcal{V} = \{ z \ge 1 \}\); and the observation reveals nothing before \(T_{\mathrm{obs}}\) and reveals \(\theta\) exactly at time \(T_{\mathrm{obs}}\). From a known \(z_{0}\) with both regimes possible, the belief after \(k < T_{\mathrm{obs}}\) blind steps contains the two extreme histories \((z_{0} - k, -1)\) and \((z_{0} + k, +1)\), and a blind policy must choose the same \(u\) for both branches at every blind step. The guaranteed blind-window survival time is therefore \(\sigma^{*}(B_{0}) = z_{0} - 1\): the best blind policy drives one branch down at unit rate, so the declining branch reaches the floor \(1\) at time \(z_{0} - 1\), and after the reveal \(u = \theta\) keeps \(z \ge 1\) forever. Hence \(B_{0}\) is robustly epistemically viable if and only if \(z_{0} - 1 \ge T_{\mathrm{obs}}\), i.e.\ \(z_{0} \ge 1 + T_{\mathrm{obs}}\) --- the ground truth against which the certificates are audited. The one-step certificate of Theorem~\ref{calc-thm:common-action} fires exactly when \(z_{0} < 2\) (one blind step already forces \(z_{0} - 1 < 1\)), and the timing bound of Theorem~\ref{calc-thm:delayed} fires exactly when \(z_{0} \ge 2\) and \(T_{\mathrm{obs}} > z_{0} - 1\).


\textbf{A coverage audit.} To quantify the gap identified in Section 6.5, we audit the certificates against the exact recursion on the delayed hidden-regime system, over the grid \(z_{0} \in \{1.0, 1.1, \dots, 2.5\}\), \(T_{\mathrm{obs}} \in \{1,2,3\}\). Table~\ref{calc-tab:coverage} and Figure~\ref{calc-fig:coverage} report the true verdict, the one-step certificate, and the timing bound. The audit is reproduced by the script \texttt{paper2\_coverage\_audit.py}, publicly available at \url{https://github.com/MIKEAA2020/general-sustainability} (folder \texttt{arena agent 1/paper rewrites/latex}), which regenerates Table~\ref{calc-tab:coverage} and Figure~\ref{calc-fig:coverage} verbatim.

\begin{table*}[t]
\centering
\footnotesize
\begin{tabular}{@{}lcccccccccccccccc@{}}
\toprule
\(T_{\mathrm{obs}}\) & \multicolumn{16}{c}{initial position \(z_{0}\)} \\
\cmidrule(lr){2-17}
 & 1.0 & 1.1 & 1.2 & 1.3 & 1.4 & 1.5 & 1.6 & 1.7 & 1.8 & 1.9 & 2.0 & 2.1 & 2.2 & 2.3 & 2.4 & 2.5 \\
\midrule
1 & \(\bullet\) & \(\bullet\) & \(\bullet\) & \(\bullet\) & \(\bullet\) & \(\bullet\) & \(\bullet\) & \(\bullet\) & \(\bullet\) & \(\bullet\) & \(\circ\) & \(\circ\) & \(\circ\) & \(\circ\) & \(\circ\) & \(\circ\) \\
2 & \(\bullet\) & \(\bullet\) & \(\bullet\) & \(\bullet\) & \(\bullet\) & \(\bullet\) & \(\bullet\) & \(\bullet\) & \(\bullet\) & \(\bullet\) & \(\ast\) & \(\ast\) & \(\ast\) & \(\ast\) & \(\ast\) & \(\ast\) \\
3 & \(\bullet\) & \(\bullet\) & \(\bullet\) & \(\bullet\) & \(\bullet\) & \(\bullet\) & \(\bullet\) & \(\bullet\) & \(\bullet\) & \(\bullet\) & \(\ast\) & \(\ast\) & \(\ast\) & \(\ast\) & \(\ast\) & \(\ast\) \\
\bottomrule
\end{tabular}
\caption{Coverage audit on the delayed hidden-regime system: \(\circ\) viable; \(\bullet\) nonviable, certified by the common-action certificate (Theorem~\ref{calc-thm:common-action}); \(\ast\) nonviable, certified only by the timing bound (Theorem~\ref{calc-thm:delayed}). Every nonviable cell is covered by exactly one certificate.}
\label{calc-tab:coverage}
\end{table*}

\begin{figure*}[t]
\centering
\includegraphics[width=0.82\linewidth]{figs_p2/fig_p2_coverage.png}
\caption{The coverage audit of Table~\ref{calc-tab:coverage}: green = viable; red = nonviable, certified by the common-action certificate; orange = nonviable, certified only by the timing bound (the residual gap of Section 6.5). The dashed curve is the true boundary \(z_{0} = 1 + T_{\mathrm{obs}}\) and the solid vertical line the one-step boundary \(z_{0} = 2\).}
\label{calc-fig:coverage}
\end{figure*}

\begin{remark}[joint completeness in the delayed class]\label{calc-rem:coverage}
In every audited cell the true verdict is reproduced by exactly one of the two certificates: of the 42 nonviable cells, 30 are certified by the common-action certificate and the remaining 12 --- precisely those with \(z_{0} \ge 2\) and \(T_{\mathrm{obs}} > z_{0} - 1\) --- by the timing bound; none is uncovered. In this class the two certificates are jointly complete, and the Section 6.5 gap reduces to the timing cells.
\end{remark}

\emph{Why joint completeness holds in this class.} The ground truth is the identity \(B_{0} \in \mathrm{ERViab}_{\mathcal{I}}(\mathcal{V}) \iff z_{0} \ge 1 + T_{\mathrm{obs}}\) established in the reproducible-model paragraph above, so nonviability holds exactly when \(z_{0} < 1 + T_{\mathrm{obs}}\). The two certificates partition that set: the common-action certificate fires exactly when \(z_{0} < 2\), and the timing bound fires exactly when \(z_{0} \ge 2\) and \(T_{\mathrm{obs}} > z_{0} - 1\), that is \(2 \le z_{0} < 1 + T_{\mathrm{obs}}\); the two conditions are disjoint and their union is \(\{z_{0} < 1 + T_{\mathrm{obs}}\}\). Every nonviable cell is therefore certified by exactly one certificate, so within the delayed hidden-regime class the pair \(\text{(common-action, timing)}\) is a complete characterization at every horizon, and the gap of Section 6.5 is exactly the timing cells.

\textbf{Calibration.} The certificate parameters correspond to quantities estimable from standard stock-assessment practice: the drift margin \(\varepsilon\) is a net growth rate at the floor; the first informative observation time \(T_{\mathrm{obs}}\) is the assessment or review periodicity; the bias \(b\) of the certainty-equivalence instance is the systematic error of the indicator; and the floor is the reference point of the management rule. Calibrating these four numbers from a published assessment is all the machinery the certificates require; the present section stays symbolic, and the empirical calibration is left to applied studies.

\textbf{A two-patch protection audit.} The aggregation obstruction is not a one-dimensional phenomenon, and the certificates read it directly in two dimensions. Two patches \(i \in \{1,2\}\) carry stocks \(z_{i} \in \{0,1,2,3\}\), each with recruitment \(+1\) per step and a protection resource that covers one patch per step: the action \(u \in \{1,2\}\) protects patch \(u\), the protected patch loses nothing, and the other loses two units, \(z_{i}^{+} = \min\{3,\, z_{i} + 1 - 2\cdot \mathbf{1}[u \neq i]\}\). The sustainability floor is \(z_{i} \ge 1\) for both patches, and the regulator observes only the aggregate reading \(y = z_{1} + z_{2}\). From the reading \(y = 3\) the compatible safe states are \((1,2)\) and \((2,1)\); the first is safe only under \(u = 1\) and the second only under \(u = 2\), so the common safe-action set is empty and the common-action certificate (Theorem~\ref{calc-thm:common-action}) fires --- while each state alone lies in the full-information kernel (protect the low patch, then alternate), the minimal epistemic-emptiness form of Section 3.4 realized in two dimensions. The certainly-safe readings are exactly \(y \ge 4\): every state compatible with \(y \le 3\) includes one with a depleted patch, so the aggregated index can certify safety only above that threshold --- the two-dimensional reading of Corollary~\ref{calc-cor:certainly-safe}. Table~\ref{calc-tab:patch} audits the aggregated beliefs against the exact belief-level recursion, in which the reading splits the belief into its \(y\)-cells at every step: every nonviable belief is certified, and all by the common-action certificate --- the aggregation obstruction is instantaneous, and no observation schedule rescues it, in contrast with the timing cells of the one-dimensional audit above, where \(T_{\mathrm{obs}} = 1\) is viable. The full-information kernel is \(\mathrm{RViab}(\mathcal{V}) = \{z_{1}, z_{2} \ge 1\} \setminus \{(1,1)\}\): the double-depleted state is nonviable outright, one protection per step cannot hold both floors --- an instance of nonviability beneath all epistemic analysis. Both computations are machine-verified in exact integer arithmetic, and the verification script regenerates the table verbatim.

\begin{table}[t]
\centering
\resizebox{\columnwidth}{!}{\begin{tabular}{@{}lcll@{}}
\toprule
initial belief & reading \(y\) & verdict & covering certificate / witness \\
\midrule
\({(2,2)}\) & \(4\) & \(\circ\) viable & witness: \(u = (1,2)\) cycling \\
\({(1,2),(2,2)}\) & \(3, 4\) & \(\circ\) viable & witness: \(u = (1,2)\) cycling \\
\({(2,1),(2,2)}\) & \(3, 4\) & \(\circ\) viable & witness: \(u = (2,1)\) cycling \\
\({(1,2),(2,1)}\) & \(3\) & \(\bullet\) & common-action (Theorem~\ref{calc-thm:common-action}) \\
\({(1,2),(2,1),(2,2)}\) & \(3, 4\) & \(\bullet\) & common-action (Theorem~\ref{calc-thm:common-action}) \\
\({(1,1)}\) & \(2\) & \(\bullet\) & common-action (Theorem~\ref{calc-thm:common-action}) \\
\bottomrule\end{tabular}}
\caption{The two-patch protection audit: aggregated beliefs under the reading
\(y = z_{1} + z_{2}\), audited against the exact belief-level recursion. Every
nonviable belief is certified by the common-action certificate
(Theorem~\ref{calc-thm:common-action}); the viable beliefs carry explicit blind
witness cycles.}
\label{calc-tab:patch}
\end{table}

The audit is symbolic and one-dimensional, with two two-dimensional instances
(the aggregated reading above and the two-patch protection audit); a full
multidimensional belief-space numerical campaign is deferred to future work,
and the case study makes no claim of a general-purpose computational calculus.

\subsection{10. A Probabilistic Development: Belief-State Safety Values}\label{calc-probabilistic-development}

The certificates of Sections 3--4 rest on set-membership semantics: the information state is a set of compatible states. This section records the stochastic counterpart --- the belief-state safety value of partially observable control --- and its obstruction certificate, connecting the calculus to the partially observable literature (\r{A}str\"om, 1965; Smallwood and Sondik, 1973).

Let \(X, A, D, Y\) be finite, with a stochastic transition \(T(x' \mid x, a)\) and an observation likelihood \(g(y \mid x, a, x')\). Augment the state space with an absorbing unsafe state \(\bot\): for \(x \in \mathcal{V}\), \(T(\bot \mid x, a) = 1 - \sum_{x' \in \mathcal{V}} T(x' \mid x, a)\), and \(T(\bot \mid \bot, a) = 1\). A belief is a distribution \(b\) over \(X \cup \{\bot\}\), updated by Bayes' rule \(\tau(b, a, y)\), with \(\mathbb{P}(y \mid b, a) = \sum_{x, x'} b(x)\, T(x' \mid x, a)\, g(y \mid x, a, x')\). The \emph{safety value} is the maximal survival probability
\[
V_{k}(b) = \max_{\pi}\; \mathbb{P}_{\pi}\bigl( x_{t} \in \mathcal{V} \text{ for } t = 0, \dots, k \;\big|\; b \bigr).
\]

\begin{theorem}[belief-state safety value]\label{calc-thm:pomdp}
\(V_{0}(b) = b(\mathcal{V})\) (the horizon counts the current state), and for \(k \ge 0\)
\[
V_{k+1}(b) = \max_{a \in A}\; \sum_{y \in Y} \mathbb{P}(y \mid b, a)\; V_{k}\bigl(\tau(b, a, y)\bigr),
\]
with the time convention that the policy acts, observes \(y\), then acts again. The recursion is exact on every finite horizon; it is the Smallwood--Sondik belief-state value iteration specialized to the absorbing safety reward (Smallwood and Sondik, 1973).
\end{theorem}

\emph{Proof.} Standard value iteration for the belief-state Markov decision process: conditioning on the observation after the action and applying the optimal continuation gives the identity, and the absorbing state makes \(V_{k}\) the \(k\)-step survival probability (Smallwood and Sondik, 1973). \hfill\(\square\)

\begin{proposition}[chance-constrained common-action obstruction]\label{calc-prop:chance}
Let \(b \in \Delta(\mathcal{V})\) (so \(b(\mathcal{V}) = 1\)) and write \(p(x,a) = 1 - \sum_{x' \in \mathcal{V}} T(x' \mid x, a)\). If
\[
\sum_{x \in \mathcal{V}} b(x)\, p(x,a) \;\ge\; \delta \qquad \text{for every } a \in A,
\]
then \(V_{k}(b) \le 1 - \delta\) for every \(k \ge 1\), and the chance constraint \(\mathbb{P}_{b,\pi}(x_{0}, \dots, x_{k} \in \mathcal{V}) \ge 1 - \varepsilon\) is infeasible under every policy whenever \(\varepsilon < \delta\).
\end{proposition}

\emph{Proof.} The one-step survival probability under action \(a\) is \(1 - \sum_{x} b(x)\, p(x,a) \le 1 - \delta\), so \(V_{1}(b) \le 1 - \delta\). Since \(V_{k+1} \le V_{k}\) --- more steps cannot raise the survival probability --- \(V_{k}(b) \le V_{1}(b) \le 1 - \delta\). \hfill\(\square\)

\begin{proposition}[degenerate (deterministic) limit]\label{calc-prop:degenerate}
Let \(X, A, D, Y\) be finite with the deterministic maps \(x^{+} = F(x,a,d)\) and \(y = O(x,a,x^{+})\) of Section 3.1, and let the stochastic transition and likelihood have supports \(\{F(x,a,d) : d \in D(x)\}\) and \(\{O(x,a,x^{+}) : x^{+} \text{ reachable}\}\). Then the posterior support under the belief-state update is exactly the set-valued post-state \(\mathrm{Post}(B,a,y)\) of Section 3.1, so the support of the belief-state recursion follows the set-valued recursion \(\mathcal{W}_{k}\). For a belief \(b\) carried by \(B \subseteq \mathcal{V}\) and the disturbance read adversarially within its support, \(V_{k}(b) = 1\) if \(B \in \mathcal{W}_{k}\), while if \(B \notin \mathcal{W}_{k}\) then \(V_{k}(b) \le 1 - \min_{x \in B} b(x) < 1\).
\end{proposition}

\emph{Proof.} Under the deterministic kernels the posterior is supported exactly on \(\mathrm{Post}(B,a,y)\), so the support of the belief evolves by the set-valued recursion. If \(B \in \mathcal{W}_{k}\), the witness actions of Theorem~\ref{calc-thm:finite-horizon} keep every compatible branch in \(\mathcal{V}\), giving \(V_{k}(b) = 1\). If \(B \notin \mathcal{W}_{k}\), every action leaves some compatible state \(x \in B\) of mass \(b(x)\) with a successor outside \(\mathcal{V}\) under an admissible disturbance, so the one-step survival probability under every action is at most \(1 - \min_{x \in B} b(x)\); since \(V_{k}\) is non-increasing in \(k\), \(V_{k}(b) \le 1 - \min_{x \in B} b(x) < 1\). \hfill\(\square\)

Proposition~\ref{calc-prop:degenerate} is the sense in which the chance-constrained obstruction degenerates to the common-action obstruction: at the deterministic limit \(V_{1}(b) = 1\) exactly when the common safe-action set is nonempty, and the deficit from \(1\) is at least the smallest compatible mass \(\min_{x \in B} b(x)\). The bound is attained on the worked instances of the Supplementary Material (S1): the two-floor conflict at \(V_{k}(b_{0}) = \tfrac{1}{2}\) for all \(k\), and the delayed hidden regime on the audit grid of Section 9, where the script regenerates all 48 cells of the safety-value table verbatim.

\subsection{11. Conclusion}\label{calc-conclusion}

Under standard finite-dimensional regularity assumptions, viability under perfect measurement has a mature kernel, tangency, and approximation theory. Viability under incomplete observation has
long been asymmetric: the sufficiency side is carried by Veliov's
output-feedback condition and the estimation-tube programme, while the failure side has received less systematic treatment. This paper supplies the missing instrument: the obstruction calculus develops sound, computationally interpretable certificates for nonviability --- that is, sufficient conditions for nonviability, equivalently necessary conditions for viability. They do not exhaust the complement of the epistemic kernel (Section 6.5), and the paper claims no complete characterization. The mechanisms
are organized by one selector principle (Proposition~\ref{calc-prop:selector}) --- observation-based viability
is the existence of a common recursively safe response --- and are
nested in an obstruction ladder (Proposition~\ref{calc-prop:ladder}); in
finite systems backward belief recursion is sound and complete
(Theorem~\ref{calc-thm:finite-horizon}), so the calculus is exact where it
can be and sound elsewhere.

With displayed proofs and explicit witnesses, the certificates identify
five mechanisms by which an observation structure can make robust
viability impossible, plus one through which a single uncorrected policy
empties the kernel: the disturbance can enforce exit in finite time
against every control; a constant observation can merge states with
incompatible admissible controls (Proposition~\ref{calc-prop:emptiness}'s admissibility form);
the compatible states can demand incompatible actions at a single
information state; the information can arrive after the enforced exit;
and the observation fibres can cross the safe-set boundary, defeating
every exact certifier --- with the uncorrected certainty-equivalence controller as the
sixth, policy-specific, mechanism. Each certificate that admits a design
remedy identifies it: a faster indicator, a separating observation, an
observation that separates the sides of each floor instead of aggregate
certification, or the correction of a known bias. For sustainability
governance, where observation is conducted through coarse indicators at
discrete reviews, the calculus states the quantitative limits of what
monitoring can deliver --- and converts each impossibility into a design
specification. Four elements complete the picture: the certificates are complete in the one-step and static-observation classes, with the residual dynamic gap stated as an open problem (Section 7); the selector principle consolidates the certificates as dual witnesses of one correspondence, with the monitoring-adequacy converse and the design rules (Section 8); a calibrated case study and a coverage audit show the certificates reproducing the exact recursion verdict on a grid of delayed hidden-regime systems, with a two-patch audit carrying the aggregation obstruction to two dimensions (Section 9); and the probabilistic belief-state lift connects the calculus to the partially observable control literature, its safety values worked in closed form on the same grid (Section 10). A fourth completion is structural: over polytope-declared
blind-window classes the timing certificate is exactly computable as a linear
program, sound beyond the declaration with an unbounded implementable-class gap
(Section 3.3), and the residual dynamic gap decomposes exactly, its
post-observation recourse phase undecidable by any window-measurable certificate
pair (Section 7).

\part*{Part II --- Probabilistic sufficiency on belief states}
\label{part:probabilistic}
\addcontentsline{toc}{part}{Part II --- Probabilistic sufficiency}

\maketitle

\begin{abstract}
The obstruction calculus characterizes nonviability through certificates on
information states; the probabilistic side asks what the certificates imply
when the information state is a distribution. This paper develops the
belief-state safety value \(V_{k}(b)\) --- the maximal probability of
remaining safe for \(k\) steps --- as a theory carried entirely in exact
rational arithmetic, with six results. First, a support identity: on
finite models with deterministic kernels and deterministic observation
maps the posterior support coincides with the set-valued post-state of the
calculus, the value-one level of the recursion coincides with the
viable-set recursion \(\mathcal{W}^{\mathrm{fb}}_{k}\), and outside it
the deficit
satisfies \(1 - V_{k}(b) \ge \min_{x \in \mathrm{supp}(b)} b(x)\), attained
on natural instances; moreover the alpha-vectors are exactly the
indicators of the maximal jointly survivable subsets of the support ---
an antichain, Sperner-bounded --- which explains why the certified witness
census realizes three types. Second, exactness and parametric closed
forms: over rational input data the finite-horizon values are piecewise
linear with rational alpha-vectors, and on the delayed hidden-regime
class both the declared hold-class and the open-loop values are proved in
closed form for all initial states, all window lengths and horizons, and
all prior weights --- value one exactly on
\(z_{0} \ge 1 + \min(k, T_{\mathrm{obs}})\) and on \(z_{0} \ge 2\)
respectively, the larger prior mass off those regions, zero below the
floor --- so the complete \(16 \times 12\) certified value table, the
36/156 census, and the 384 hold-direction witnesses (348 distinct after
deduplication) become one-line corollaries. Third, the class lattice:
declared hold \(\subseteq\) open-loop \(\subseteq\) sequential
\(\subseteq\) fully observed orders the values, and both gaps are
realized exactly --- the value of flexibility on the timing cells and the
value of observation on the common-action cells, each \(\tfrac{1}{2}\) at
the symmetric prior and zero elsewhere. Fourth, agreement and
stabilization: the segment-form enumeration reproduces the closed forms
on all 48 certified cells and horizons \(k \le 4\), the certificate
partition coincides with the value boundaries, the hold-class value jumps
by exactly \(\tfrac{1}{2}\) at \(z_{0} = 1 + \min(k, T_{\mathrm{obs}})\),
the two classes differ exactly where
\(2 \le z_{0} < 1 + \min(k, T_{\mathrm{obs}})\) (on the certified grid,
exactly the timing cells, 36 combinations), and the value is frozen once
\(k \ge T_{\mathrm{obs}}\) on the finite survivable-set lattice, so the
certified tables are complete. Fifth, a genuinely stochastic layer: an
exact weighted-path deficit identity under noisy observation --- the
deficit is the prior mass weighted by fatal-reading probabilities, and a
two-state sensor instance attains deficit \(1/20\) against a minimal mass
of \(1/2\), so the min-mass bound is not universal there --- a
noisy-probe deadline family with three-valued exact values, and the
repetition trade-off: a second probe cuts the misread probability from
\(1/10\) to \(7/250\), and each probe step costs \(1/10\) of floor
margin. Sixth, a distributionally robust interpolation between the
expected and adversarial operators over contamination ambiguity sets ---
\(V^{\rho}\) is the exact \(\rho\)-mixture of the two, recovering
expectation at \(\rho = 0\) and the worst-case reading at
\(\rho = 1\), with an exact closed form on the delayed class --- an
exact probe-count bound \(p_{\mathrm{wrong}}(n) \le (4\varepsilon(1 -
\varepsilon))^{n/2}\) giving the \(\log(1/\delta)/\mathrm{sep}^{2}\)
probe law, and the general additive deadline law: the threshold is the
floor plus the probe excursion plus drift times deadline, verified on two
instances. Every closed form is verified against brute-force enumeration
from the definitions, including a seeded randomized battery.
\end{abstract}

\textbf{Keywords:} partially observable control; viability kernels;
belief states; safety values; exact arithmetic; adaptive observation;
POMDP

\subsection{Introduction}\label{prob-introduction}

Viability under incomplete observation divides into a necessity side and a
sufficiency side. The obstruction calculus treats the necessity side:
certificates whose firing proves nonviability of an information state
(Abaee, 2026, An obstruction calculus). The probabilistic question is the quantitative counterpart.
If the controller holds a distribution \(b\) over states rather than a set,
how safe is safe? The answer is a function: the belief-state safety value
\(V_{k}(b)\), the maximal worst-case probability of keeping the state
trajectory in the safe set for \(k\) steps. Its value-one level is the
viability statement of the set-valued theory; everything below level one
is measured by the deficit \(1 - V_{k}(b)\), and the obstruction
certificates reappear as lower bounds on that deficit.

This paper develops the probabilistic layer of the viability theory of
Abaee (2026, An obstruction calculus): the calculus there is set-valued, and the present theory
quantifies what remains safe once the information state is a
distribution. It brings together the stochastic selector development, the
belief-state value campaign, and the hidden-parameter learning-deadline
audit into one theory with parametric closed forms, a class lattice, and
an exact stochastic layer; the companion papers (the calculus, its
computational certification, and the worked-systems audit) are unaffected,
and nothing verified in them is re-proved differently here.

The development proceeds in seven steps. Section \ref{prob-model} states the
augmented model --- total transition rows, an absorbing unsafe state
carrying its own observation, masked witnesses --- and the value recursion,
distinguishing the selector's attaining set from its value-one level.
Section \ref{prob-support} connects the belief-state recursion to the
set-valued recursion of the calculus exactly, through the support of the
belief. Section \ref{prob-exact} shows the finite-horizon values are piecewise
linear and exactly representable, and carries the complete certified value
table over the delayed grid. Section \ref{prob-deterministic} evaluates the
deterministic limit, at which the recursion degenerates to a
survived-mass formula over blind action sequences, with the asymmetric
audit carried sequence by sequence. Section \ref{prob-deficit} reads the
obstruction certificates systematically as deficit lower bounds. Section \ref{prob-agreement} proves agreement between the declared-class
segment enumeration and the closed-form values on the certified
delayed-regime class,
quantifies what the class declaration contributes, locates the value's
jump loci, and carries the derivations. Section \ref{prob-adaptive} treats
observation as a design variable: probes, learning deadlines, and the
exact kernel cost of late resolution of a hidden parameter. The running
instances are the delayed hidden-regime instance, with the two-floor
conflict as the second worked case --- both shared with the obstruction
calculus's Supplementary Material. The contributions are: the
belief-state value recursion with its operator laws and the selector's
viable-response structure (Section \ref{prob-model}); the support identity
and exact sufficiency, with the piecewise-linear closed forms on the
delayed hidden-regime class --- hold-class and open-loop values proved
for all initial states, window lengths, horizons, and prior weights ---
and the deficit lower-bound reading of obstruction certificates (Sections
\ref{prob-support}--\ref{prob-deficit}); and the agreement theorem between the
declared-class enumeration and the closed-form values, together with the
adaptive-observation layer and the distributionally robust interpolation
family (Sections \ref{prob-agreement}--\ref{prob-dr}).

\emph{Relation to prior work.} The belief-support reduction at the core
of Theorem \ref{prob-thm:support} is the finite-horizon qualitative analysis
of partially observable Markov decision processes: almost-sure safety
depends only on belief supports, as established for qualitative
objectives by Chatterjee, Doyen, and Henzinger (2009). What is added
here is quantitative and structural: the deficit bound and its witness
reading (Propositions \ref{prob-prop:pl} and \ref{prob-prop:deficit}), the
survivable-set (antichain) formula behind the witness census
(Proposition \ref{prob-prop:antichain}), the class lattice with its realized
gaps (Theorem \ref{prob-thm:lattice}), the parametric closed forms
(Proposition \ref{prob-cor:closed} and Theorem \ref{prob-thm:parametric}), and the
class-scoped timing obstruction. On complexity, finite-horizon partially
observable problems are PSPACE-complete and their no-observation variant
NP-complete (Papadimitriou and Tsitsiklis, 1987); the exponential
witness-growth law of Proposition \ref{prob-prop:pl} is the general behaviour, and the certified two-witness families are structural rather
than representative. Runtime safety filters for learning agents
(Alshiekh et al., 2018) and distributionally robust partially observable
models over ambiguity sets (Nakao, Jiang, and Shen, 2021), whose value
functions remain piecewise linear, bound the present development from
the computational and the modelling side; Section \ref{prob-scope} records
the ambiguity-set interpolation as the formal home of worst-case
wording.

All certified computations are exact: integer and rational arithmetic,
standard library only, deterministic scripts (Section \ref{prob-methods}).

\subsection{The augmented model, the safety value, and the selector}\label{prob-model}

\emph{The augmented state space.} Let the physical safe set be
\(\mathcal{V}\) and let the state space of the model be
\(X^{\perp} = \mathcal{V} \cup \{\bot\}\), where \(\bot\) is an absorbing
unsafe state. For \(x \in \mathcal{V}\) the transition row splits between
\(\mathcal{V}\) and \(\bot\),
\[
\textstyle\sum_{x' \in \mathcal{V}} T(x' \mid x, a) \le 1,\quad
T(\bot \mid x, a) = 1 - \sum_{x' \in \mathcal{V}} T(x' \mid x, a),
\]
and the bottom row is \(T(\bot \mid \bot, a) = 1\): every row of \(T\)
sums to one and \(\bot\) is absorbing, so exit from \(\mathcal{V}\) is
absorption and no trajectory re-enters \(\mathcal{V}\) from outside it.
The observation model is completed on all of \(X^{\perp}\): there is a
distinguished observation \(y_{\bot}\) with
\(g(y_{\bot} \mid \bot, a, \bot) = 1\) (and
\(g(y \mid \cdot, \cdot, \bot) = 0\) for \(y \ne y_{\bot}\)), while
\(g(\cdot \mid x, a, x')\) on \(\mathcal{V}\) is arbitrary non-negative
with the usual normalization per \((x, a)\). A belief is a distribution
\(b\) over \(X^{\perp}\), updated by Bayes' rule \(\tau(b, a, y)\), with
\(\mathbb{P}(y \mid b, a) = \sum_{x, x'} b(x)\, T(x' \mid x, a)\,
g(y \mid x, a, x')\); because every row sums to one and every
\((x, a, x')\) emits exactly one observation with probability one, the
laws \(\mathbb{P}(y \mid b, a)\) sum to one over \(y\) at every belief ---
no sub-probability bookkeeping is needed. On the deterministic kernels of
the certified instances this formalism converts set-valued nonviability into
almost-sure nonviability on the augmented chain.

\begin{definition}[safety value]\label{prob-def:value}
Fix a declared policy class \(\Pi\). The safety value at horizon \(k\) is
the maximal class-survival probability
\[
V^{\Pi}_{k}(b) = \max_{\pi \in \Pi}\; \mathbb{P}_{\pi}\bigl( x_{t} \in
\mathcal{V} \text{ for } t = 0, \dots, k \;\big|\; b \bigr),
\]
the horizon counting the current state. The unrestricted class
\(\Pi = \Pi_{\mathrm{seq}}\) is the class of all sequential policies.
\end{definition}

\emph{The selector and its levels.} In the selector framework this is the
distributional counterpart of the recursive correspondence: the
information state is \(b\), a response is a class-admissible policy
segment. Two objects must be distinguished. The \emph{attaining set} at
\(b\),
\[
\Gamma^{\mathrm{s}}_{k}(b) = \Bigl\{ a \in A : \textstyle\sum_{y}
\mathbb{P}(y \mid b, a)\, V^{\Pi}_{k-1}\bigl(\tau(b, a, y)\bigr)
= V^{\Pi}_{k}(b) \Bigr\},
\]
collects the actions attaining the value; for finite \(A\) it is nonempty
at every belief and horizon --- nonemptiness is automatic and carries no
viability content. The \emph{value-one level},
\[
S^{\Pi}_{k}(b) = \Bigl\{ a \in A : \textstyle\sum_{y} \mathbb{P}(y \mid
b, a)\, V^{\Pi}_{k-1}\bigl(\tau(b, a, y)\bigr) = 1 \Bigr\},
\]
is the selector's viable-response set, and its nonemptiness corresponds
exactly to \(V^{\Pi}_{k}(b) = 1\). Obstruction certificates correspond to
deficits \(1 - V^{\Pi}_{k}(b) > 0\). Notation is fixed as follows:
\(\Gamma^{\Pi}_{k}\) denotes alpha-vector witness sets,
\(\Gamma^{\mathrm{s}}_{k}(b)\) the attaining action set,
\(\mathcal{W}^{\mathrm{fb}}_{k}\) the calculus's feedback viable-set
recursion (a family of state sets), \(\mathcal{W}_{k}\) the blind
recursion (the same, but with one action for the whole set, chosen
before the observation arrives), and
\(\mathcal{W}^{\mathrm{bel}}_{k}\) the belief-space analogue of
\(\mathcal{W}_{k}\) (a family of beliefs); unsuperscripted values do
not appear.

\medskip
\noindent\textbf{Shared symbols across the companion papers.} The
companion papers share this vocabulary with paper-local meanings; the
letters below carry meanings here that differ there:
\begin{center}\small
\begin{tabular}{@{}p{0.42\linewidth}p{0.52\linewidth}@{}}
\toprule
\textbf{Here} & \textbf{In the companions} \\
\midrule
\(\Gamma^{\Pi}_{k}\), \(\Gamma^{\mathrm{s}}_{k}\): the alpha-vector witness sets and the attaining action set & the certificate margin \(\Gamma_{\mathcal{I}}\) (certification); the recourse certificate \(\Gamma_{h}\) (obstruction calculus; minimax duals) \\
\(\rho\): the contamination/interpolation radius of \(\mathcal{D}_{\rho}\) (expected value at \(\rho = 0\), worst case at \(\rho = 1\)) & the mesh relaxation parameter (certification); the net biomass multiplier (applied regime viability) \\
\bottomrule
\end{tabular}
\end{center}


\begin{theorem}[belief-state value recursion]\label{prob-thm:recursion}
\(V^{\Pi}_{0}(b) = b(\mathcal{V})\) (the horizon counts the current
state), and for \(k \ge 0\),
\[
V^{\Pi}_{k+1}(b) = \max_{a \in A(\Pi)}\; \sum_{y \in Y} \mathbb{P}(y \mid
b, a)\; V^{\Pi}_{k}\bigl(\tau(b, a, y)\bigr),
\]
where \(A(\Pi)\) is the class-admissible action set at the current stage
for Markov classes \(\Pi\). The recursion is exact on every finite
horizon; over the unrestricted class it is the Smallwood--Sondik
belief-state value iteration specialized to the absorbing safety reward
(Smallwood and Sondik, 1973).
\end{theorem}

\begin{proof}
The proof is standard value iteration for the belief-state Markov decision
process, with the absorbing state making \(V^{\Pi}_{k}\) the \(k\)-step
class-survival probability (Bertsekas and Shreve, 1978). For non-Markov
classes (the hold and open-loop classes of Section \ref{prob-agreement}) the
class constraint couples stages; the certified computations there use the
equivalent segment form of Proposition \ref{prob-prop:pl} below.
\end{proof}

\subsection{The support identity and exact sufficiency}\label{prob-support}
Sperner, E. (1928). Ein Satz ueber Untermengen einer endlichen Menge.
Mathematische Zeitschrift 27, 544-548.


The bridge to the set-valued theory runs through the support. Write
\(\mathrm{supp}(b)\) for the states carrying positive mass, and let
\(\mathcal{W}^{\mathrm{fb}}_{k}\) denote the calculus's feedback
viable-set recursion:
\(\mathcal{W}^{\mathrm{fb}}_{0} = \{B \subseteq \mathcal{V}\}\), and
\(B \in \mathcal{W}^{\mathrm{fb}}_{k}\) if some admissible \(a\)
admits, for every observation \(y\) that \((B,a)\) can produce, a
\(\mathcal{W}^{\mathrm{fb}}_{k-1}\)-continuation of
\(\mathrm{Post}(B,a,y)\). The blind recursion \(\mathcal{W}_{k}\) is
the special case in which the continuation may not depend on \(y\):
\(\mathcal{W}_{0} = \{B \subseteq \mathcal{V}\}\) and
\(\mathcal{W}_{k} = \{B : \text{some admissible } a \text{ maps every }
x \in B \text{ into } \mathcal{W}_{k-1}\text{-supported posteriors}\}\).
Thus \(\mathcal{W}_{k} \subseteq \mathcal{W}^{\mathrm{fb}}_{k}\),
strictly in general (Remark~\ref{prob-rem:feedback-strict}).

\begin{theorem}[support identity]\label{prob-thm:support}
On finite models with deterministic kernels and deterministic
observation maps, for every action \(a\) and
observation \(y\),
\[
\mathrm{supp}\bigl(b^{+}(\cdot \mid a, y)\bigr)
\;=\; \mathrm{Post}\bigl(\mathrm{supp}(b), a, y\bigr),
\]
the set-valued post-state of the calculus. Consequently, for the
unrestricted sequential class,
\(V_{k}(b) = 1\) if and only if
\(\mathrm{supp}(b) \in \mathcal{W}^{\mathrm{fb}}_{k}\), while if
\(\mathrm{supp}(b) \notin \mathcal{W}^{\mathrm{fb}}_{k}\) then every policy
loses a
compatible state of positive mass along its realized observation path
(the loss can occur at any stage), so
\[
1 - V_{k}(b) \;\ge\; \min_{x \in \mathrm{supp}(b)} b(x).
\]
Moreover \(V_{k}(b)\) is non-increasing in \(k\).\end{theorem}

\begin{proof}
Support identity: a state \(x'\) carries posterior mass exactly when
some \(x \in \mathrm{supp}(b)\) has
\(T(x' \mid x, a)\, g(y \mid x, a, x') > 0\); with deterministic
kernels and observation maps the indicators make this set exactly
\(\mathrm{Post}(\mathrm{supp}(b), a, y)\). Sufficiency is induction on
\(k\): the calculus's witness sequence keeps every compatible trajectory
inside \(\mathcal{V}\) and the posterior supports inside the calculus's
viable sets at every stage, so every branch survives and
\(V_{k}(b) = 1\). For the converse, induct on \(k\); for \(k = 0\)
there is nothing to prove. Let
\(\mathrm{supp}(b) \notin \mathcal{W}^{\mathrm{fb}}_{k}\) and let
\(\pi\) be any policy, with first action \(a\). For each observation
\(y\) that \(a\) produces with positive probability from \(b\), the
posterior \(b^{+}(\cdot \mid a, y)\) has support
\(\mathrm{Post}(\mathrm{supp}(b), a, y)\) by the support identity; if
that support lay in \(\mathcal{W}^{\mathrm{fb}}_{k-1}\) for every such
\(y\), then \(\mathrm{supp}(b)\) would lie in
\(\mathcal{W}^{\mathrm{fb}}_{k}\) by the recursion's definition.
Hence for some \(y\) the posterior support lies outside
\(\mathcal{W}^{\mathrm{fb}}_{k-1}\), and by the induction hypothesis
the continuation of \(\pi\) after \((a,y)\) loses a branch of that
posterior of positive mass. Composing, \(\pi\) loses a branch of
\(b\) of positive mass; every lost branch carries mass at least
\(\min_{x \in \mathrm{supp}(b)} b(x)\), and there are finitely many
branches, giving
\(1 - V_{k}(b) \ge \min_{x \in \mathrm{supp}(b)} b(x)\). Monotonicity:
absorption at \(\bot\) is permanent, so a policy safe for \(k+1\)
stages is safe for \(k\) (equivalently, the survivable families are
prefix-nested).
\end{proof}

\begin{remark}[observation-dependent continuation is strictly
stronger]\label{prob-rem:feedback-strict}
Both recursions are downward closed and
\(\mathcal{W}_{k} \subseteq \mathcal{W}^{\mathrm{fb}}_{k}\); the
inclusion can be strict, even under the hypotheses of
Theorem~\ref{prob-thm:support}. Let \(p \neq q\) be safe states and
\(a_{0}\) the only admissible action that keeps both alive, sending
them to distinct successors \(p^{+} \neq q^{+}\) that the observation
map distinguishes; let the unique safe continuation at \(p^{+}\) be
\(b\) and the unique safe continuation at \(q^{+}\) be
\(c \neq b\). Then
\(\{p,q\} \in \mathcal{W}^{\mathrm{fb}}_{2}\) --- play \(a_{0}\),
read the observation, continue with \(b\) or \(c\) accordingly ---
while \(\{p,q\} \notin \mathcal{W}_{2}\): the first action must be
\(a_{0}\), and no single second action is safe at both \(p^{+}\) and
\(q^{+}\).
For a belief with support \(\{p,q\}\) it follows that
\(V^{\Pi_{\mathrm{seq}}}_{2}(b) = 1\) while
\(V^{\Pi_{B}}_{2}(b) < 1\): the value of observation is strict here in
the strongest form available, all or nothing. The two recursions agree
when no information arrives --- on a blind window,
\(\mathcal{W}^{\mathrm{fb}}_{k} = \mathcal{W}_{k}\), the set-level
content of Theorem~\ref{prob-thm:lattice}'s
\(V^{\mathrm{ol}} = V^{\mathrm{seq,blind}}\). Theorems~\ref{prob-thm:support}
and~\ref{prob-prop:degen} therefore differ in their class, not only in their
normalization: the first is a feedback theorem, the second a blind one,
and Proposition~\ref{prob-cor:closed} is accordingly not a corollary of
Theorem~\ref{prob-thm:support}.
\end{remark}

\begin{proposition}[survivable sets, witnesses, and the census]\label{prob-prop:antichain}
Let the kernels and observation maps be deterministic, fix a declared
class, and let \(\mathcal{S}_{k}\) be the family of subsets
\(S\) of the support that are jointly survivable by one admissible
class-element (sequence or policy). Then:
(i) \(V_{k}(b) = \max_{S \in \mathcal{S}_{k}} b(S)\);
(ii) the alpha-vectors are exactly the indicators of the maximal elements
of \(\mathcal{S}_{k}\) in the componentwise order;
(iii) the maximal elements form an antichain, so their number is at
most \(\binom{n}{\lfloor n/2 \rfloor}\) for a support of size
\(n\) (Sperner);
(iv) the deficit is \(1 - V_{k}(b) = \min \, b(S^{c})\), the minimum
over the maximal survivable sets, and the min-mass bound is the special
case \(|S^{c}| = 1\).
On the delayed instance the maximal survivable sets are
\(\{+, -\}\) exactly when \(z_{0} \ge 2\) and the two singletons
otherwise (at \(z_{0} \ge 1\)), giving the empty, mirror-pair, and
merged witness families --- the three-type census of the certified campaign
is this antichain read off the grid.
\end{proposition}

Branches survive a fixed sequence independently, survived prior mass sums
over surviving branches, and maximizing over class-elements gives (i);
the indicator of a survivable set is a witness and dominance gives (ii)
with deduplication; maximality gives (iii) with Sperner's theorem; and
monotonicity of \(b(S)\) in \(S\) gives (iv).

\begin{proof}
(i) Fix a class element; the compatible branches it keeps alive form a
set \(S \in \mathcal{S}_{k}\) by definition and its value at \(b\) is
\(b(S)\); conversely every \(S \in \mathcal{S}_{k}\) is realized by
some element, so maximizing over elements gives
\(V_{k}(b) = \max_{S} b(S)\). (ii) The element's alpha-vector is its
survival indicator \(\mathbf{1}_{S}\); the indicator of a non-maximal
\(S\) is pointwise dominated by that of a maximal \(S' \supseteq S\),
so the attained maximum is always attained on a maximal element and the
undominated vectors --- the alpha-set --- are exactly the indicators of
the maximal elements. (iii) Maximality is inclusion-maximality, so the
maximal sets form an antichain, and their count is at most
\(\binom{n}{\lfloor n/2 \rfloor}\) by Sperner's theorem (Sperner,
1928). (iv) \(1 - \max_{S} b(S) = \min_{S} b(S^{\complement})\) over
the same maximal sets; when a maximal set misses exactly one state this
is the min-mass bound. The delayed-instance classification is the
computed one (the verification script).
\end{proof}

\begin{proposition}[closed form on the delayed class, declared hold class]
\label{prob-cor:closed}
On the delayed hidden-regime instance --- states \((z, \theta)\) with
\(\theta \in \{-1, +1\}\), prior \(b_{0} = (\tfrac{1}{2}, \tfrac{1}{2})\)
from \(z_{0}\), no observation before \(T_{\mathrm{obs}}\), exact
revelation at \(T_{\mathrm{obs}}\), and unit drift against the floors ---
the belief support follows the two branches \(z_{0} \mp k\), and the
\emph{declared hold class} value (\(\Pi_{\mathrm{hold}}\): one blind
hold \(u \equiv \pm 1\) for the whole window) is
\[
\small
V^{\Pi_{\mathrm{hold}}}_{k}(b_{0}) \;=\;
\begin{cases}
\tfrac{1}{2}\,\bigl(\mathbb{1}[z_{0} - k \ge 1] + \mathbb{1}[z_{0} + k \ge 1]\bigr), & k \le T_{\mathrm{obs}},\\[1ex]
\tfrac{1}{2}\,\bigl(1 + \mathbb{1}[z_{0} \ge 1 + T_{\mathrm{obs}}]\bigr), & k > T_{\mathrm{obs}}.
\end{cases}
\]
In particular \(V^{\Pi_{\mathrm{hold}}}_{k}(b_{0}) = 1\) exactly on
\(z_{0} \ge 1 + T_{\mathrm{obs}}\), recovering the declared-class
deterministic verdict as the value-one level; Theorem \ref{prob-thm:class}
quantifies where the unrestricted class does better. This is a statement
about the declared class, proved by direct enumeration below --- not a
corollary of Theorem \ref{prob-thm:support}, whose value is the unrestricted
one.
\end{proposition}
\begin{theorem}[parametric closed forms]\label{prob-thm:parametric}
On the delayed hidden-regime instance with prior mass \(w\) on the
\(\theta = +1\) branch (arbitrary rational \(w \in [0,1]\)), initial
state \(z_{0}\) (arbitrary real), and effective window
\(\ell = \min(k, T_{\mathrm{obs}}) \ge 1\):
\begin{align*}
V^{\Pi_{\mathrm{hold}}}_{k}(b_{0}) &=
\begin{cases} 0, & z_{0} < 1,\\
1, & z_{0} \ge 1 + \ell,\\
\max(w,\, 1 - w), & 1 \le z_{0} < 1 + \ell,
\end{cases}
\\[0.4ex]
V^{\Pi_{\mathrm{ol}}}_{k}(b_{0}) &=
\begin{cases} 0, & z_{0} < 1,\\
1, & z_{0} \ge 2,\\
\max(w,\, 1 - w), & 1 \le z_{0} < 2.
\end{cases}
\end{align*}
For \(k = 0\), \(V_{0}(b_{0}) = \mathbb{1}[z_{0} \ge 1]\). At the
symmetric prior both off-viability values are \(\tfrac12\), and
Proposition \ref{prob-cor:closed}, the complete certified table, the
36/156 census, and the witness types are one-line corollaries (the grid
has \(z_{0} \ge 1\), so the zero level is vacuous there).
\end{theorem}

\begin{proof}
Survival requires the state in the safe set at every stage, in particular
at \(t = 0\): if \(z_{0} < 1\) both branches are below the floor
immediately and the value is zero. Under a hold \(u\) both branches
move monotonically, so each survives the window exactly if it is at or
above the floor at its horizon position: under \(u = +1\) the
\(+\) branch rises (safe) and the \(-\) branch ends at
\(z_{0} - \ell\), giving survived mass
\(w + (1 - w)\,\mathbb{1}[z_{0} \ge 1 + \ell]\); the mirror holds
for \(u = -1\); maximizing over the two holds gives the first display.
Under open-loop play the partial sums of any \(\pm1\)-sequence move the
branches in opposition, so the first move sends one branch to
\(z_{0} - 1\) and joint survival requires \(z_{0} \ge 2\); the
alternating sequence keeps both branches in
\([z_{0} - 1, z_{0} + 1]\), which lies above the floor exactly when
\(z_{0} \ge 2\); a single branch is survivable by the constant
sequence away from it (which makes that branch rise), so off the joint
region the value is the larger prior mass.
\end{proof}

\begin{proof}[Derivation of Proposition \ref{prob-cor:closed}]
While \(k \le T_{\mathrm{obs}}\) no observation arrives, so a
declared-class policy commits to one blind hold \(u \equiv \pm1\): under \(u \equiv +1\) the
\(\theta = +1\) branch sits at \(z_{0} + k\) at the horizon and the
\(\theta = -1\) branch at \(z_{0} - k\); under \(u \equiv -1\) the roles
swap. Each branch carries mass \(\tfrac{1}{2}\) and survives exactly when
it stays at or above the floor, giving
\(\tfrac{1}{2}(\mathbb{1}[z_{0} - k \ge 1] + \mathbb{1}[z_{0} + k \ge 1])\),
and the expectation over the prior retains exactly the surviving
mass --- the lost branch simply carries the missing half; no adversarial
selection is involved.
For \(k > T_{\mathrm{obs}}\), revelation at \(T_{\mathrm{obs}}\) splits
the belief: the posterior collapses to the realized branch, and the
rising control \(u = +1\) keeps it safe exactly when it sits at or
above the floor at revelation --- so survival reduces to surviving the
blind window of length \(T_{\mathrm{obs}}\) on both branches, the
mass-\(\tfrac12\) falling branch sitting at \(z_{0} - T_{\mathrm{obs}}\)
at revelation, whence the indicator
\(\mathbb{1}[z_{0} \ge 1 + T_{\mathrm{obs}}]\) weighting the
certainly-saved mass \(\tfrac{1}{2}\). (The rising continuation is the
point: a falling continuation would exit later, and Proposition
\ref{prob-prop:freeze} pins the frozen value for all horizons.)
\end{proof}

\subsection{Exact piecewise-linear values}\label{prob-exact}

\begin{proposition}[rational alpha-vectors, class-restricted]\label{prob-prop:pl}
Fix a declared class \(\Pi\). At every horizon \(k\),
\(V^{\Pi}_{k}\) is piecewise linear and convex:
\(V^{\Pi}_{k}(b) = \max_{\alpha \in \Gamma^{\Pi}_{k}} \alpha^{\top} b\)
for a finite witness set \(\Gamma^{\Pi}_{k}\) on
\(\mathbb{Q}^{|\mathcal{V}|+1}\) (coordinates indexed by
\(\mathcal{V} \cup \{\bot\}\)), constructed by the masked backup
\[
\alpha^{a,\,\gamma_{\cdot}}(x) =
\sum_{x'} T(x' \mid x, a) \sum_{y} g(y \mid x, a, x')\, \gamma_{y}(x'),\quad
\gamma_{y} \in \Gamma^{\Pi}_{k},
\]
with \(\Gamma^{\Pi}_{0} = \{ \mathbf{1}_{\mathcal{V}} \}\) and the
boundary condition \(\alpha(\bot) = 0\) applied at every backup (the
backup acts on \(\mathcal{V}\) only; the \(\bot\)-coordinate of every
witness is anchored at zero and remains zero under the backup) and
\(\Gamma^{\Pi}_{k+1} = \{ \alpha^{a, \gamma_{\cdot}} : a \in A(\Pi),\
\gamma_{\cdot} \text{ a selection} \}\), duplicates removed by exact
equality. For a blind open-loop window the equivalent segment form
applies: each admissible action segment \((u_{0}, \dots, u_{m}) \in \Pi\)
contributes the witness
\(\alpha^{\mathrm{seg}} = (\mathbf{1}[\text{the segment saves branch } x])_{x}\)
and \(V^{\Pi}_{k}(b) = \max_{\mathrm{seg}} \alpha^{\mathrm{seg}\top} b\).
Over rational \(T\), \(g\), and priors, every \(\alpha\) is an exact
rational vector, every \(V^{\Pi}_{k}(b)\) at a rational belief is an exact
rational, and the attaining witness is part of the computation. The
worst-case growth is
\(|\Gamma^{\Pi}_{k+1}| \le |A|\,|\Gamma^{\Pi}_{k}|^{|Y|}\);
deduplication and pruned selections are the practical controls, and the
exponential law is the exact complexity statement for the general class.
\end{proposition}

The PL structure is Smallwood--Sondik's; the masking factor anchors
\(\alpha(\bot) = 0\) (survival requires the current state in
\(\mathcal{V}\), and \(\bot\) is absorbing, so no backup can revive a lost
branch). The exact-arithmetic discipline is the point: the backup maps
rational vectors to rational vectors by finite sums and products,
selections are finite, and equality tests for deduplication are exact on
rationals. In the verification campaign the delayed grid enumerates the two
hold-direction witnesses at each of its 192 cell-horizon combinations ---
384 exact rational vectors (the two-floor instance's witness set is
carried separately) --- and 348 distinct vectors remain after
exact-equality deduplication.

\begin{proof}
Induction on \(k\). At \(k = 0\),
\(V^{\Pi}_{0}(b) = b(\mathcal{V}) = \mathbf{1}_{\mathcal{V}}^{\top} b\)
with \(\alpha(\bot) = 0\). For the step, substitute the induction
hypothesis into the recursion of Theorem~\ref{prob-thm:recursion}: the finite
maximum over \(a \in A(\Pi)\) and witness tuples
\((\gamma_{y})_{y}\), \(\gamma_{y} \in \Gamma^{\Pi}_{k}\), commutes
with the observation sum, and collecting the coefficient of \(b(x)\)
gives exactly the masked backup \(\alpha^{a, \gamma_{\cdot}}\); the
\(\bot\)-coordinate stays zero because the backup sums over
\(\mathcal{V}\) only and the anchor \(\alpha(\bot) = 0\) is applied
at every stage. Finiteness: finitely many action choices and witness
tuples. Rationality: rational transition and observation data with
rational \(\Gamma^{\Pi}_{k}\) give rational backups, by induction.
Piecewise linearity and convexity are the form "finite maximum of linear
functionals"; for non-Markov classes the same construction runs with
history-indexed witnesses (the segment form), the class constraint
pinning the available action at each stage, which restricts the backup
tuples without changing the argument.
\end{proof}

\begin{remark}[dimension]\label{prob-rem:dim}
On the two-floor instance the value function on the one-dimensional
simplex is \(\max(b_{0}, b_{1})\) for all certified horizons (Figure
\ref{prob-fig:twofloor}); the witness alpha-vectors live in the two-coordinate
display, with the suppressed bottom coordinate identically zero --- padded
to \((1,0,0), (0,1,0)\) on \(\mathcal{V} \cup \{\bot\}\).
\end{remark}

On the delayed grid the campaign evaluates
\(\alpha_{k}(z_{0}, T_{\mathrm{obs}})\) at all
\(16 \times 3 \times 4 = 192\) cell-horizon combinations of
\(z_{0} \in \{1.0, 1.1, \dots, 2.5\}\),
\(T_{\mathrm{obs}} \in \{1, 2, 3\}\), \(k \in \{1, 2, 3, 4\}\). The
realized vectors take exactly three types: \((1, 1)\) (both branches
savable in the blind window), \(72\) occurrences; and the mirror pair
\((1, 0)\) and \((0, 1)\) (one branch savable), \(156\) occurrences each
--- the symmetry between the two mirror types is exact,
\(72 + 156 + 156 = 384\). The type counts are the pre-deduplication
census over the two hold directions: on the \(36\) combinations where
both holds save both branches the two witnesses coincide at \((1, 1)\)
(\(2 \times 36 = 72\)), and exact-equality deduplication merges them,
leaving \(36 + 156 + 156 = 348\) distinct vectors --- one per
unit-value combination, two per the others. Table \ref{prob-tab:campaign}
carries the complete value table; Figure \ref{prob-fig:staircase} displays
the closed forms.

\begin{table*}[t]
\centering
\footnotesize
\caption{The complete certified value table on the delayed grid: the exact
\(V_{k}(b_{0})\) (values \(1\) or \(\tfrac{1}{2}\)) for all 16 grid
states, three blind-window lengths, four horizons --- 192 cell-horizon
combinations, each regenerated by script and each equal to both the
declared-class segment enumeration and the closed form of Proposition
\ref{prob-cor:closed}.}
\label{prob-tab:campaign}
\begin{tabular}{@{}l cccc c@{\hspace{9pt}}cccc c@{\hspace{9pt}}cccc@{}}
\toprule
 & \multicolumn{4}{c}{$T_{\mathrm{obs}} = 1$}
 & \multicolumn{4}{c}{$T_{\mathrm{obs}} = 2$}
 & \multicolumn{4}{c}{$T_{\mathrm{obs}} = 3$}\\
\cmidrule(lr){2-5}\cmidrule(lr){6-9}\cmidrule(lr){10-13}
$z_0$ & $k{=}1$ & $k{=}2$ & $k{=}3$ & $k{=}4$
 & $k{=}1$ & $k{=}2$ & $k{=}3$ & $k{=}4$
 & $k{=}1$ & $k{=}2$ & $k{=}3$ & $k{=}4$\\
\midrule
1.0 & $\tfrac12$ & $\tfrac12$ & $\tfrac12$ & $\tfrac12$ & $\tfrac12$ & $\tfrac12$ & $\tfrac12$ & $\tfrac12$ & $\tfrac12$ & $\tfrac12$ & $\tfrac12$ & $\tfrac12$\\
1.1 & $\tfrac12$ & $\tfrac12$ & $\tfrac12$ & $\tfrac12$ & $\tfrac12$ & $\tfrac12$ & $\tfrac12$ & $\tfrac12$ & $\tfrac12$ & $\tfrac12$ & $\tfrac12$ & $\tfrac12$\\
1.2 & $\tfrac12$ & $\tfrac12$ & $\tfrac12$ & $\tfrac12$ & $\tfrac12$ & $\tfrac12$ & $\tfrac12$ & $\tfrac12$ & $\tfrac12$ & $\tfrac12$ & $\tfrac12$ & $\tfrac12$\\
1.3 & $\tfrac12$ & $\tfrac12$ & $\tfrac12$ & $\tfrac12$ & $\tfrac12$ & $\tfrac12$ & $\tfrac12$ & $\tfrac12$ & $\tfrac12$ & $\tfrac12$ & $\tfrac12$ & $\tfrac12$\\
1.4 & $\tfrac12$ & $\tfrac12$ & $\tfrac12$ & $\tfrac12$ & $\tfrac12$ & $\tfrac12$ & $\tfrac12$ & $\tfrac12$ & $\tfrac12$ & $\tfrac12$ & $\tfrac12$ & $\tfrac12$\\
1.5 & $\tfrac12$ & $\tfrac12$ & $\tfrac12$ & $\tfrac12$ & $\tfrac12$ & $\tfrac12$ & $\tfrac12$ & $\tfrac12$ & $\tfrac12$ & $\tfrac12$ & $\tfrac12$ & $\tfrac12$\\
1.6 & $\tfrac12$ & $\tfrac12$ & $\tfrac12$ & $\tfrac12$ & $\tfrac12$ & $\tfrac12$ & $\tfrac12$ & $\tfrac12$ & $\tfrac12$ & $\tfrac12$ & $\tfrac12$ & $\tfrac12$\\
1.7 & $\tfrac12$ & $\tfrac12$ & $\tfrac12$ & $\tfrac12$ & $\tfrac12$ & $\tfrac12$ & $\tfrac12$ & $\tfrac12$ & $\tfrac12$ & $\tfrac12$ & $\tfrac12$ & $\tfrac12$\\
1.8 & $\tfrac12$ & $\tfrac12$ & $\tfrac12$ & $\tfrac12$ & $\tfrac12$ & $\tfrac12$ & $\tfrac12$ & $\tfrac12$ & $\tfrac12$ & $\tfrac12$ & $\tfrac12$ & $\tfrac12$\\
1.9 & $\tfrac12$ & $\tfrac12$ & $\tfrac12$ & $\tfrac12$ & $\tfrac12$ & $\tfrac12$ & $\tfrac12$ & $\tfrac12$ & $\tfrac12$ & $\tfrac12$ & $\tfrac12$ & $\tfrac12$\\
2.0 & $1$ & $1$ & $1$ & $1$ & $1$ & $\tfrac12$ & $\tfrac12$ & $\tfrac12$ & $1$ & $\tfrac12$ & $\tfrac12$ & $\tfrac12$\\
2.1 & $1$ & $1$ & $1$ & $1$ & $1$ & $\tfrac12$ & $\tfrac12$ & $\tfrac12$ & $1$ & $\tfrac12$ & $\tfrac12$ & $\tfrac12$\\
2.2 & $1$ & $1$ & $1$ & $1$ & $1$ & $\tfrac12$ & $\tfrac12$ & $\tfrac12$ & $1$ & $\tfrac12$ & $\tfrac12$ & $\tfrac12$\\
2.3 & $1$ & $1$ & $1$ & $1$ & $1$ & $\tfrac12$ & $\tfrac12$ & $\tfrac12$ & $1$ & $\tfrac12$ & $\tfrac12$ & $\tfrac12$\\
2.4 & $1$ & $1$ & $1$ & $1$ & $1$ & $\tfrac12$ & $\tfrac12$ & $\tfrac12$ & $1$ & $\tfrac12$ & $\tfrac12$ & $\tfrac12$\\
2.5 & $1$ & $1$ & $1$ & $1$ & $1$ & $\tfrac12$ & $\tfrac12$ & $\tfrac12$ & $1$ & $\tfrac12$ & $\tfrac12$ & $\tfrac12$\\
\bottomrule
\end{tabular}
\end{table*}

\begin{figure}[t]
\centering
\includegraphics[width=\linewidth]{figs_bs2/fig_staircase.pdf}
\caption{The closed-form values of Proposition \ref{prob-cor:closed} on the
delayed grid: flat at \(\tfrac{1}{2}\) while a branch falls, jumping by
exactly \(\tfrac{1}{2}\) at \(z_{0} = 1 + \min(k, T_{\mathrm{obs}})\),
constant beyond.}
\label{prob-fig:staircase}
\end{figure}

\subsection{The deterministic limit and the two operators}\label{prob-deterministic}

\begin{remark}[the two operators]\label{prob-rem:operators}
The expectation over the disturbance law and the adversarial reading of
its support are \emph{different Bellman operators}. Proposition
\ref{prob-prop:degen} states the deterministic degeneration as a change of
operator: when the kernels are deterministic the stochastic expectation
is replaced by the robust (adversarial-support) reading, and the
resulting value is the robust value of the companion papers. The
degeneration is a theorem about the robust operator; it is not a claim
about the stochastic plant, whose expectation and robust readings differ
in general (Section \ref{prob-scope}).
\end{remark}

\begin{proposition}[survived-mass formula]\label{prob-prop:degen}
Let the kernels be deterministic, \(x^{+} = F(x, a, d)\), and let the
disturbance be read adversarially within its support --- the robust
operator of Remark \ref{prob-rem:operators}. Let the window be blind with
declared sequential class \(\Pi_{B}\). Then the degenerated value is the
\emph{survived-mass formula}
\[
V_{k}(b) = \max_{(u_{0},\dots) \in \Pi_{B}}\;
\sum_{x \in \mathrm{supp}\, b} b(x)\,
\mathbf{1}\bigl[ x \text{ survives } (u_{0}, \dots) \bigr],
\]
a maximum of rational linear functions of \(b\) (the surviving indicators
are the degenerate alpha-vectors). Consequently: if the support \(B\)
carries a jointly surviving blind sequence, \(V_{k}(b) = 1\); and if every
declared sequence loses some branch, then
\(V_{k}(b) \le 1 - \min_{x \in B} b(x)\) --- the robust common-action
verdicts and the minimal-mass deficit bound of the companion papers,
recovered as exact values of the robust operator.
\end{proposition}

With deterministic kernels the posterior support follows the set-valued
recursion, survival from a branch under a fixed sequence is a
deterministic event, and prior mass sums over surviving branches;
maximizing over declared sequences gives the formula, from which both
consequences are immediate (a lost branch carries at least the minimal
mass).

\begin{proof}
Under the robust operator a branch survives a declared sequence exactly
when it survives under every disturbance in its support, so survival is
deterministic per branch and the value is the maximum over declared
sequences of the total survived mass --- the displayed formula, whose
summands are the survived indicators, the degenerate alpha-vectors.
Each indicator is a \(\{0,1\}\)-valued linear function of \(b\), so
the value is a maximum of finitely many rational linear functions. If
some declared sequence keeps every branch of \(B\), its indicator is
identically one on \(B\) and \(V_{k}(b) = 1\). If every declared
sequence loses some branch, fix a maximizing sequence and a lost branch
\(x^{*}\): the survived mass is \(1 - b(x^{*}) \le 1 -
\min_{x \in B} b(x)\) --- the robust common-action verdict with the
minimal-mass deficit bound.
\end{proof}

\begin{proposition}[stabilization on the survivable-set lattice]\label{prob-prop:freeze}
With deterministic kernels the survivable-set families satisfy
\(\mathcal{S}_{\ell+1} \subseteq \mathcal{S}_{\ell}\); hence
\(V_{\ell}(b)\) is non-increasing in the horizon and stabilizes after
at most \(2^{|\mathrm{supp}\, b|}\) strict decreases, the frozen value
being \(\max_{S}\, b(S)\) over the maximal sets of the frozen family
--- an exact computable infinite-horizon value. On the delayed instance
the open-loop family is frozen from \(\ell = 1\) (oscillation and
single-branch rises are length-independent), and the declared class
freezes once \(k \ge T_{\mathrm{obs}}\) (the window ends and the
rising continuation \(u = +1\) then preserves safety indefinitely
above the floor); consequently every certified table at
\(T_{\mathrm{obs}} \le 3\), \(k \le 4\) already records the
full-horizon value.
\end{proposition}

Survivability over a longer window is the conjunction of survivability
over the shorter, giving the decreasing lattice; a monotone sequence on a
finite lattice freezes; the instance's length-independence and the
window's end give the two freezing statements, verified to horizon
\(40\) by the verification script.

On the two-floor chance instance --- states \(\{0, 1\}\), actions
\(u \in \{\tfrac{2}{5}, \tfrac{3}{5}\}\), state \(0\) surviving exactly
the lower action and state \(1\) exactly the higher --- the one-step
witness set is \(\Gamma^{\Pi}_{1} = \{(1,0), (0,1)\}\) and
\(V_{k}(b_{0}) = \tfrac{1}{2}\) for all certified \(k \le 5\): the chance
deficit \(\delta = \tfrac{1}{2}\) is attained, and the deficit equals the
minimal mass (Figure \ref{prob-fig:twofloor}).

The asymmetric audit prices the deficit beyond the symmetric case. At
\((z_{0}, T_{\mathrm{obs}}) = (\tfrac{3}{2}, 2)\) with prior masses
\(w = \tfrac{3}{10}\) and \(\tfrac{7}{10}\), the four blind two-step
sequences survive
\(\{ \tfrac{3}{10}, \tfrac{3}{10}, \tfrac{7}{10}, \tfrac{7}{10} \}\)
(Table \ref{prob-tab:masses}): every sequence loses a branch, so
\(V_{2}(b_{0}) = \tfrac{7}{10}\) and the deficit \(\tfrac{3}{10}\) equals
the minimal mass --- the degenerate bound of Proposition
\ref{prob-prop:deficit} attained off the symmetric prior (Figure
\ref{prob-fig:masses}).

\begin{table}[t]
\centering
\footnotesize
\caption{The asymmetric audit at
\((z_{0}, T_{\mathrm{obs}}) = (\tfrac{3}{2}, 2)\),
\(b_{0} = (\tfrac{3}{10}, \tfrac{7}{10})\): survived mass of each blind
two-step sequence. Every sequence loses a branch at step \(1\) --- at
\(z_{0} = \tfrac{3}{2}\) the first move sends one branch to
\(\tfrac{1}{2} < 1\) ---; the value is
\(\tfrac{7}{10}\) and the deficit \(\tfrac{3}{10}\) equals the minimal
mass.}
\label{prob-tab:masses}
\begin{tabular}{@{}l c l@{}}
\toprule
Sequence & Survived mass & Lost branch\\
\midrule
\((+1, +1)\) & $\tfrac{3}{10}$ & heavier branch, step 1\\
\((+1, -1)\) & $\tfrac{3}{10}$ & heavier branch, step 1\\
\((-1, +1)\) & $\tfrac{7}{10}$ & lighter branch, step 1\\
\((-1, -1)\) & $\tfrac{7}{10}$ & lighter branch, step 1\\
\bottomrule
\end{tabular}
\end{table}

\begin{figure}[t]
\centering
\includegraphics[width=\linewidth]{figs_bs2/fig_masses.pdf}
\caption{The asymmetric audit: survived mass by blind two-step sequence
at \((\tfrac{3}{2}, 2)\) with \(w = \tfrac{3}{10}\). The best sequences
save the heavy branch only; the deficit \(\tfrac{3}{10}\) is the minimal
mass, attained.}
\label{prob-fig:masses}
\end{figure}

\begin{proof}
A class element surviving \(\ell + 1\) stages survives its prefix of
\(\ell\) stages, and declared classes are closed under prefixes, so
\(\mathcal{S}_{\ell+1} \subseteq \mathcal{S}_{\ell}\); by
Proposition~\ref{prob-prop:antichain}(i), \(V_{\ell}(b)\) is non-increasing
in \(\ell\). The families live in the finite lattice
\(2^{\mathrm{supp}\, b}\), so each strict decrease retires at least
one subset and at most \(2^{|\mathrm{supp}\, b|}\) strict decreases
occur; the frozen value is the maximal-set mass of the frozen family. On
the delayed instance the closed forms of
Theorem~\ref{prob-thm:parametric} are length-independent in the stated
regimes (oscillation and single-branch rises do not change with added
stages), so the open-loop family is frozen from \(\ell = 1\); the
declared class freezes once \(k \ge T_{\mathrm{obs}}\) because the
window has ended and the rising continuation keeps every revealed branch
above the floor indefinitely; the verification script certifies the frozen
values on the certified tables.
\end{proof}

\subsection{Certificates as deficit bounds}\label{prob-deficit}

The obstruction certificates bound the deficit from below; their attaining
witnesses are the deficit's dual objects --- computed lower-bound sides of
the exact value, not surrogates for it.

\begin{proposition}[deficit lower bounds]\label{prob-prop:deficit}
Write \(\Delta_{k}(b) = 1 - V_{k}(b)\) for the deficit, and let
\(\mathcal{W}^{\mathrm{bel}}_{k}\)
denote the \(k\)-step robust kernel in belief space: the family of beliefs
whose support admits a jointly surviving declared-class sequence. Each
companion certificate is a computable lower bound on the deficit.
(i) \emph{Chance-constrained common action:} if the one-step violation
probabilities satisfy \(\sum_{x} b(x)\, p(x, a) \ge \delta\) for every
\(a \in A\), then \(\Delta_{k}(b) \ge \delta\) for every \(k \ge 1\).
(ii) \emph{Degenerate minimal mass:} at the deterministic limit, if
\(b \notin \mathcal{W}^{\mathrm{bel}}_{k}\) (the support admits no jointly surviving declared
sequence), then \(\Delta_{k}(b) \ge \min_{x \in \mathrm{supp}\, b} b(x)\)
--- the two-floor instance attains \(\delta = \tfrac{1}{2}\), and on the
delayed grid at \((z_{0}, T_{\mathrm{obs}}) = (\tfrac{3}{2}, 2)\) every
blind sequence loses a branch, so the deficit is at least the minimal
mass \(\tfrac{3}{10}\) of the certified asymmetric prior, attained
(Table \ref{prob-tab:masses}).
(iii) \emph{Monotone deficit:} \(V_{k+1}(b) \le V_{k}(b)\), so the
deficit is non-decreasing in the horizon; on the certified nonviability
region of the delayed grid the declared-class deficit is constant in
\(k\) once \(k > T_{\mathrm{obs}}\) (and at every \(k \ge 1\) on
the common-action cells), while on the timing cells it steps from
\(0\) at \(k = 1\) to \(\tfrac{1}{2}\) at \(k \ge 2\) (Theorem
\ref{prob-thm:agree}d).
For finite horizons the exact deficit
\(1 - \max_{\alpha \in \Gamma_{k}} \alpha^{\top} b\) is computable, and
the bounds are witnesses of it, not surrogates.
\end{proposition}

\begin{proof}
(i) At \(k = 1\), survival under \(a\) has probability
\(1 - \sum_{x} b(x)\, p(x, a) \le 1 - \delta\), and the hypothesis
makes the bound uniform over \(a\), so \(V_{1}(b) \le 1 - \delta\);
monotonicity of \(V_{k}\) in \(k\) (Theorem~\ref{prob-thm:support}) gives
\(\Delta_{k}(b) \ge \Delta_{1}(b) \ge \delta\) for every
\(k \ge 1\). (ii) At the deterministic limit
Proposition~\ref{prob-prop:degen} gives \(V_{k}(b) = \max_{S} b(S)\) over
the jointly surviving sets; if
\(b \notin \mathcal{W}^{\mathrm{bel}}_{k}\) no full-mass set
survives, so every surviving \(S\) misses a state of mass at least
\(\min_{x} b(x)\) and \(\Delta_{k}(b) \ge \min_{x} b(x)\). The
two-floor attainment \(\delta = \tfrac12\) and the delayed-instance
attainments are the computed records of the verification script.
\end{proof}

\subsection{Agreement on the delayed-regime class}\label{prob-agreement}

The delayed hidden-regime instance of the calculus paper --- regimes
\(\theta \in \{-1, +1\}\), blind window \(T_{\mathrm{obs}}\), declared
hold class \(\{u \equiv \pm 1\}\), prior \((\tfrac12, \tfrac12)\) from
\(z_{0}\) --- is the agreement test bed. The class declaration is part of
the model: the blind-window policy may be required to hold one action
throughout the window (the declared hold class), or may act sequentially
without observation, or may act with open-loop time-variation.

\begin{theorem}[certificate--value agreement]\label{prob-thm:agree}
On the certified grid (\(z_{0} \in \{1.0, \dots, 2.5\}\),
\(T_{\mathrm{obs}} \in \{1,2,3\}\); 48 cells), over the declared hold
class, and horizons \(k \le 4\):
(a) the segment-form enumeration of the declared hold class
(Proposition \ref{prob-prop:pl}, segment form) returns exactly the closed
form of Proposition
\ref{prob-cor:closed} (on the certified grid the second indicator of the
\(k \le T_{\mathrm{obs}}\) line is vacuous, \(z_{0} + k \ge 2\); it is
retained for branch symmetry). The class constraint is load-bearing
here: a stagewise recursion over \(A(\Pi) = \{\pm 1\}\) would compute
the unrestricted value, which is \(1\) on the timing cells --- the
agreement is a statement about the segment enumeration;
(b) the certificate partition coincides with the value boundaries:
``viable'' iff \(V_{k} \equiv 1\); common-action-certified iff
\(V_{1}(b_{0}) < 1\); timing-certified iff \(V_{1}(b_{0}) = 1\) but
\(V_{T_{\mathrm{obs}}+1}(b_{0}) < 1\);
(c) at the one-step horizon both mirror witnesses \((1,0)\) and
\((0,1)\) attain the value \(\tfrac{1}{2}\) below \(z_{0} = 2\), and
the merged \((1,1)\) attains \(1\) at and above it; at horizon
\(k \le T_{\mathrm{obs}}\) the family changes at \(z_{0} = 1 + k\) ---
an off-grid consequence of the closed form (in particular \(z_{0} = 3\),
not \(2\), for \(k = 2\) with \(T_{\mathrm{obs}} \ge 2\)), verified
by the verification script beyond the certified grid, which stops at
\(z_{0} = 2.5\); and read as a function of
the continuous parameter \(z_{0}\) the hold-class value has jump
discontinuities exactly at \(z_{0} = 1 + \min(k, T_{\mathrm{obs}})\) --- a
jump of exactly \(\tfrac{1}{2}\), from \(\tfrac{1}{2}\) to \(1\) --- these
loci being the certificate edges on the grid (kinks in the strict sense
belong to the piecewise-linear value as a function of the belief, not of
\(z_{0}\));
(d) on the nonviable cells the hold-class deficit equals
\(\tfrac{1}{2}\) at every horizon \(k > T_{\mathrm{obs}}\) and at
every \(k \le T_{\mathrm{obs}}\) with \(k > z_{0} - 1\), while it is
\(0\) at horizons \(k \le z_{0} - 1\) (on the certified grid: on the
timing cells the deficit is \(0\) at \(k = 1\) and \(\tfrac{1}{2}\)
at \(k \ge 2\); on the common-action cells it is \(\tfrac{1}{2}\) at
every \(k \ge 1\)) --- the minimal-mass bound attained with the
minimal prior from the first horizon at which the falling branch
exits.
\end{theorem}

\begin{proof}[Derivation of the jump loci]
Fix \(T_{\mathrm{obs}}\) and \(k\). On the closed form: if
\(k \le T_{\mathrm{obs}}\) the value is
\(\tfrac{1}{2}(\mathbb{1}[z_{0} \ge 1 + k] + 1)\), the second indicator
being identically one on the grid; it jumps by exactly \(\tfrac{1}{2}\)
at \(z_{0} = 1 + k\). If \(k > T_{\mathrm{obs}}\) the value is
\(\tfrac{1}{2}(1 + \mathbb{1}[z_{0} \ge 1 + T_{\mathrm{obs}}])\),
jumping by exactly \(\tfrac{1}{2}\) at \(z_{0} = 1 + T_{\mathrm{obs}}\).
Both loci are \(z_{0} = 1 + \min(k, T_{\mathrm{obs}})\); between jumps
both expressions are constant in \(z_{0}\). The witness family changes
exactly at the same loci: the attaining one-step vectors are both
mirrors \((1, 0)\) and \((0, 1)\) below \(z_{0} = 2\) and the merged
\((1, 1)\) at and above, and at horizon \(k\) the change
sits at \(z_{0} = 1 + \min(k, T_{\mathrm{obs}})\).
\end{proof}

\begin{theorem}[class lattice and realized gaps]\label{prob-thm:lattice}
For the value of Definition \ref{prob-def:value}, \(\Pi \subseteq \Pi'\)
implies \(V^{\Pi}_{k}(b) \le V^{\Pi'}_{k}(b)\). On the delayed
instance with \(\ell \ge 1\) and \(z_{0} \ge 1\):
\[
V^{\Pi_{\mathrm{hold}}} \;\le\; V^{\Pi_{\mathrm{ol}}}
\;=\; V^{\Pi_{\mathrm{seq, blind}}}
\;\le\; V^{\Pi_{\mathrm{obs}}},
\]
the middle equality holding because no observation arrives inside the
window, so sequential blind policies reduce to open-loop sequences; and
\(V^{\Pi_{\mathrm{obs}}} = \mathbb{1}[z_{0} \ge 1]\) (the controller
knows \(\theta\) at \(t = 0\) and plays the rising control). The gaps
are exact: \(V^{\Pi_{\mathrm{ol}}} - V^{\Pi_{\mathrm{hold}}} =
1 - \max(w, 1 - w)\) exactly on \(\{2 \le z_{0} < 1 + \ell\}\) ---
the \emph{value of flexibility} --- and
\(V^{\Pi_{\mathrm{obs}}} - V^{\Pi_{\mathrm{ol}}} = 1 - \max(w, 1 -
w)\) exactly on \(\{1 \le z_{0} < 2\}\) --- the \emph{value of
observation}; at the symmetric prior each gap is \(\tfrac12\) on its
cells and \(0\) elsewhere, and the two cell families are disjoint.
\end{theorem}

Inclusion monotonicity is immediate from the maximizing definition; the
instance values are Theorem \ref{prob-thm:parametric}'s and the observed
value is the rising computation; the gap arithmetic is
\(1 - \max(w, 1 - w) = \min(w, 1 - w)\).

\begin{proof}
Monotonicity: \(\Pi \subseteq \Pi'\) makes every admissible class
element of \(\Pi\) admissible for \(\Pi'\), and the value is a
maximum over class elements, so
\(V^{\Pi}_{k} \le V^{\Pi'}_{k}\). Middle equality: inside the blind
window no observation arrives, so a sequential blind policy's stage
actions depend only on the observation-free history --- the initial
belief --- and the policy is one fixed sequence: sequential blind
collapses to open-loop. Observed value: knowing \(\theta\) at
\(t = 0\), the controller plays the rising control on every branch
(drift \(+\tfrac3{10}\)), viable exactly from \(z_{0} \ge 1\),
giving \(V^{\Pi_{\mathrm{obs}}} = \mathbb{1}[z_{0} \ge 1]\). The
gaps: hold and open-loop values are the closed forms of
Theorem~\ref{prob-thm:parametric}; subtracting the two case displays, the gap
\(V^{\Pi_{\mathrm{ol}}} - V^{\Pi_{\mathrm{hold}}} = 1 - \max(w,
1 - w)\) survives exactly on the cases where the displays differ ---
\(\{2 \le z_{0} < 1 + \ell\}\) --- and vanishes off them; the
observation gap's locus is the same case analysis against
\(\mathbb{1}[z_{0} \ge 1]\). The script verifies both loci cell by
cell on the certified grid.
\end{proof}

\begin{theorem}[class declaration]\label{prob-thm:class}
The declared hold class and the unrestricted sequential blind class
coincide at the one-step horizon. For horizons \(k \ge 2\) they differ
exactly where \(2 \le z_{0} < 1 + \min(k, T_{\mathrm{obs}})\) --- on
the certified grid \((z_{0} \le 2.5)\) exactly the timing cells, since
\(1 + \min(k, T_{\mathrm{obs}}) \ge 3\) there --- and there open-loop
time-variation within the blind window lifts the value to one: an
alternating hold sequence keeps both branches at or above the floor
whenever \(z_{0} \ge 2\), which no constant hold does (off the grid
the hold class already attains one at \(z_{0} \ge 1 + k\): e.g. at
\((z_{0}, T_{\mathrm{obs}}, k) = (3.5, 3, 2)\) both classes attain
one, so no difference). On the
common-action cells \(z_{0} < 2\) the restriction costs nothing --- no
blind sequence saves both branches --- and on the viable cells both
classes attain one. On the certified grid the difference comprises exactly
the 12 timing cells at each of \(k = 2, 3, 4\): 36 cell-horizon
combinations in all (Figure \ref{prob-fig:classdiff}).
\end{theorem}

\begin{proof}[Derivation of Theorem \ref{prob-thm:class}]
At \(k = 1\) both classes choose one blind action, so they coincide. Let
\(k \ge 2\) and write the two branch positions after a blind sequence
\(u_{1:\ell}\), \(\ell = \min(k, T_{\mathrm{obs}})\), as
\(z_{0} \pm \sum_{t} u_{t}\). If \(z_{0} < 2\), some branch sits below
\(2 - (z_{0} - 1)\) steps from its floor, and since the partial sums of
any \(\pm1\)-sequence move one branch down whenever they move the other
up, one branch reaches its floor: no blind sequence saves both, both
classes give \(\tfrac{1}{2}\), no difference. If
\(z_{0} \ge 1 + \min(k, T_{\mathrm{obs}})\) the constant hold already
survives the window on both branches (the falling branch exits only
after \(k > z_{0} - 1\) steps) and both classes attain one. On the
difference cells \(2 \le z_{0} < 1 + \min(k, T_{\mathrm{obs}})\): a
constant hold drives the
falling branch to \(z_{0} - \ell < 1\) (deficit), while the alternating
sequence has partial sums in \(\{0, \pm 1\}\), so both branches stay at
or above \(z_{0} - 1 \ge 1\) through the window, and revelation at
\(T_{\mathrm{obs}}\) then regenerates both: the unrestricted class
attains one while the declared class does not. This locates the
difference setwise.
\end{proof}

\begin{remark}[the timing obstruction is class-scoped]\label{prob-rem:scope}
The calculus paper's timing certificates --- \(\sigma^{*}(B_{0}) =
z_{0} - 1\), with the certificate firing exactly when
\(T_{\mathrm{obs}} > \sigma^{*}\) (equivalently
\(z_{0} < 1 + T_{\mathrm{obs}}\); viability at the deadline holds when
\(T_{\mathrm{obs}} \le \sigma^{*}\)) --- are statements about the
declared hold-until-review class --- the institutional reading of a
sampled-review regulator. Theorem \ref{prob-thm:class} quantifies the
contrast: within the declared class the timing cells are nonviable with
deficit \(\tfrac{1}{2}\) at every horizon \(k \ge 2\) (and \(0\) at
\(k = 1\)); under the unrestricted sequential blind class the same
cells attain value \(1\) by oscillation, and on the certified grid no
other cell differs. The
timing obstruction therefore certifies exactly what it claims ---
nonviability \emph{within the declared class} --- and no statement of the
companion papers is altered by this reading. The word ``adaptive'' is
avoided for blind windows: while the window is blind, no observation
updates the belief, and the only distinction available is hold-constant
versus time-varying open-loop control.
\end{remark}

\begin{figure}[t]
\centering
\includegraphics[width=\linewidth]{figs_bs2/fig_classdiff.pdf}
\caption{The class declaration on the certified grid: red cells mark
cell-horizon combinations where the declared hold class and the
unrestricted open-loop class differ --- exactly the timing cells
(\(2 \le z_{0} < 1 + T_{\mathrm{obs}}\)) at horizons \(k = 2, 3, 4\)
(three sub-cells per column); 36 combinations in all. The classes agree
at \(k = 1\), on the common-action cells, and on the viable cells.}
\label{prob-fig:classdiff}
\end{figure}

\subsection{Adaptive observation and hidden parameters}\label{prob-adaptive}

Observation is a design variable. Structural uncertainty --- a hidden
parameter resolved too late, rather than noisy data --- is the
obstruction: the state augments with a hidden parameter
\(\theta \in \Theta\), the belief lives on the product
\((x, \theta)\), and the obstruction is the emptiness of the joint
floor-servable set over \((x, \theta)\) cells. On the declining instance
the parameter \(\theta \in \{-1, +1\}\) gates the correctable component
of a \(-\tfrac{1}{10}\) drift; the list of affected sustainability
parameters --- regeneration rates, climate sensitivity, institutional
effectiveness --- motivates the model class, and the paper verifies on
exact instances.

\begin{proposition}[probes and learning deadlines]\label{prob-prop:learn}
On the certified grid, with quantized control \(u \in \{-1, 0, 1\}\) and
hold-then-probe-then-act schedules: with \(\theta\) observed, every cell
is viable (16 of 16). Under one-probe endogenous learning with free probe timing (probe at
the start), the probe
reveals \(\theta\) exactly --- the branch separation is \(2\), so one
observation splits the belief with no statistics --- and survival forces
\(z_{0} \ge 21/10\), independent of any deadline. Under the certified
pinned-probe family --- hold \(u = 0\) to the probe horizon and probe
at \(T_{\mathrm{learn}}\), so no earlier probe is available within the
family --- viability holds exactly when
\[
z_{0} \;\ge\; \frac{21}{10} \;+\; \frac{T_{\mathrm{learn}}}{10},
\qquad T_{\mathrm{learn}} = 0, 1, 2, 3, 4,
\]
the pinned-probe learning-deadline identity: the deadline prices the
forced late probe, and the kernel shrinks exactly
linearly and antitone in it, \(|K| = 5, 4, 3, 2, 1\) on the grid for
\(T_{\mathrm{learn}} = 0, 1, 2, 3, 4\) (Table \ref{prob-tab:deadline} and
Figure \ref{prob-fig:deadline}). Two exact benchmarks bound the family from
both sides: with free probe timing the threshold is flat at
\(21/10\), and with probe-free exogenous revelation at
\(T_{\mathrm{learn}}\) (hold, then the learned act; no probe)
viability holds exactly when \(z_{0} \ge 1 + T_{\mathrm{learn}}/10\)
--- kernel sizes \(16, 15, 14, 13, 12\) on the grid for
\(T_{\mathrm{learn}} = 0, 1, 2, 3, 4\); the certified family sits
between the two. A two-patch instance with unknown
regeneration rate \(g \in \{1, 2\}\) shows the same phenomenon on the
certified two-patch belief pair: unknown-regeneration nonviable,
\((x, g)\)-observed viable.
\end{proposition}

\begin{proof}
The schedule pins the probe to \(T_{\mathrm{learn}}\): within the
family no earlier probe is admissible, so the free-probe and probe-free
thresholds are not attainable in it. Before \(T_{\mathrm{learn}}\) the
drift is \(-\tfrac{1}{10}\) (hold
\(u=0\)), so survival to the probe forces
\(z_{0}-\tfrac{T_{\mathrm{learn}}}{10}\ge\) the probe threshold; the
probe \(u=+1\) moves \(z\) by \(-\tfrac{1}{10}+\theta\), which must stay
safe on both branches, giving the threshold \(\tfrac{21}{10}\) at
\(T_{\mathrm{learn}}=0\); after the probe the learned law \(u=\theta\)
has drift \(-\tfrac{1}{10}+\theta^2=\tfrac{9}{10}>0\), so survival is
indefinite. The branch positions
\(z_{0}-\tfrac{1}{10}T_{\mathrm{learn}}-\tfrac{1}{10}\pm1\) differ by
\(2\), so the post-probe reading determines \(\theta\) uniquely.
Antitonicity is the threshold identity; strictness on the grid is the
computed counts \(5,4,3,2,1\).
\end{proof}

\begin{table}[t]
\centering
\footnotesize
\caption{The pinned-probe learning-deadline audit: the exact threshold
\(\tfrac{21}{10} + \tfrac{T_{\mathrm{learn}}}{10}\) and the resulting
kernel on the grid at each deadline.}
\label{prob-tab:deadline}
\begin{tabular}{@{}l c c l@{}}
\toprule
\(T_{\mathrm{learn}}\) & Threshold & \(|K|\) & Kernel \(K\) on the grid\\
\midrule
0 & $21/10$ & 5 & $\{2.1, 2.2, 2.3, 2.4, 2.5\}$\\
1 & $22/10$ & 4 & $\{2.2, 2.3, 2.4, 2.5\}$\\
2 & $23/10$ & 3 & $\{2.3, 2.4, 2.5\}$\\
3 & $24/10$ & 2 & $\{2.4, 2.5\}$\\
4 & $25/10$ & 1 & $\{2.5\}$\\
\bottomrule
\end{tabular}
\end{table}

\begin{figure}[t]
\centering
\includegraphics[width=\linewidth]{figs_bs2/fig_deadline.pdf}
\caption{The pinned-probe learning-deadline identity: the viability
threshold increases exactly linearly in \(T_{\mathrm{learn}}\) (left),
and the kernel on the grid shrinks linearly and antitone, \(|K| = 5, 4,
3, 2, 1\) (right); the probe-free benchmark \(1 +
T_{\mathrm{learn}}/10\) (kernels \(16\) down to \(12\)) and the
flat free-probe threshold \(21/10\) bound it from below and above.}
\label{prob-fig:deadline}
\end{figure}

\begin{figure}[t]
\centering
\includegraphics[width=0.88\linewidth]{figs_bs2/fig_twofloor.pdf}
\caption{The two-floor chance instance: the one-step value is
\(\max(b_{0}, b_{1})\), attained by the alpha-vectors of
\(\Gamma^{\Pi}_{1} = \{(1,0), (0,1)\}\) with zero bottom coordinate (green
probes); at the symmetric prior the value is \(\tfrac{1}{2}\) and the
deficit equals the minimal mass.}
\label{prob-fig:twofloor}
\end{figure}

\subsection{Genuinely stochastic instances}\label{prob-stochastic}

Everything to this point has deterministic kernels and observations; the
values are prior masses of surviving branches. This section opens the
stochastic layer, where values cease to be two-valued and the min-mass
bound ceases to be universal.

\begin{proposition}[weighted-path deficit identity]\label{prob-prop:pathdeficit}
At the one-step horizon, with finite \(X\), deterministic per-reading
actions \(a(y')\), and arbitrary normalized observations, the deficit
is exactly
\[
\Delta_{1}(b) \;=\; \sum_{x} b(x)\, p_{x},
\qquad p_{x} = \mathbb{P}\bigl(a(y') \text{ unsafe at } x \,\big|\, x\bigr),
\]
the prior mass weighted by fatal-reading probabilities. The min-mass
bound is the case \(p_{x} \equiv 1\) (deterministic readings, no
common action) and is \emph{not universal}: in the two-state instance
with \(a_{i}\) safe only at \(x_{i}\) and sensor error \(\varepsilon\)
(\(\mathbb{P}(\text{reads } j \ne i \mid x_{i}) = \varepsilon\)),
the deficit is \(\varepsilon\) for every prior; at
\(b = (\tfrac12, \tfrac12)\), \(\varepsilon = \tfrac1{20}\), the
deficit \(\tfrac1{20}\) is strictly below the minimal mass
\(\tfrac12\). This delimits the deterministic-limit reach of the
calculus's probabilistic development (its Section~10): there the
deficit from unity is bounded below by the smallest compatible mass
under adversarial disturbance readings, an estimate exact at the
degeneration this paper's Proposition~\ref{prob-prop:degen} records; the
sensor instance shows the estimate is scoped to that limit, not a law
of the noisy layer. Conversely, the calculus's delayed-hidden-regime
instance is the deadline law of
Proposition~\ref{prob-prop:additivelaw} at \(d_{0} = 1\), \(e_{p} = 0\):
its viability boundary \(z_{0} \ge 1 + T_{\mathrm{obs}}\) is the
law's threshold form, and the tabulated agreement there is
cell-by-cell.
\end{proposition}

\begin{proof}
Fix a deterministic per-reading action map \(a(\cdot)\). One-step
survival is \(\sum_{x} b(x)\, \mathbb{P}(a(y') \text{ safe} \mid x)
= 1 - \sum_{x} b(x)\, p_{x}^{a(\cdot)}\), and optimizing over maps
gives \(V_{1}(b) = 1 - \min_{a(\cdot)} \sum_{x} b(x)\,
p_{x}^{a(\cdot)}\) --- the displayed identity at the optimizing map. In
the two-state instance the reading-own-action map makes
\(p_{x} = \varepsilon\) for both states (the wrong reading arrives
with probability \(\varepsilon\) and its action is fatal), and no map
does better: one of the two actions is fatal at each state whichever
reading carries it with positive probability. Hence
\(\Delta_{1}(b) = \varepsilon\) at every prior, and at
\(b = (\tfrac12, \tfrac12)\), \(\varepsilon = \tfrac1{20}\), the
deficit \(\tfrac1{20}\) is strictly below the minimal mass
\(\tfrac12\).
\end{proof}

\begin{proposition}[noisy probes and the repetition trade-off]\label{prob-prop:noisyprobe}
On the declining instance (blind drift \(-\tfrac1{10}\)), pin a single
probe to horizon \(t_{p}\) whose reading is correct with probability
\(1 - \varepsilon\) and flipped with probability \(\varepsilon\),
then act on the reading to horizon \(k \ge t_{p}\). With
\(z_{p} = z_{0} - t_{p}/10\):
\[
V_{k}(b_{0}) \;=\; (1 - \varepsilon)\,\mathbb{1}[z_{p} \ge 1]
\;+\; \varepsilon\,\mathbb{1}\bigl[z_{p} \ge 1 + \tfrac{11}{10}(k - t_{p})\bigr],
\]
exactly (the correct reading acts under drift \(\tfrac9{10}\), the
flipped reading under drift \(-\tfrac{11}{10}\)). The values realize
exactly three levels \(\{0,\, 1 - \varepsilon,\, 1\}\) ---
genuinely non-two-valued --- and the min-mass bound survives here as a
bound (the deficit \(\varepsilon\) or \(1\) dominates the
post-probe minimal mass \(\varepsilon\)) while failing as an
identity. With three pinned probes and majority reading, the
wrong-majority probability is
\(3\varepsilon^{2}(1 - \varepsilon) + \varepsilon^{3} = \varepsilon^{2}(3 - 2\varepsilon)\)
--- \(\tfrac{7}{250}\) at \(\varepsilon = \tfrac1{10}\), against
\(\tfrac1{10}\) for one probe --- and the same display holds with
\(p = \tfrac{7}{250}\) in place of \(\varepsilon\): each
additional probe step costs \(\tfrac1{10}\) of floor margin to the
blind drift, the exact explore--exploit exchange at this scale.
\end{proposition}

\begin{proof}
Condition on \((\theta, \text{reading})\): the four joint cases carry
probabilities \(\tfrac12(1-\varepsilon)\) and \(\tfrac12\varepsilon\)
per branch; the correct reading acts under drift
\(-\tfrac1{10} + \theta^{2} = \tfrac9{10}\) and survives from
\(z_{p} \ge 1\); the flipped reading acts under drift
\(-\tfrac1{10} - 1 = -\tfrac{11}{10}\) and survives exactly when
\(z_{p} - \tfrac{11}{10}(k - t_{p}) \ge 1\); summing gives the
display. For the majority family, the wrong-majority event is two or
three flipped readings, of probability
\(3\varepsilon^{2}(1-\varepsilon) + \varepsilon^{3}\); the falling
continuation is as above. The identity's proof is the law of total
probability over readings at the one-step horizon.
\end{proof}

\subsection{Distributionally robust interpolation}\label{prob-dr}

The expected and adversarial readings of Remark \ref{prob-rem:operators} are
the two ends of one exact family. Fix the contamination ambiguity set
\(\mathcal{D}_{\rho} = \{q = (1 - \rho)\,p + \rho\, r : r \in
\Delta(Y)\}\) at radius \(\rho \in [0, 1]\).

\begin{proposition}[contamination interpolation]\label{prob-prop:dr}
Let \(\mathcal{D}_{\rho}\) be as above and define
\(V^{\rho}_{k+1}(b) = \max_{a}\inf_{q \in \mathcal{D}_{\rho}}
\sum_{y} q(y)\, V^{\rho}_{k}(\tau(b, a, y))\). Then:
(i) the inner infimum is exact,
\[
\inf_{q \in \mathcal{D}_{\rho}} \textstyle\sum_{y} q(y)\, v_{y}
\;=\; (1 - \rho) \sum_{y} p(y)\, v_{y} \;+\; \rho \min_{y} v_{y},
\]
the adversary spending the whole contamination \(\rho\) on the worst
outcome;
(ii) \(V^{\rho}_{k}(b) = (1 - \rho)\, V^{\,0}_{k}(b) + \rho\,
V^{\,1}_{k}(b)\) at every stage: the interpolation is pointwise exact,
with \(\rho = 0\) the expected value of Definition \ref{prob-def:value}
and \(\rho = 1\) the adversarial reading;
(iii) \(V^{\rho}_{k}\) is concave and piecewise linear in \(b\) at
every \(\rho\), so the exact-arithmetic discipline of Proposition
\ref{prob-prop:pl} carries over verbatim.
\end{proposition}

\begin{proof}
The infimum of a linear functional over the contamination set is attained
at a vertex \(r = \delta_{y^{*}}\), giving (i); (ii) follows by
induction with (i) at every stage, since the adversary's optimal
contamination is stage-wise; concavity is preserved under pointwise
mixtures of concave piecewise-linear functions. On the delayed class
with prior contamination \(q(+1) = (1 - \rho)w + \rho\, r\), the
hold-class value is exactly
\[
V^{\rho,\, \Pi_{\mathrm{hold}}}_{k}(b_{0}) =
\begin{cases} 0, & z_{0} < 1,\\
1, & z_{0} \ge 1 + \ell,\\
(1 - \rho) \max(w, 1 - w), & 1 \le z_{0} < 1 + \ell,
\end{cases}
\]
the adversary pouring the contamination onto a falling branch; at the
symmetric prior the three \(\rho\)-levels are \(0\), \(
\tfrac{1 - \rho}{2}\), and \(1\) (verified by exact enumeration on
the swept grid).
\end{proof} This proposition is the formal home of worst-case
wording: an ambiguity-minded reviewer may read every expected value here
at any \(\rho\), and every table carries the \(\rho = 0\) column of
a one-parameter family.

\begin{proposition}[probe-count bound]\label{prob-prop:probecount}
Let \(n = 2m + 1\) pinned probes with per-probe error
\(\varepsilon < 1/2\) be read by majority. The wrong-majority
probability satisfies
\[
p_{\mathrm{wrong}}(n) \;\le\; \bigl(4\varepsilon(1 - \varepsilon)\bigr)^{n/2}
\;\le\; e^{-2 n\, \mathrm{sep}^{2}},
\qquad \mathrm{sep} = \tfrac12 - \varepsilon,
\]
so \(n \ge \ln(1/\delta) / (2\,\mathrm{sep}^{2})\) probes achieve
wrong probability \(\le \delta\): the probe count grows like
\(\log(1/\delta)\) over the squared separation, and on the declining
instance each probe step costs \(\tfrac1{10}\) of floor margin to the
blind drift (the exploration--safety exchange of Proposition
\ref{prob-prop:noisyprobe}, now with its rate). At
\(\varepsilon = \tfrac1{10}\) the bound reads
\(p_{\mathrm{wrong}}(n) \le (3/5)^{n}\), verified exactly at
\(n = 1, 3, 5, 7, 9\) (the exact values \(\tfrac1{10},
\tfrac7{250}, \tfrac{107}{12500}, \tfrac{341}{125000},
\tfrac{22273}{25000000}\)).
\end{proposition}

\begin{proof}
The bound is the standard misselection probability of the binomial
head: \(p_{\mathrm{wrong}}(n) \le (4\varepsilon(1-\varepsilon))^{n/2}\)
by pairing each wrong-majority string with its complement and
Stirling's estimate, and \(4\varepsilon(1 - \varepsilon) =
1 - 4\,\mathrm{sep}^{2} \le e^{-4\,\mathrm{sep}^{2}}\).
\end{proof}

\begin{remark}[noiseless limit of the probe-count law]
\label{prob-rem:eps0}
The exact split of Proposition~\ref{prob-prop:learn} is the
\(\varepsilon \to 0\) degeneration of the probe-count bound: the
separation \(\mathrm{sep} = \tfrac12 - \varepsilon\) rises to
\(\tfrac12\) and \(p_{\mathrm{wrong}}(1) = \varepsilon\) falls to
\(0\), so one probe already achieves \(\delta = 0\) --- the branch
split of separation \(2\) needs no statistics precisely because the
bound's exponent vanishes with \(\varepsilon\). For
\(\varepsilon > 0\) the count law prices what the exact split gets
for free, and on the declining instance each probe step additionally
pays its \(\tfrac1{10}\) of floor margin
(Proposition~\ref{prob-prop:noisyprobe}).
\end{remark}

\begin{proposition}[general additive deadline law]\label{prob-prop:additivelaw}
Let the hidden parameter enter the drift additively:
\(\dot z = d(u, \theta) = c(u) + h(u)\,\theta\) with
\(\theta \in \Theta\) finite, floor \(z \ge 1\), and consider the
hold--probe--learn schedule family: hold \(u_{0}\) to the probe horizon
\(T\) (the probe pinned there, no earlier probe in the family), probe
with \(u_{p}\), then play \(u^{*}(\theta)\) with
\(\min_{\theta} d(u^{*}(\theta), \theta) = \mu > 0\). Set
\(d_{0} = \max_{\theta}(-d(u_{0}, \theta))\) and
\(e_{p} = \max_{\theta}(-d(u_{p}, \theta))\). Then indefinite
viability holds exactly when
\[
z_{0} \;\ge\; 1 \;+\; T\, d_{0} \;+\; e_{p},
\]
the floor plus drift times deadline plus probe excursion; after the
learned act the positive drift \(\mu\) preserves safety forever, so
the deadline identity is the whole answer. The certified instance has
\(d_{0} = \tfrac1{10}\), \(e_{p} = \tfrac{11}{10}\), giving
\(\tfrac{21}{10} + \tfrac{T}{10}\); a second instance with drift
\(-\tfrac{3}{20}\) and excursion \(\tfrac{21}{20}\) gives
\(\tfrac{41}{20} + \tfrac{3T}{20}\) --- both verified exactly at all
swept \((z_{0}, T)\).
\end{proposition}

\begin{proof}
Necessity: some \(\theta\) realizes the hold displacement \(d_{0}\)
and the probe excursion \(e_{p}\) at the adversarial stages (the maxima
are attained), dropping the state to
\(z_{0} - T d_{0} - e_{p}\), below the floor exactly when the display
fails. Sufficiency: every \(\theta\) keeps the state at or above the
floor through hold and probe, and the learned drift \(\mu > 0\) then
lifts it away from the floor forever.
\end{proof}

\medskip
\noindent\textbf{Reciprocal check (dual-certificate companion).} On the
two-window residue instance of the minimax companion (Abaee, 2026, Minimax dual certificates), the
interpolation family above is exact at its coupling extremum:
\(\sup_{q \in \mathcal{D}_{\rho}} \mathbb{E}_{q}[G] = \rho\) at
every probed \(\rho \in \{\tfrac14, \tfrac12, \tfrac34, 1\}\),
the product coupling attaining the \(\rho = 1\) end, with the
one-step marginals preserved at every \(\rho\) and the
coupling-selection difficulty left orthogonal to the mixture parameter;
the joint check is archived with the companion's verification record.

\subsection{Verification methods}\label{prob-methods}

Every value, bound, identity, threshold, figure, and tabulated entry in
this paper is regenerated by deterministic scripts in exact integer and
rational arithmetic, standard library only. Three per-system scripts cover the worked instances; one master entry point executes all of them.
The computational companion papers (exact belief computation and its certification) carry the scaling record --- exact
point-based evaluation at probed beliefs (in the POMDP literature's
sense, Lovejoy, 1991 --- here exact at the probed beliefs, not bounds),
the antichain compression of
survivable-set lattices, and the four-parameter cube's Hamming
classification --- built on the alpha-vector representation of
Proposition~\ref{prob-prop:pl}; its scripts chain this paper's seeds but
are not part of this paper's chain.
The belief-state campaign script runs sixteen check families on the
stochastic values (the two-floor instance, piecewise-linearity, the
deterministic limit, the 48-cell agreement, the partition coincidence,
the witness family, the deficit identity, the unrestricted class,
rationality of all 384 campaign alpha-vectors, deficit monotonicity, the
oscillation identity, model completeness, the selector split, dimension
consistency, the jump loci, and the declarations), chaining two
per-instance scripts of twenty-one further checks. The stochastic
selector script re-derives every selector claim in fifteen check
families (S1--S15: the two-floor backup, eleven-belief piecewise
linearity, degeneration against brute force, 48-cell closed-form
agreement, partition coincidence, witness-family loci, the exact
deficit on all 42 nonviable cells, the class declaration by brute
force, rationality of the 384 vectors, monotonicity, the oscillation
identity to horizon 24, model completeness, the selector split,
dimension padding, and the one-sided jump values). The hidden-parameter
script runs six check families on the adaptive-observation instance
(theta observed on all 16 cells, the probe threshold
\(\tfrac{21}{10}\) as a branchwise identity, the deadline identity at
\(T_{\mathrm{learn}} = 0, \dots, 4\), injective revelation, the
antitone kernel chain, and the unknown-regeneration two-patch
instance). The master script for this paper chains the three, then
recomputes independently from the definitions: the closed-form values on
the whole \(16 \times 12\) grid by direct branch-mass enumeration, the
resulting census (\(36\) unit values and \(156\) half values on the
declared class), the identity of the class-difference set with the region
\(2 \le z_{0} < 1 + \min(k, T_{\mathrm{obs}})\), checked on an
extended grid to \(z_{0} = 4.0\) and equal to the timing-cell
combinations on the certified grid at \(k \ge 2\), by brute-force
enumeration over blind sequences; the \(384\)-vector type census
(\(156\), \(156\), \(72\)) and its deduplicated \(348\)-vector
form; the hold-class deficit profile on all 48 cells (\(0\) before the
falling-branch exit, \(\tfrac{1}{2}\) after); the lost-branch step of
the asymmetric audit (step \(1\) for all four sequences, heavier
branch under \(u_{1} = +1\), lighter under \(u_{1} = -1\)); the
pinned-probe learning-deadline chain
\(\tfrac{21}{10} + \tfrac{T_{\mathrm{learn}}}{10}\) with the
free-probe (\(21/10\), flat) and probe-free
(\(1 + T_{\mathrm{learn}}/10\), kernels \(16\) down to \(12\))
benchmarks; and every headline
number asserted present in the text and bound to the scripts' certified families. Two independent paths back the closed forms: a
brute-force path that enumerates trajectories straight from the
definitions (all holds and all open-loop sequences, pathwise survival)
over a structured sweep of 2268 triples
\((z_{0}, w, \ell)\) including \(z_{0} < 1\), and a seeded
randomized battery of 500 further cases; the survivable-set antichain,
the class ordering and its gaps, the stabilization freeze (spot
enumeration to horizon \(15\) and value agreement to horizon
\(40\)), the weighted-path identity, the sensor counterexample, and
the noisy-probe and majority closed forms are each re-derived by
independent enumeration; the contamination lemma and the delayed-class
\(\rho\)-closed form are checked against vertex enumeration, the
probe-count bound at \(n = 1, 3, 5, 7, 9\) in exact rationals, and the
additive deadline law on both instances.

\subsection{Delimitations and outlook}\label{prob-scope}

The semantics here is stochastic expectation: the value is a survival
\emph{probability}, and the robust quantifiers of the companion papers
are recovered only at the deterministic degeneration, which is a change
of operator (Remark \ref{prob-rem:operators}); the stochastic and robust
readings of a genuinely stochastic plant differ, and this paper does not
conflate them. The timing layer's statements, here and in the companion
papers, are class-scoped (Remark \ref{prob-rem:scope}). No learning theory is
claimed beyond the one-probe exact split: Bayesian updating, repeated
experiments, and exploration--exploitation trade-offs are out of scope;
the deadline identity is for the pinned-probe schedule family (hold to
the probe horizon, probe, learned act), with the free-probe and
probe-free benchmarks stated alongside, and other schedules are not
classified; the parameter
\(\theta\) enters additively, and state-dependent parameter dynamics
belong to the hybrid companion's setting. The alpha-set growth is
exponential in the worst case --- and finite-horizon partially
observable problems are PSPACE-complete, their no-observation variant
NP-complete (Papadimitriou and Tsitsiklis, 1987) --- so the certified
two-witness families are structural, not representative. The contamination interpolation of Section \ref{prob-dr} is the formal home
of worst-case wording; the moment- and distance-based ambiguity families
(piecewise-linear values persist there; Nakao, Jiang, and Shen, 2021)
are the natural next ring and are not developed here.
Scaling past the certified instances --- exact point-based value iteration
with rational solvers, survivable-set compression over the antichain of
Proposition \ref{prob-prop:antichain}, and medium instances with several
hidden parameters --- the probe-count asymptotics of the repetition
trade-off, the general learning-deadline law for additively entering
hidden parameters, are the natural next steps. A third item stood on that
list when this paper was drafted --- a proof-assistant
formalization of the closed forms --- and is now complete:
the support identity, the freeze count and its frozen value,
and the rationality of the \(\alpha\)-vectors are machine-checked
(Lean~4; no axioms, no admitted gaps, ordered-field interface
only), the first on the repaired reading of
Theorem~\ref{prob-thm:support} recorded in
Remark~\ref{prob-rem:feedback-strict}. One bound remains cited
rather than formalized: the Sperner bound of
Proposition~\ref{prob-prop:antichain} (iii), for which the layer
supplies the binomial machinery needed to \emph{state} it but
no proof. The deficit-bound reading of
Proposition \ref{prob-prop:deficit} remains the interface through which this
theory attaches to the certificate families already established, and the
exact rational values of the certified campaign are the benchmark any
relaxation must meet.

\subsection{Conclusion}\label{prob-conclusion}

The belief-state safety value carries the obstruction calculus into the
probabilistic setting without loss of exactness. The support identity
ties the value-one level to the viable-set recursion; rational
alpha-vectors keep every certified value exact, and the complete value
table shows the two-valued structure the certificates predict; the
certificates survive as deficit lower bounds, attained on natural
instances on both the symmetric and the asymmetric prior; the certified
delayed-regime class agrees with its closed form cell by cell, with the class declaration responsible for exactly the 36 timing
combinations on the grid; and the adaptive-observation audit prices late resolution exactly ---
for the pinned-probe family a linear learning deadline on the viable
set, \(|K| = 5, 4, 3, 2, 1\), bracketed by the flat free-probe
threshold and the probe-free benchmark --- while the stochastic layer
prices observation noise exactly: a weighted-path deficit identity, a
sensor instance breaking the min-mass bound's universality, and a
repetition trade-off at \(7/250\) per extra probe pair. The
observation layer, not the forecast layer, sets the value; the design
rules are the same ones the deterministic audits impose, now with a
price attached and every number in the tables open to script-level
re-inspection.

\section{Related work}
\label{prob-priorart}

This paper asks when an observation-based policy is \emph{sufficient}, exactly, and what
remains safe when it is not. Both halves have near neighbours, and this section names them
before stating what is claimed.

\subsection{The belief-state reduction and the $\alpha$-vector representation}

That a partially observed problem can be written as a fully observed one on the space of
beliefs is classical (Astr\"{o}m, 1965), and the finite-horizon value function of a
partially observed Markov decision process is piecewise linear and convex in the belief,
represented by a finite set of $\alpha$-vectors (Smallwood and Sondik, 1973; see Bertsekas
and Shreve, 1978, for the measure-theoretic setting). The safety value $V_k(b)$ studied here
is a belief-space value in that lineage, with one structural change: the recursion
maximises over actions but \emph{minimises} over the disturbance, so the object is a
worst-case safety probability rather than an expected total reward. Consequently the
piecewise-linear representation survives but its identifying structure does not: here the
$\alpha$-vectors are exactly the indicators of the maximal jointly survivable subsets of the
support, an antichain (Proposition~\ref{prob-prop:antichain}), rather than the
reward-derived vectors of the expected-criterion theory.

\subsection{Sufficiency under imperfect measurement: the two nearest neighbours}

The closest work is the viability-theoretic treatment of imperfect information, and it comes
in two forms that this paper sits between.

\emph{Veliov (1993).} For a time-invariant system with incomplete and inexact measurement,
Veliov gives a \textbf{sufficient} condition for the existence of an output feedback
regulation map, and shows it is equivalent to Haddad's viability condition when the
measurement is perfect. The condition is sufficient and not necessary: a system may admit an
adequate output feedback that the condition does not detect.

\emph{Doyen (2000).} \textbf{(Luc Doyen, the viability theorist --- not Laurent Doyen of the
qualitative partially-observed literature cited below; the surnames coincide and the two
should not be conflated.)} For uncertain nonlinear systems under state and control
constraints with bounded disturbances on both the dynamics and the output, Doyen gives
conditions that are \textbf{necessary and sufficient} for the guaranteed viability property
--- defined as the existence of a \emph{Lipschitz closed-loop} maintaining the state
\emph{exactly} in a given closed domain --- and derives equivalent geometric and
Hamilton--Jacobi--Isaacs characterisations via the Lipschitz kernel of a set-valued map.

These two delimit the contribution. Relative to Veliov, what is added is exactness: the
condition here is necessary as well as sufficient, so its failure is informative rather than
merely inconclusive. Relative to Doyen, what differs is the admissible policy class and the
criterion. Doyen's loops are Lipschitz and \emph{memoryless}, of the form $u(h(x,w))$; the
recursion here admits policies that depend on the whole observation history, and the
distinction is strict --- the feedback viable-set recursion
$\mathcal{W}^{\mathrm{fb}}_k$ contains the blind recursion $\mathcal{W}_k$, strictly in
general (Remark~\ref{prob-rem:feedback-strict}). And where Doyen's criterion is the exact
maintenance of a closed domain, the safety value here is quantitative: when a belief is not
fully safe, the deficit satisfies
$1 - V_k(b) \ge \min_{x \in \mathrm{supp}(b)} b(x)$
(Proposition~\ref{prob-prop:deficit}), so near-misses are measured rather than discarded.

\subsection{The estimation-space and qualitative programmes}

A distinct strategy absorbs imperfect information into an enlarged state: the estimation-tube
programme formulates the problem on a space of estimators or tubes rather than on the belief
simplex (Cardaliaguet, Quincampoix and Saint-Pierre, 2007, for the survey of viability
methods in this line). That programme is set in continuous time and targets differential
games; the present development is discrete, finite-horizon, and computational, and it does
not compete with it.

On the qualitative side, the decidability and complexity landscape for partially observed
Markov decision processes under qualitative objectives --- reachability, safety,
B\"{u}chi --- is established (Chatterjee, Doyen and Henzinger, 2009), and it is known that
almost-sure safety depends only on belief supports. What is added here is quantitative and
structural: the deficit bound and its witness reading, the survivable-set (antichain) formula
behind the witness census, the class lattice with its realized gaps, the parametric closed
forms, and the class-scoped timing obstruction --- together with the observation that the
support-only reduction, valid for almost-sure objectives, does \emph{not} extend to the
quantitative criterion, which is why the value, not just the support, has to be carried.

\subsection{Complexity, robustness and shielding}

That finite-horizon partially observed problems are computationally hard is classical
(Papadimitriou and Tsitsiklis, 1987), and the practical response has been bounds and
approximations (Lovejoy, 1991). This paper does not improve those complexity results; it
identifies classes on which exact answers are nonetheless available in closed form, and
records where exactness stops. Distributionally robust formulations place an ambiguity set
over the transition or the observation kernel (Nakao, Jiang and Shen, 2021); the disturbance
model here is worst-case over a support rather than distributional, so the two are not
comparable and neither subsumes the other. Safe reinforcement learning via shielding
(Alshiekh, Bloem, Ehlers, K\"{o}nighofer, Niekum and Topcu, 2018) enforces safety by
constraining a learner's actions; the setting there is learning with an unknown model, here
the model is known and the question is exact sufficiency.

\subsection{What is claimed}

Given the above, the contribution is five things and nothing more.

\begin{enumerate}
\item \textbf{Exact sufficiency where the nearest sufficient condition is only sufficient.}
  The value-one level of the recursion coincides with the feedback viable-set recursion, so
  the condition is necessary as well as sufficient (Theorem~\ref{prob-thm:support}).
\item \textbf{Admission of history-dependent policies}, with the feedback/blind distinction
  shown to be strict (Remark~\ref{prob-rem:feedback-strict}) --- a repaired reading; the
  mechanized support identity is checked on the repaired statement.
\item \textbf{A quantitative deficit rather than a binary verdict}:
  $1 - V_k(b) \ge \min_{x \in \mathrm{supp}(b)} b(x)$, with the witness reading of
  Proposition~\ref{prob-prop:deficit}.
\item \textbf{The structural identification of the $\alpha$-vectors} as the indicators of
  the maximal jointly survivable subsets of the support, Sperner-bounded
  (Proposition~\ref{prob-prop:antichain}), the class lattice with realized gaps
  (Theorem~\ref{prob-thm:lattice}), and parametric closed forms on named classes.
\item \textbf{Machine-checking}: the support identity, the freeze count and its frozen value,
  and the rationality of the $\alpha$-vectors are formalized in Lean~4 with no axioms and no
  admitted gaps. One bound remains cited rather than formalized: the Sperner bound of
  Proposition~\ref{prob-prop:antichain}(iii).
\end{enumerate}

\noindent Not claimed: the belief-state reduction (Astr\"{o}m), the $\alpha$-vector
representation as such (Smallwood and Sondik), the qualitative complexity landscape
(Chatterjee, Doyen and Henzinger), or the hardness results that motivate approximation
(Papadimitriou and Tsitsiklis; Lovejoy).

\section{Cross-part conclusion}
\label{conclusion}

The two parts above ask one question --- \emph{can any observation-based policy keep this system
within its constraints?} --- in the two languages available for it, and the answer is the same in
both. That agreement is the paper's central result, and it is worth stating precisely, because it
is easy to overstate.

\emph{What the two readings share, exactly.} The probabilistic theory is not an approximation to
the combinatorial one, and the combinatorial one is not the degenerate case of the probabilistic
one. They meet at an identity: the value-one level of the belief-state safety recursion
coincides with the viable-set recursion of the calculus, the posterior support coincides with the
set-valued post-state on finite models with deterministic kernels and observation maps, and
outside the value-one level the deficit satisfies
$1 - V_k(b) \ge \min_{x \in \mathrm{supp}(b)} b(x)$, attained on natural instances. Moreover the
$\alpha$-vectors are exactly the indicators of the maximal jointly survivable subsets of the
support --- an \emph{antichain}, Sperner-bounded, which is why the certified witness census
realizes three types rather than the continuum a generic belief-space value would permit. So the
probabilistic reading does not add a degree of confidence to a set-valued verdict; it \emph{prices}
the same verdict, and the price is exact.

\emph{Why this matters more than the identity itself.} The direction that has received the
systematic treatment in the literature is \emph{sufficiency}: Veliov's output-feedback regulation
condition, and the estimation-tube reduction, give the canonical answer to ``is there a policy
that works?''. The complementary direction --- certifying that \emph{no} observation-based policy
is viable --- has received less. That asymmetry is a practical problem, because a negative verdict
is what licenses a redesign: it is the difference between ``no policy we tried worked'' and
``no policy can work, and here is the finite object that proves it.'' Both parts supply finite,
checkable, \emph{exact} witnesses for that negative direction, and the fact that they do so in
integer and rational arithmetic means the witness can be re-inspected at script level rather than
trusted. \emph{A nonviability certificate is an artefact that a reader can check, not a claim a
reader must accept.}

\emph{What does not transfer.} The taxonomy is explicitly nonexhaustive: the mechanisms do not
exhaust the complement of the epistemic kernel, and only in finite systems is backward recursion
complete. So a system on which every mechanism here fails to fire is not thereby shown to be
viable --- the certificates are sound sufficient conditions for nonviability, not necessary ones,
and the paper states the gap rather than papering over it. Correspondingly, the agreement between
the two readings is proved under the stated structural hypotheses (finite models, deterministic
kernels and observation maps, rational input data); where those fail the two readings may come
apart, and nothing here asserts otherwise.

\emph{The design consequence, stated generally.} Every mechanism in Part I empties a specific
response correspondence, and each therefore licenses a specific \emph{design} response rather
than a generic appeal to better information: enlarging the command set when the admissible safe
actions are empty; a separating observation, or a shorter review interval, when the common safe
action set is empty; an earlier informative observation when the deadline is exceeded; a refined
index when the fibre criterion fails; correcting a known bias when the certainty-equivalence trap
binds. Part II's contribution to this is to attach a \emph{price} to each --- the adaptive-observation
audit prices late resolution exactly, and the stochastic layer prices observation noise. The
general lesson is that \emph{the observation layer, not the forecast layer, sets the value}: when
a managed system is failing, the remedial leverage is usually in what is measured and when, not
in predicting better from what is already measured.

\subsection*{References}
\label{references}
Abaee, A., 2026. Robust viability of the 2J3KL limit reference point
under a surplus-production map: policy scoring, expansion, and when
catch cannot help. Zenodo.
\url{https://doi.org/10.5281/zenodo.22552060}

Abaee, A., 2026. Minimax dual certificates for the common-action
obstruction (submitted).

Academic Press, New York.

\`{A}str\"{o}m, K.J., 1965. Optimal control of Markov processes with
incomplete state information. J. Math. Anal. Appl. \textbf{10},
174--205.

%% ------------------------------------------------------------------
%% Added 2026-09-29 to support section \ref{priorart}. All three VERIFIED
%% against publisher records (volume, pages, venue confirmed):
%%   Veliov, Set-Valued Anal. 1, 305-317 (1993), DOI 10.1007/BF01027640
%%   Doyen, Set-Valued Anal. 8, 149-162 (2000)
%%   Cardaliaguet/Quincampoix/Saint-Pierre, Advances in Dynamic Game Theory,
%%     Ann. ISDG vol. 9, Birkhauser Boston, pp. 3-35 (2007)
%% NOTE: Luc Doyen (2000, viability) and Laurent Doyen (2009, qualitative
%% POMDPs) are DIFFERENT AUTHORS. Both are cited; do not merge them.
%% ------------------------------------------------------------------
Cardaliaguet, P., Quincampoix, M., Saint-Pierre, P., 2007. Differential games through
viability theory: old and recent results. In: J{\o}rgensen, S., Quincampoix, M., Vincent,
T.L. (Eds.), \emph{Advances in Dynamic Game Theory}. Annals of the International Society of
Dynamic Games, vol. 9, pp. 3--35.

Alshiekh, M., Bloem, R., Ehlers, R., K\"{o}nighofer, B., Niekum, S.,
Topcu, U., 2018. Safe reinforcement learning via shielding. Proc. AAAI
Conf. Artif. Intell. \textbf{32}(1).

Aubin, J.-P.: Viability Theory.

Bertsekas, D.P., Shreve, S.E., 1978. Stochastic Optimal Control: The
Discrete-Time Case.

Birkh\"{a}user, Boston.

Birkh\"auser, Boston (2011)

Aubin, J.-P., Catt\'e, F.: Bilateral fixed-points and algebraic properties of viability kernels and capture basins of sets. Set-Valued Anal. \textbf{10}, 379--416 (2002)

Aubin, J.-P., Frankowska, H.: Set-Valued Analysis.

Birkh\"auser, Boston (1991)

Aubin, J.-P.: Viability kernels and capture basins of sets under differential inclusions. SIAM J. Control Optim. \textbf{40}(3), 853--871 (2001)

Aubin, J.-P., Bayen, A.M., Saint-Pierre, P.: Viability Theory: New Directions, 2nd edn.

Birkh\"auser, Boston (1990)

Bansal, S., Chen, M., Herbert, S.L., Tomlin, C.J.: Hamilton--Jacobi reachability: a brief overview and recent advances. In: Proc. 56th IEEE Conf. on Decision and Control (CDC) (2017)

Bansal, S., Tomlin, C.J.: DeepReach: a deep learning approach to high-dimensional reachability. In: Proc. IEEE Int. Conf. on Robotics and Automation (ICRA) (2021)

B\'en\'e, C., Doyen, L., Gabay, D.: A viability analysis for a bio-economic
model. Ecol. Econ. \textbf{36}, 385--396 (2001)

Cardaliaguet, P., Quincampoix, M., Saint-Pierre, P.: Differential games
through viability theory: old and recent results. In: J\o{}rgensen, S.,
Quincampoix, M., Vincent, T.L. (eds.) Advances in Dynamic Game Theory.
Annals of the International Society of Dynamic Games, vol.~9, pp.~3--35.

Birkh\"auser, Boston (2007)

Clark, C.W.: Mathematical Bioeconomics: The Optimal Management of Renewable Resources, 2nd edn.

Birkh\"auser, Boston (1997)

Maghenem, M., Sanfelice, R.G.: Characterizations of safety in hybrid
inclusions via barrier functions. In: Proceedings of the 22nd ACM
International Conference on Hybrid Systems: Computation and Control
(HSCC), pp.~109--118 (2019)

Margellos, K., Lygeros, J.: Hamilton--Jacobi formulation for reach--avoid differential games. IEEE Trans. Automat. Control \textbf{56}(8), 1849--1861 (2011)

Martinez-Alier, J., Munda, G., O'Neill, J.: Weak comparability of values
as a foundation for ecological economics. Ecol. Econ. \textbf{26}(3),
277--286 (1998)

Mitchell, I.M., Bayen, A.M., Tomlin, C.J.: A time-dependent Hamilton--Jacobi formulation of reachable sets for continuous dynamic games. IEEE Trans. Automat. Control \textbf{50}(7), 947--957 (2005)

Nagumo, M.: \'Uber die Lage der Integralkurven gew\'ohnlicher
Differentialgleichungen. Proc. Phys.-Math. Soc.

Chatterjee, K., Doyen, L., Henzinger, T.A., 2009. Qualitative analysis
of partially-observable Markov decision processes. In: Int. Symp. Math.
Found. Comput. Sci. (MFCS 2009).

Doyen, L., 2000. Guaranteed output feedback control for uncertain systems under control and
state constraints. \emph{Set-Valued Analysis} \textbf{8}, 149--162.

Japan 24, 551--559
(1942)

Prajna, S., Jadbabaie, A.: Safety verification of hybrid systems using
barrier certificates. In: Alur, R., Pappas, G.J. (eds.) Hybrid Systems:
Computation and Control (HSCC 2004).

LNCS, vol.~2993, pp.~477--492.

LNCS, vol. 5734, pp. 635--646.

Lovejoy, W.S., 1991. Computationally feasible bounds for partially
observed Markov decision processes.

McGraw-Hill, New York
(1960)

Haddad, G.: Monotone viable trajectories for functional-differential
inclusions. J. Differ. Equ. \textbf{42}, 1--24 (1981)

Hale, J.K., Verduyn Lunel, S.M.: Introduction to Functional Differential
Equations.

Nakao, H., Jiang, R., Shen, S., 2021. Distributionally robust partially
observable Markov decision process with moment-based ambiguity. SIAM J.
Optim. \textbf{31}(1), 461--488.

Operations Research 39, 162--175.

Papadimitriou, C.H., Tsitsiklis, J.N., 1987. The complexity of Markov
decision processes. Math. Oper. Res. \textbf{12}(3), 441--450.

\r{A}str\"om, K.J.: Optimal control of Markov processes with incomplete state information. J. Math. Anal. Appl. \textbf{10}, 174--205 (1965)

Abaee, A.: Aggregate indices and transition safety: a quantifier-order separation between scalarized and coordinate-wise feasibility. figshare. https://doi.org/10.6084/m9.figshare.33764023 (2026).

Smallwood, R.D., Sondik, E.J., 1973. The optimal control of partially
observable Markov processes over a finite horizon. Oper. Res.
\textbf{21}(5), 1071--1088.

Springer, Berlin (2008)

Doyen, L.: Guaranteed output feedback control for uncertain systems under control and state constraints. Set-Valued Anal. \textbf{8}, 149--162 (2000)

Doyen, L., Gajardo, P.: Sustainability standards, multicriteria maximin,
and viability. Nat. Resour. Model. \textbf{33}(3), e12250 (2020)

Doyen, L., Th\'ebaud, O., B\'en\'e, C., Martinet, V., Gourguet, S., Bertignac,
M., Fifas, S., Blanchard, F.: A stochastic viability approach to
ecosystem-based fisheries management. Ecol. Econ. \textbf{75}, 32--42
(2012)

Farkas, J.: Theorie der einfachen Ungleichungen. J. Reine Angew. Math.
\textbf{124}, 1--27 (1902)

Frankowska, H.: Optimal trajectories associated with a solution of
contingent Hamilton--Jacobi equations. Appl. Math. Optim. \textbf{19},
291--311 (1989)

Gale, D.: The Theory of Linear Economic Models.

Springer, Berlin (2004)

Prajna, S., Jadbabaie, A., Pappas, G.J.: A framework for worst-case and
stochastic safety verification using barrier certificates. IEEE Trans.
Autom. Control \textbf{52}(8), 1415--1428 (2007)

Prajna, S., Rantzer, A.: On the necessity of barrier certificates. IFAC
Proc. Vol. \textbf{38}(1), 526--531 (2005)

Quincampoix, M., Veliov, V.M.: Viability with a target: theory and
applications. In: Control Theory, Multivariate Analysis and
Applications, pp.~47--63. Springer (1994)

Saint-Pierre, P.: Approximation of the viability kernel. Appl. Math.
Optim. \textbf{29}, 187--209 (1994)

Smallwood, R.D., Sondik, E.J.: The optimal control of partially
observable Markov processes over a finite horizon. Oper. Res. 21(5),
1071--1088 (1973)

Sontag, E.D.: Mathematical Control Theory: Deterministic Finite Dimensional Systems, 2nd edn.

Springer, New York (1993)

Kurzhanski, A.B., V\'alyi, I.: Ellipsoidal Calculus for Estimation and
Control.

Springer, New York (1998)

Veliov, V.M.: Sufficient conditions for viability under imperfect
measurement. Set-Valued Anal. \textbf{1}, 305--317 (1993)

Witsenhausen, H.S.: A counterexample in stochastic optimum control.
SIAM J. Control \textbf{6}(1), 131--147 (1968)

Veliov, V.M., 1993. Sufficient conditions for viability under imperfect measurement.
\emph{Set-Valued Analysis} \textbf{1}, 305--317.


{\footnotesize
Abaee, A., 2026. An obstruction calculus for viability under incomplete
observation (submitted).

Wiley, New York (1990)

Crandall, M.G., Ishii, H., Lions, P.-L.: User's guide to viscosity solutions of second order partial differential equations. Bull. Amer. Math. Soc. (N.S.) \textbf{27}, 1--67 (1992)

De Lara, M., Doyen, L.: Sustainable Management of Natural Resources:
Mathematical Models and Methods.

\section*{Declarations}

\textbf{Funding.} None.\quad \textbf{Competing interests.} None.\quad \textbf{Data availability.} No external data were used.\quad \textbf{Code availability.} The coverage-audit script \texttt{paper2\_coverage\_audit.py} is publicly available at \url{https://github.com/MIKEAA2020/general-sustainability} (folder \texttt{arena agent 1/paper rewrites/latex}); it regenerates Table~\ref{calc-tab:coverage} and Figure~\ref{calc-fig:coverage} verbatim.\quad \textbf{AI declaration.} GLM (Z.ai), Qwen (Alibaba Cloud) and DeepSeek AI assisted with drafting and iterative review.
\subsection*{Declarations}

\subsection*{Funding}
None.

\subsection*{Competing interests}
None.

\subsection*{Data availability}
No external data were used.

\subsection*{Code availability}
The verification script
\url{paper2_probabilistic_sufficiency_v14_verification.py} --- which chains
the per-system scripts \url{paper2_belief_state_v2_verification.py},
\url{paper2_stochastic_selector_v2_verify.py}, and
\url{hidden_parameter_learning_v1_verify.py} --- and the figure-and-table
generation script \url{paper2_belief_state_figures.py} are publicly
available at \url{https://github.com/MIKEAA2020/general-sustainability}
(folder \texttt{arena agent 1/paper rewrites/latex}, at the commit
pinned to this article's submission); they reproduce all values, bounds,
thresholds, figures, and tabulated entries verbatim.

\subsection*{AI declaration}
GLM (Z.ai), Qwen (Alibaba Cloud) and DeepSeek AI assisted with drafting
and iterative review. The author reviewed and edited outputs and takes
responsibility for the final work.

\end{document}
